\documentclass[11pt,a4paper,oneside,openany]{book}
\usepackage[T1]{fontenc}
\usepackage[utf8]{inputenc}
\usepackage[english]{babel}
\usepackage{lmodern,amsmath,amssymb,amsthm,mathtools,bm}
\usepackage{booktabs,longtable,geometry,graphicx,microtype,enumitem}
\usepackage{hyperref,xcolor,fancyhdr,listings}
\geometry{margin=2.55cm,headheight=14pt}
\hypersetup{colorlinks=true,linkcolor=blue!55!black,citecolor=blue!55!black,urlcolor=blue!55!black,pdftitle={From Ordinary Language to Computational Neuroscience},pdfauthor={CON2EPHYS research notebook}}
\lstdefinelanguage{LeanFour}{
  morekeywords={import,namespace,end,theorem,lemma,def,by,variable,open,
    where,fun,have,show,exact,intro,apply,example},
  sensitive=true,
  morecomment=[l]--,
  morestring=[b]"
}
\lstset{
  literate={¬}{{$\neg$}}1
           {∧}{{$\land$}}1
           {∨}{{$\lor$}}1
           {↔}{{$\leftrightarrow$}}1
           {→}{{$\to$}}1
           {∀}{{$\forall$}}1
           {∃}{{$\exists$}}1
           {α}{{$\alpha$}}1
           {⊆}{{$\subseteq$}}1
           {⟨}{{$\langle$}}1
           {⟩}{{$\rangle$}}1
           {·}{{$\cdot$}}1
}
\numberwithin{equation}{chapter}
\newtheorem{definition}{Definition}[chapter]
\newtheorem{proposition}[definition]{Proposition}
\newtheorem{theorem}[definition]{Theorem}
\newtheorem{remark}[definition]{Remark}
\theoremstyle{definition}
\newtheorem{workedexample}[definition]{Worked example}
\newtheorem{readergoal}[definition]{Reading goal}
\newenvironment{exercises}{\par\medskip\noindent\textbf{Exercises}\begin{enumerate}[leftmargin=2em,itemsep=0.18em]}{\end{enumerate}\medskip}
\newcommand{\E}{\mathbb E}
\newcommand{\Var}{\operatorname{Var}}
\newcommand{\Cov}{\operatorname{Cov}}
\newcommand{\R}{\mathbb R}
\newcommand{\C}{\mathbb C}
\newcommand{\Prob}{\mathbb P}
\newcommand{\diag}{\operatorname{diag}}
\newcommand{\rank}{\operatorname{rank}}
\newcounter{bridgecount}
\newcommand{\bridgeentry}[2]{%
  \refstepcounter{bridgecount}%
  \section*{Bridge \thebridgecount: #2}%
  \addcontentsline{toc}{section}{Bridge \thebridgecount: #2}%
  \label{bridge:#1}}

\pagestyle{fancy}\fancyhf{}\fancyhead[L]{\small CON$^2$EPHYS Mathematics}\fancyhead[R]{\small\nouppercase{\leftmark}}\fancyfoot[C]{\thepage}
\setlength{\parskip}{0.25em}
\title{\Huge From Ordinary Language to\\ Computational Neuroscience\\[1em]\Large A Mathematical Foundation for Electrophysiological Analysis\\[0.5em]\normalsize A modular companion to CON$^2$EPHYS \texttt{arsenal/path}}
\author{CON$^2$EPHYS Research Manuscript}
\date{Bridge edition --- October 2026\\\small Independent exposition inspired by Halmos and Mendelson}
\begin{document}
\frontmatter
\maketitle
\chapter*{Preface}\addcontentsline{toc}{chapter}{Preface}
This book begins with statements in ordinary English, not with calculus, matrix notation, or a list of methods. Its reader need not know what a set, function, derivative, matrix, probability, or oscillation is. We will name each mathematical object before asking the reader to manipulate it, and we will distinguish statements about abstract models from statements about nervous systems. The purpose is not to make mathematics superficial: it is to make the deepest methods accessible by building their prerequisites instead of assuming them.

We follow two complementary pedagogical inspirations. Halmos develops the set-theoretic language of mathematics through informal but precise exposition; Mendelson connects logic, sets, and successive number systems to the foundations of analysis. Neither book is copied here. Their themes guide an original scientific text culminating in signal processing, neuronal point processes, statistical inference, and population dynamics.

Every chapter includes formal mathematics, an electrophysiological interpretation, limitations, and exercises. A reader should practice proving elementary results as well as simulating advanced algorithms. Source-data quality, estimand definitions, and null hypotheses receive the same attention as derivations. The earlier edition presented a provisional \texttt{arsenal/path} mapping. The repository is now verified; for actual MATLAB filenames and method status, consult the companion theorem crosswalk and \texttt{arsenal/METHODS.tsv}.
\tableofcontents
\chapter{How to read this book}
The intended reader begins with ordinary English and curiosity about neural signals. Part I requires no mathematical notation beyond what it introduces. Parts II and III build number systems, limits, and finite-dimensional algebra before the book assumes calculus, probability, or complex phase. Later chapters use the same format: meaning in words, formal definition, calculation or derivation, a neuroscience connection, a failure mode, and exercises. A novice may work every exercise; an experienced reader can use the cross-reference path below to fill only missing prerequisites.

\section*{Suggested prerequisite chains}
\begin{description}[leftmargin=2cm,style=nextline]
\item[Hilbert and PAC] English $\to$ functions $\to$ complex numbers $\to$ sequences/integrals $\to$ Fourier $\to$ analytic signal $\to$ circular inference.
\item[PCA and manifolds] Relations $\to$ functions $\to$ real numbers $\to$ vector spaces $\to$ covariance $\to$ SVD/PCA $\to$ dynamical systems.
\item[Spike GLMs] Sets $\to$ functions $\to$ counting $\to$ probability $\to$ conditional probability $\to$ Poisson processes $\to$ likelihood and GLM.
\item[Granger analysis] Covariance $\to$ stochastic processes $\to$ cross spectra $\to$ multivariate linear prediction $\to$ null models and validation.
\item[Current source density] Coordinates $\to$ derivatives $\to$ multivariable calculus $\to$ forward measurement physics $\to$ regularized inverse methods.
\end{description}

\section*{Technical vocabulary}
An \emph{estimand} is the quantity the analysis intends to estimate; an \emph{estimator} is a data-dependent rule for estimating it. A \emph{model} is a set of assumed mathematical relationships; a \emph{measurement} is an observation through an instrument. \emph{Association} is not automatically \emph{causation}. \emph{Prediction} need not identify an intervention effect. The word ``significant'' should be accompanied by the hypothesis and statistical criterion being used.

\section*{About the two foundational books}
The structure draws inspiration from Halmos' conversational construction of sets and relations \cite{HalmosNaive}, and Mendelson's progression from logic and set theory to the number systems and real-number completeness \cite{MendelsonNumbers}. The explanations, examples, and proofs here are written afresh for CON$^2$EPHYS and are not reproduced sections of either text.
\mainmatter
\part{The language of mathematics}
\chapter{From ordinary English to mathematical statements}\label{ch:english}
\begin{readergoal}Read, question, and translate statements before manipulating symbols. No mathematical prerequisite is assumed.\end{readergoal}
\section{Why begin with sentences?}
A scientist looks at an experiment and makes claims: ``this animal ran faster,'' ``these spikes followed the stimulus,'' or ``these electrodes recorded the same rhythm.'' At first, the claims are expressed in ordinary language. Mathematics begins when we decide exactly what would make a claim true, how it could be false, and what follows from it. The symbols are a shorthand for decisions already made in words. Learning a formula without learning its grammar is like memorizing a sentence in a language one cannot read.

Consider ``every recorded neuron fired.'' Does this refer to every neuron in the brain, every electrode channel, or every neuron retained after spike sorting? The statement is incomplete until we specify its domain of discourse. Consider ``some neurons were silent.'' The word \emph{some} means \emph{at least one}; it does not exclude the possibility that many neurons were silent. The word \emph{or} in standard mathematics ordinarily permits both alternatives. When a manuscript says ``the signal is oscillatory,'' we must also decide what counts as an oscillation. A label is not yet an operational test.

\section{Statements, questions, definitions, and hypotheses}
A \emph{proposition} is a declarative statement to which a truth value can be assigned, given a specified interpretation. ``Two is even'' is a proposition. ``Is two even?'' is a question, while ``Let $n$ be the number of spikes'' is an instruction introducing notation. ``For every $t$ in the recording interval, $x(t)>0$'' becomes a proposition once $x$ and the recording interval are fixed. An \emph{open sentence} such as ``$n$ is even'' has a free variable and does not yet assert anything about a particular number.

A \emph{definition} establishes how a term or symbol will be used; a \emph{hypothesis} asserts a fact that may or may not hold; a \emph{theorem} is a proposition established from previously stated assumptions by valid reasoning. In applied science we must further distinguish a mathematical theorem, which follows within a model, from an empirical assertion about the world, which requires observations.

\section{Translating ordinary English}
Let $U$ be the set of units retained after quality control. Introduce the predicate $F(u)$ for ``unit $u$ fired in a designated 100-ms window.'' Then ``every retained unit fired'' becomes
\[ \forall u\in U,\ F(u). \]
``At least one retained unit did not fire'' becomes $\exists u\in U$ such that $\neg F(u)$. These are logical negations of one another. ``No unit fired'' instead means $\forall u\in U,\neg F(u)$. Notice the difference between a negative statement about \emph{all} units and a statement that \emph{not all} units fired.

In data analysis, ``higher gamma amplitude when theta phase is near its peak'' must be turned into an explicit comparison: define phase bins, summarize amplitudes, select a test statistic, and specify a null hypothesis. This transition from sentence to mathematical object is the beginning of sound methodology.

\section{A claim is only as good as its scope}
Suppose 40 of 50 detected units increase firing rate during a task. It is legitimate to say ``80 percent of these detected units increased.'' It is not automatically legitimate to say ``80 percent of neurons in this brain area increased.'' Electrode sampling, sorting errors, and selection rules change the population to which a claim refers. We will repeatedly separate \emph{a set of observed objects}, \emph{a mathematical model}, and \emph{a physiological population}.

\begin{workedexample}[The ambiguity of ``only'']
Compare ``only neurons in CA1 show a theta rhythm'' with ``neurons in CA1 show only a theta rhythm.'' The first restricts the location of units exhibiting theta; the second restricts the rhythms observed in CA1 units. Write down what observation would contradict each statement. The two claims can disagree even on the same recording.
\end{workedexample}
\begin{exercises}
\item Translate ``not every trial contains a ripple'' into a quantified sentence and simplify its negation.
\item Explain why ``all spikes occur in bursts'' can be vacuously true for an empty spike train.
\item Identify the observational unit in each claim: neuron, trial, event, animal, and session.
\item Write one empirically falsifiable claim about coherence and one purely mathematical definition of coherence.
\end{exercises}
\chapter{Logic: the grammar of proof}\label{ch:logic}
\begin{readergoal}Learn the connectives, truth conditions, quantifiers, and the difference between logical implication and scientific causation.\end{readergoal}
\section{And, or, and not}
Let $P$ mean ``the animal is running'' and $Q$ mean ``theta power exceeds a specified threshold.'' The conjunction $P\land Q$ is true precisely when both claims are true. The disjunction $P\lor Q$ is false only when both are false; in mathematical usage it is \emph{inclusive}. Negation $\neg P$ reverses the truth value of $P$. The entire truth table is small enough to check directly:
\[\begin{array}{cc|ccc}P&Q&P\land Q&P\lor Q&P\Rightarrow Q\\\hline
T&T&T&T&T\\T&F&F&T&F\\F&T&F&T&T\\F&F&F&F&T\end{array}\]
The conditional $P\Rightarrow Q$ is false only when $P$ holds but $Q$ does not. This is the \emph{material implication}; it is not an assertion of a mechanism or a causal influence. Indeed, a false antecedent makes the material implication true without producing any biological event.

\section{Converse, inverse, and contrapositive}
From $P\Rightarrow Q$ one may infer the logically equivalent contrapositive $\neg Q\Rightarrow\neg P$. The converse $Q\Rightarrow P$ and inverse $\neg P\Rightarrow\neg Q$ are different claims. A low firing rate might be observed during inactivity, but inactivity need not be the only circumstance producing that rate. We may write $P\Leftrightarrow Q$ for a \emph{biconditional}, meaning that both implications hold.

De Morgan's rules supply a useful bridge to data selection: $\neg(P\land Q)\Leftrightarrow(\neg P\lor\neg Q)$ and $\neg(P\lor Q)\Leftrightarrow(\neg P\land\neg Q)$. Thus the complement of trials satisfying both inclusion criteria is the collection failing at least one criterion.

\section{Quantifiers and their order}
The phrase ``for every'' is denoted $\forall$, and ``there exists'' by $\exists$. These operations are not interchangeable: $\forall i\,\exists t\,F(i,t)$ says each neuron fires at some time, possibly a different time for every neuron. The stronger statement $\exists t\,\forall i\,F(i,t)$ claims there is one time at which every neuron fires. Neither form alone specifies simultaneity to a millisecond: that depends on how $t$ and $F$ are defined.

Negation propagates through quantifiers by the rules
\[\neg\forall x\,P(x)\ \Leftrightarrow\ \exists x\,\neg P(x),\qquad
\neg\exists x\,P(x)\ \Leftrightarrow\ \forall x\,\neg P(x).\]
For a finite dataset, verifying a universal claim requires checking all included cases; refuting it requires only one counterexample.

\section{Necessary and sufficient conditions}
If $P\Rightarrow Q$, then $P$ is \emph{sufficient} for $Q$, while $Q$ is \emph{necessary} for $P$. An example: if a matrix $A$ is invertible, its columns are linearly independent. In the square finite-dimensional case, the converse also holds; the statements are equivalent. Contrast this with ``if the signal is a sinusoid, it has a spectral peak,'' whose converse may fail for many reasons. Whenever an algorithm claims to ``detect oscillations,'' ask whether its acceptance rule is necessary, sufficient, both, or neither for the physiological interpretation.

\section{Formal proof laboratory: propositional and quantified logic}
\label{sec:lean-logic-pilot}
The five mathematical statements below are shown alongside their
Lean 4 implementations. A successful Lean build means the kernel
accepts the propositions \emph{as typed}; it does not establish
the scientific interpretation of a null hypothesis.

\begingroup
\setlength{\columnsep}{0.05\textwidth}
\noindent\begin{minipage}[t]{0.45\textwidth}
\textbf{English mathematics}

\textbf{LOGIC-001.} De Morgan's conjunction law:
for arbitrary propositions $P,Q$,
$\neg(P\land Q)\iff\neg P\lor\neg Q$.
If not both criteria hold, at least one fails. \emph{Proof:} by excluded middle, either $P$ is false, or $P$ holds, in which case $Q$ must be false. Conversely, either negated conjunct contradicts $P\land Q$.

\medskip\textbf{LOGIC-002.} De Morgan's disjunction law:
$\neg(P\lor Q)\iff\neg P\land\neg Q$.
To fail either-of-two inclusion rules is to fail both. \emph{Proof:} $\neg(P\lor Q)$ separately contradicts each disjunct. Conversely, $\neg P\land\neg Q$ refutes every proposed witness for $P\lor Q$.

\medskip\textbf{LOGIC-003.} Contraposition:
$(P\Rightarrow Q)\iff(\neg Q\Rightarrow\neg P)$.
The backward direction is classical: assuming
$\neg Q\Rightarrow\neg P$, if $P$ held but $Q$ did not,
one would obtain $\neg P$, a contradiction; by excluded middle, $Q$ must hold. The forward implication follows constructively: from $P\Rightarrow Q$ and $\neg Q$, assuming $P$ would contradict $\neg Q$.

\medskip\textbf{LOGIC-004.} Quantifier negation:
$\neg(\forall x\,R(x))\iff\exists x\,\neg R(x)$.
A counterexample disproves a universal claim. \emph{Proof:} a failing witness contradicts the universal assertion. Conversely, if no witness failed, classical double-negation elimination would give $R(x)$ for every $x$, contradicting $\neg\forall x\,R(x)$.

\medskip\textbf{LOGIC-005.} Existential negation:
$\neg(\exists x\,R(x))\iff\forall x\,\neg R(x)$. \emph{Proof:} if a witness existed, it would contradict the assumption that none exists. Conversely, if every candidate fails, an assumed witness immediately yields a contradiction.

\medskip\textbf{Proof strategy.} The statements about negating
conjunction, contraposition and negating a universal
quantifier use classical reasoning. The Lean proof
invokes Mathlib tactics that elaborate these elementary
steps; the file prints their axiom dependencies.
\end{minipage}\hfill
\begin{minipage}[t]{0.50\textwidth}
\textbf{Lean 4 / Mathlib}

\lstinputlisting[language=LeanFour,basicstyle=\ttfamily\tiny,
breaklines=true,columns=fullflexible,keepspaces=true,
frame=single]{../../formal/lean/Con2physFormal/Logic.lean}

\textit{Correspondence and provenance:}
each Lean declaration has the same quantification
over propositions or predicates as the adjacent English
statement. Mathlib provides the underlying logic and tactics.
See the pinned dependency in
\texttt{\detokenize{formal/lean/lakefile.lean}}.
\end{minipage}
\endgroup


\begin{workedexample}[A null hypothesis in plain language]
Suppose $H_0$ states that the temporal positions of spikes inside specified jitter windows carry no additional information beyond their counts. Rejecting $H_0$ is evidence against this particular conditional model. It does \emph{not} establish that a synapse connects the neurons, or that all alternative null models would also be rejected. Compare \cite{Amarasingham2012}.
\end{workedexample}
% --- expanded edition (October 2026) ---
\section{Tautologies, contradictions, and valid arguments}
A compound formula is a \emph{tautology} if it is true under every assignment of truth values to its atomic letters, and a \emph{contradiction} if it is false under every assignment. The law of excluded middle $P\lor\neg P$ is a tautology; $P\land\neg P$ is a contradiction. Two formulas $\varphi$ and $\psi$ are \emph{logically equivalent}, written $\varphi\equiv\psi$, when $\varphi\Leftrightarrow\psi$ is a tautology.

An \emph{argument} consists of premises $\varphi_1,\ldots,\varphi_k$ and a conclusion $\psi$. It is \emph{valid} when every assignment making all premises true also makes $\psi$ true; equivalently, when $(\varphi_1\land\cdots\land\varphi_k)\Rightarrow\psi$ is a tautology. Validity concerns form, not content: a valid argument with a false premise can have a false conclusion, and an invalid argument can happen to have a true conclusion.

\begin{proposition}[Four standard inference rules]
The following argument forms are valid:
\begin{enumerate}[leftmargin=2em]
\item \emph{modus ponens}: from $P$ and $P\Rightarrow Q$, infer $Q$;
\item \emph{modus tollens}: from $\neg Q$ and $P\Rightarrow Q$, infer $\neg P$;
\item \emph{hypothetical syllogism}: from $P\Rightarrow Q$ and $Q\Rightarrow R$, infer $P\Rightarrow R$;
\item \emph{disjunctive syllogism}: from $P\lor Q$ and $\neg P$, infer $Q$.
\end{enumerate}
\end{proposition}
\begin{proof}
We check (2); the others are similar. Suppose $\neg Q$ and $P\Rightarrow Q$ are both true. If $P$ were true, then by the truth table of the conditional $Q$ would be true, contradicting $\neg Q$. Hence $P$ is false, i.e.\ $\neg P$ is true. For (3), suppose both premises are true and $P$ is true; then $Q$ is true by the first premise and $R$ by the second, so $P\Rightarrow R$ cannot be falsified. For (4), if $P\lor Q$ is true and $P$ is false, the only remaining way for the disjunction to hold is that $Q$ is true.
\end{proof}

Two \emph{invalid} forms are so common that they have names. \emph{Affirming the consequent} infers $P$ from $Q$ and $P\Rightarrow Q$; \emph{denying the antecedent} infers $\neg Q$ from $\neg P$ and $P\Rightarrow Q$. In both cases the assignment $P=F$, $Q=T$ makes the premises true and the conclusion false. A frequent scientific instance: ``if the region generates gamma oscillations, the LFP shows a gamma peak; the LFP shows a gamma peak; therefore the region generates gamma.'' Volume conduction from another region, or a harmonic of a lower rhythm, also produces the peak.

\section{Useful equivalences}
The following equivalences can each be verified with a four- or eight-row truth table and will be used without comment later.
\[
\begin{array}{rcl@{\qquad}l}
P\Rightarrow Q&\equiv&\neg P\lor Q & \text{(material implication)}\\
P\Rightarrow Q&\equiv&\neg Q\Rightarrow\neg P & \text{(contraposition)}\\
\neg(P\Rightarrow Q)&\equiv&P\land\neg Q & \text{(negated conditional)}\\
P\land(Q\lor R)&\equiv&(P\land Q)\lor(P\land R) & \text{(distributivity)}\\
(P\land Q)\Rightarrow R&\equiv&P\Rightarrow(Q\Rightarrow R) & \text{(exportation)}\\
P\Leftrightarrow Q&\equiv&(P\Rightarrow Q)\land(Q\Rightarrow P) & \text{(biconditional)}
\end{array}
\]
The negated conditional is the most practically useful: to show that ``$P$ implies $Q$'' fails, one must exhibit a case where $P$ holds and $Q$ does not. A case where $P$ fails is irrelevant.

\begin{proposition}
For any propositions, $\neg(P\Rightarrow Q)\equiv P\land\neg Q$.
\end{proposition}
\begin{proof}
By material implication, $\neg(P\Rightarrow Q)\equiv\neg(\neg P\lor Q)$. By De Morgan's rule this is $\neg\neg P\land\neg Q$, and double negation gives $P\land\neg Q$.
\end{proof}

\section{Nested quantifiers and their negations}
Many definitions in analysis have three or more alternating quantifiers. The definition of a limit, met in Chapter~\ref{ch:limits}, reads
\[
\forall\varepsilon>0\ \exists N\ \forall n\ge N:\ |x_n-L|<\varepsilon.
\]
Its negation is obtained mechanically by moving $\neg$ inward one quantifier at a time, swapping $\forall$ and $\exists$:
\[
\exists\varepsilon>0\ \forall N\ \exists n\ge N:\ |x_n-L|\ge\varepsilon.
\]
In words: there is a fixed tolerance that infinitely late terms keep violating. Note that the bounded quantifiers $\forall\varepsilon>0$ and $\exists n\ge N$ keep their bounds under negation; only the quantifier type and the inner predicate change.

The \emph{uniqueness quantifier} $\exists!x\,P(x)$ (``there is exactly one $x$ with $P$'') abbreviates
\[
\exists x\,\bigl(P(x)\land\forall y\,(P(y)\Rightarrow y=x)\bigr).
\]
Existence and uniqueness are proved separately: one exhibits an $x$, then shows that any two such objects coincide.

\begin{workedexample}[Negating a selectivity criterion]
Suppose a unit is called ``direction selective'' if there is a direction $d^\ast$ such that for every other tested direction $d$, the mean rate at $d^\ast$ exceeds the rate at $d$ by at least 5 spikes/s. Symbolically, $\exists d^\ast\,\forall d\ne d^\ast:\ r(d^\ast)-r(d)\ge5$. Its negation, ``not direction selective'', is $\forall d^\ast\,\exists d\ne d^\ast:\ r(d^\ast)-r(d)<5$: every candidate preferred direction has at least one competitor within 5 spikes/s. This is a much weaker statement than ``the unit responds equally to all directions,'' and the two must not be conflated in a results section.
\end{workedexample}

\section{Logic and uncertainty}
Classical logic assigns each proposition exactly one of two truth values. Empirical claims usually come with uncertainty, and it is tempting to treat a $p$-value or posterior probability as a ``degree of truth.'' This is a category mistake. A probability is assigned by a model to an event; the proposition ``this event occurred'' remains true or false. Probability theory (Part~V) extends logic by attaching numbers to sets of outcomes, but every probabilistic statement is itself a classical proposition about those numbers, and proofs about probabilities use exactly the rules of this chapter.
\begin{exercises}
\item Construct a complete truth table for $(P\Rightarrow Q)\Leftrightarrow(\neg Q\Rightarrow\neg P)$.
\item Negate ``for every frequency there exists a channel with nonzero power.''
\item Give a counterexample to the inference from $P\Rightarrow Q$ and $Q$ to $P$.
\item Explain why statistical dependence is not equivalent to causal influence.
\item Use a truth table to show that $(P\Rightarrow Q)\land(Q\Rightarrow R)\Rightarrow(P\Rightarrow R)$ is a tautology.
\item Show that affirming the consequent is invalid by exhibiting a falsifying truth assignment.
\item Negate $\forall\varepsilon>0\ \exists\delta>0\ \forall x\ (|x-a|<\delta\Rightarrow|f(x)-f(a)|<\varepsilon)$ and state the result in words.
\item Write ``exactly one channel exceeds the artifact threshold'' using $\exists!$ and then without it.
\item Prove the exportation law $(P\land Q)\Rightarrow R\equiv P\Rightarrow(Q\Rightarrow R)$.
\end{exercises}

\chapter{Definitions, proof, and mathematical maturity}\label{ch:proof}
\begin{readergoal}Learn to construct and audit arguments, identify hidden hypotheses, and write counterexamples.\end{readergoal}
\section{The anatomy of an argument}
A proof begins with an exact statement and a list of assumptions. It proceeds through steps justified by a definition, an axiom, or an earlier result. The conclusion is valid \emph{conditional on those assumptions}. Mathematical evidence and empirical evidence answer different questions: a covariance matrix is positive semidefinite by construction, but a measured covariance need not describe only neuronal interactions.

\section{Direct proof and proof by contrapositive}
We say an integer $n$ is even if $n=2k$ for some integer $k$. To prove that the sum of two even integers is even, write $a=2k$ and $b=2m$. Then $a+b=2(k+m)$; since $k+m$ is an integer, the sum is even. Notice the explicit use of the definition. An assertion such as ``the result is obvious'' hides exactly the step a beginner may need.

For the contrapositive method, consider: if $n^2$ is even, then $n$ is even. Prove instead that odd $n=2k+1$ implies $n^2=4k(k+1)+1$, which is odd. This avoids assuming the conclusion.

\section{Contradiction, induction, and counterexamples}
Proof by contradiction begins with the negation of the proposed conclusion, then derives an impossibility. The classic irrationality argument for $\sqrt2$ assumes $\sqrt2=a/b$ with coprime integers $a,b$, derives that both must be even, and contradicts coprimality. We shall revisit this when constructing the real numbers.

A single valid counterexample refutes a universal proposition. To refute ``every high-power spectral peak is neuronal,'' one can simulate a sinusoidal electrical artifact. To refute ``correlation implies coupling,'' one can simulate two independently noisy signals influenced by a common source. These examples do not show that neuronal oscillations or coupling never occur; they refute the unrestricted inference.

\section{Existence, uniqueness, and identifiability}
Three questions recur throughout this book: Does the object exist? Is it uniquely determined? Can we infer it from the observations? The first two are mathematical; the third also depends on a measurement model. For example, different neuronal current distributions may produce the same measured extracellular potential, so a valid forward model does not guarantee a unique inverse solution.

\begin{proposition}If $A$ and $B$ are subsets of a fixed universe $U$, then $A=B$ if and only if both $A\subseteq B$ and $B\subseteq A$.\end{proposition}
\begin{proof}If $A=B$, either inclusion follows immediately. Conversely, the first inclusion implies every element of $A$ belongs to $B$, and the second implies every element of $B$ belongs to $A$. The two sets therefore have exactly the same elements.\end{proof}
\begin{exercises}
\item Prove that the intersection of two subsets is contained in each subset.
\item Give two distinct functions with the same value at zero.
\item Distinguish ``the algorithm converges'' from ``the result estimates the intended biological quantity.''
\item Write a three-line counterexample to ``a zero Pearson correlation implies independence.''
\end{exercises}
\chapter{Sets and the operations that create them}\label{ch:sets}
\begin{readergoal}Use membership, inclusion, unions, intersections, set differences, power sets, and the idea of an admissible collection.\end{readergoal}
\section{From a collection to a set}
A set is a mathematical collection whose members are determined by a membership rule. If $U=\{a,b,c\}$, then $a\in U$ and $d\notin U$. In the ordinary language of mathematics, $a\in U$ means that $a$ \emph{belongs to} $U$; $A\subseteq U$ means every element of $A$ belongs to $U$. These statements differ: $\{a\}\subseteq U$, but typically $\{a\}\notin U$. The empty set $\varnothing$ has no members; $\{\varnothing\}$ has one member, which happens to be the empty set.

Two sets are equal exactly when they have identical members (the extensional criterion). Mathematical set construction must take place within an explicit foundational framework; one cannot freely declare every grammatical property to specify a set of all possible objects. The unrestricted collection ``all sets that are not members of themselves'' generates Russell's paradox. Following the accessible axiomatic spirit of \cite{HalmosNaive,MendelsonNumbers}, we allow subsets defined by conditions on an already established set and use conventional constructions for other sets.

\section{Set operations from ordinary English}
Let $A$ be the trials in which the animal runs, and $B$ those with a detected theta peak. Then $A\cap B$ consists of trials satisfying \emph{both} criteria, while $A\cup B$ contains those satisfying \emph{at least one}. Set difference $A\setminus B$ contains trials that are in $A$ but not $B$. Relative complement $A^c$ is meaningful only after choosing the universe $U$: here $A^c=U\setminus A$. We must never conflate the universe of trials in one session with all possible trials.

De Morgan's identities are
\[(A\cup B)^c=A^c\cap B^c,\qquad(A\cap B)^c=A^c\cup B^c.\]
For $U=\{1,2,3,4,5\}$, $A=\{1,3,5\}$, and $B=\{3,4\}$, we have $A\cap B=\{3\}$ and $A\cup B=\{1,3,4,5\}$. Work through the identities element by element before trusting the symbols.

\section{Families, indexing, and power sets}
A \emph{family} $(A_i)_{i\in I}$ is an indexed collection of sets. The union $\bigcup_{i\in I}A_i$ collects elements present in at least one $A_i$; the intersection $\bigcap_{i\in I}A_i$ collects elements present in every $A_i$. The \emph{power set} $\mathcal P(A)$ is the set of all subsets of $A$. For $A=\{u,v\}$ it is $\{\varnothing,\{u\},\{v\},\{u,v\}\}$. If $A$ has $n$ elements, its power set contains $2^n$ elements: each original element can independently be included or excluded.

\section{The significance for scientific data}
Many preprocessing operations change the set of data included in analysis. If $U$ is the set of detected units, then $U_{\rm quality}$ might retain only units meeting refractory-violation and signal-to-noise thresholds. Estimating a correlation on $U_{\rm quality}$ is not the same question as estimating it on $U$. Record the inclusion criteria as mathematical predicates; this makes sampling assumptions auditable.

\begin{workedexample}[Counting selections]
Suppose four electrodes pass quality control. A subset represents the electrodes chosen for a multichannel analysis; there are $2^4=16$ possible subsets, including the empty selection. Requiring at least one electrode yields 15. Requiring exactly two yields $\binom42=6$. This distinction prepares us for combinatorics and multiple-comparison correction.
\end{workedexample}
% --- expanded edition (October 2026) ---
\section{Algebra of sets}
The operations of the previous sections obey identities that parallel the logical equivalences of Chapter~\ref{ch:logic}. This parallel is not accidental: $x\in A\cap B$ means $x\in A\land x\in B$, $x\in A\cup B$ means $x\in A\lor x\in B$, and $x\in A^c$ means $\neg(x\in A)$. Every set identity can therefore be proved by translating to a logical equivalence about membership.

\begin{proposition}[Laws of the algebra of sets]
For subsets $A,B,C$ of a universe $U$:
\begin{enumerate}[leftmargin=2em]
\item commutativity: $A\cup B=B\cup A$ and $A\cap B=B\cap A$;
\item associativity: $(A\cup B)\cup C=A\cup(B\cup C)$, and similarly for $\cap$;
\item distributivity: $A\cap(B\cup C)=(A\cap B)\cup(A\cap C)$ and $A\cup(B\cap C)=(A\cup B)\cap(A\cup C)$;
\item identity and absorption: $A\cup\varnothing=A$, $A\cap U=A$, $A\cup(A\cap B)=A$;
\item complement: $A\cup A^c=U$, $A\cap A^c=\varnothing$, $(A^c)^c=A$.
\end{enumerate}
\end{proposition}
\begin{proof}
We prove the second distributive law; the others follow the same pattern. Let $x\in U$. Then $x\in A\cup(B\cap C)$ iff $x\in A\lor(x\in B\land x\in C)$. By the logical distributive law this is equivalent to $(x\in A\lor x\in B)\land(x\in A\lor x\in C)$, i.e.\ $x\in(A\cup B)\cap(A\cup C)$. Since the two sets have the same members, they are equal.
\end{proof}

\begin{proposition}[De Morgan for families]
For any indexed family $(A_i)_{i\in I}$ of subsets of $U$,
\[
\Bigl(\bigcup_{i\in I}A_i\Bigr)^c=\bigcap_{i\in I}A_i^c,\qquad
\Bigl(\bigcap_{i\in I}A_i\Bigr)^c=\bigcup_{i\in I}A_i^c.
\]
\end{proposition}
\begin{proof}
For the first identity, $x\notin\bigcup_iA_i$ iff $\neg\exists i\,(x\in A_i)$ iff $\forall i\,(x\notin A_i)$ iff $x\in\bigcap_iA_i^c$. The second follows by applying the first to the family $(A_i^c)$ and taking complements.
\end{proof}

\section{Symmetric difference and indicator functions}
The \emph{symmetric difference} $A\,\triangle\,B=(A\setminus B)\cup(B\setminus A)$ contains the elements in exactly one of the two sets. It measures disagreement between two classifications: if $A$ and $B$ are the spikes assigned to a unit by two sorting algorithms, $|A\triangle B|$ counts the spikes on which they disagree.

The \emph{indicator function} of $A\subseteq U$ is $\mathbf 1_A:U\to\{0,1\}$, with $\mathbf 1_A(x)=1$ if $x\in A$ and $0$ otherwise. Indicators turn set operations into arithmetic:
\[
\mathbf 1_{A\cap B}=\mathbf 1_A\mathbf 1_B,\qquad
\mathbf 1_{A^c}=1-\mathbf 1_A,\qquad
\mathbf 1_{A\cup B}=\mathbf 1_A+\mathbf 1_B-\mathbf 1_A\mathbf 1_B,\qquad
\mathbf 1_{A\triangle B}=|\mathbf 1_A-\mathbf 1_B|.
\]
Summing the union identity over a finite universe gives the \emph{inclusion--exclusion principle} $|A\cup B|=|A|+|B|-|A\cap B|$. Logical masks in numerical software are exactly indicator functions stored as arrays; an element-wise product of two masks is an intersection.

\section{Partitions}
A \emph{partition} of a set $U$ is a family of nonempty, pairwise disjoint subsets whose union is $U$. Trials partitioned by stimulus condition, spikes partitioned by unit identity, and time partitioned into bins are all partitions. If $\{B_1,\ldots,B_k\}$ partitions a finite $U$, then for any $A\subseteq U$,
\[
|A|=\sum_{j=1}^k|A\cap B_j|,
\]
because the sets $A\cap B_j$ are disjoint and their union is $A$. This additivity, generalized from counting to probability and integration, is the root of the law of total probability (Chapter~\ref{ch:probability}).

\section{Russell's paradox in detail}
Suppose there were a set $R$ whose members are exactly the sets that are not members of themselves: $x\in R\Leftrightarrow x\notin x$. Taking $x=R$ yields $R\in R\Leftrightarrow R\notin R$, a contradiction. Hence no such set exists. The \emph{axiom of separation} avoids the paradox by allowing $\{x\in A:P(x)\}$ only for an already given set $A$. Applied to any set $A$, the same argument shows that $\{x\in A:x\notin x\}$ is not a member of $A$; consequently, there is no set of all sets. In applied work this matters less as a danger than as a habit: every set in an analysis should be carved out of a specified universe.

\begin{workedexample}[Masks and exclusion rules]
A session has 400 trials. Let $M$ be trials with movement artifacts (37 trials), $L$ trials with lick contamination (52 trials), and suppose 15 trials have both. The excluded set is $M\cup L$, with $|M\cup L|=37+52-15=74$ trials, leaving $|(M\cup L)^c|=326$. Reporting ``37 and 52 trials were excluded'' would suggest 89 exclusions; inclusion--exclusion gives the correct count.
\end{workedexample}
\begin{exercises}
\item Prove $A\cap(B\cup C)=(A\cap B)\cup(A\cap C)$ from element membership.
\item Determine whether $\varnothing\in\{\varnothing\}$ and whether $\varnothing\subseteq\{\varnothing\}$.
\item Describe the difference between a population of neurons and its sampled subset.
\item Given three data-exclusion rules, express the accepted trials with set operations.
\item Prove $A\setminus(B\cup C)=(A\setminus B)\cap(A\setminus C)$.
\item Show that $A\triangle B=\varnothing$ if and only if $A=B$, and that $\triangle$ is associative.
\item Derive inclusion--exclusion for three sets from the indicator identity.
\item Show that $\{x\in A:x\notin x\}\notin A$ for every set $A$.
\item Two sorters assign 1200 and 1100 spikes to a unit and agree on 1000. Compute $|A\triangle B|$ and a normalized disagreement $|A\triangle B|/|A\cup B|$.
\end{exercises}

\chapter{Ordered pairs, products, and relations}\label{ch:relations}
\begin{readergoal}Understand how ordered observations lead to relations and why correlation, equality, and ordering are mathematically different relations.\end{readergoal}
\section{The order of information matters}
The collection $\{2,7\}$ does not specify which element came first. An \emph{ordered pair} $(2,7)$ does. Formally, an ordered-pair construction must satisfy
\[(a,b)=(c,d)\quad\Longleftrightarrow\quad a=c\ \text{and}\ b=d.\]
One set-theoretic realization is $(a,b)=\{\{a\},\{a,b\}\}$, although the particular encoding is much less important than the stated property. A spike observation might be the pair $(t,u)$ of timestamp and unit identifier; reversing it gives a different type of object.

The \emph{Cartesian product} $A\times B$ contains all ordered pairs $(a,b)$ with $a\in A$ and $b\in B$. If $A$ has $m$ elements and $B$ has $n$, then $A\times B$ has $mn$ elements. More generally, $A_1\times\cdots\times A_k$ defines ordered tuples. A multichannel recording may be indexed by channel $c$, time sample $n$, and trial $r$: its index space is $C\times T\times R$.

\section{What is a relation?}
A \emph{binary relation} from $A$ to $B$ is a subset $R\subseteq A\times B$. Writing $aRb$ is shorthand for $(a,b)\in R$. One relation on the integers is $a\le b$; another is ``$a$ divides $b$.'' In neuroscience, an empirically defined relation may connect two neurons if their estimated spike-count correlation is above a threshold. This makes a graph edge, not necessarily a synapse.

A relation $R\subseteq A\times A$ is \emph{reflexive} if $aRa$ for every $a$, \emph{symmetric} if $aRb$ implies $bRa$, and \emph{transitive} if $aRb$ and $bRc$ imply $aRc$. A relation satisfying all three is an \emph{equivalence relation}; it partitions $A$ into disjoint equivalence classes. A \emph{partial order} is reflexive, antisymmetric, and transitive. Equality is an equivalence relation; subset inclusion is a partial order; temporal precedence is an order relation that requires a convention for simultaneous events.

\section{Composing and inverting relations}
For $R\subseteq A\times B$ and $S\subseteq B\times C$, define
\[S\circ R=\{(a,c):\exists b\in B,\ (a,b)\in R\ \land\ (b,c)\in S\}.\]
The inverse relation is $R^{-1}=\{(b,a):(a,b)\in R\}$. Relation inversion is not yet the same as inversion of a function. For instance, a relation recording which neurons are active during which trials can be reversed to ask which trials contain activity from a neuron.

\begin{workedexample}[Why graph clustering needs care]
Let $uRv$ mean ``estimated Pearson correlation between $u$ and $v$ exceeds 0.5.'' The rule is symmetric whenever correlations are defined, and self-pairs have correlation one for nonconstant variables, but it need not be transitive: neuron 1 may correlate strongly with 2, 2 with 3, and 1 weakly with 3. Thus its connected components are not the same as equivalence classes generated by the direct threshold relation.
\end{workedexample}
% --- expanded edition (October 2026) ---
\section{Equivalence relations and partitions}
Equivalence relations and partitions are two descriptions of the same structure. Given an equivalence relation $\sim$ on $A$, the \emph{equivalence class} of $a$ is $[a]=\{x\in A:x\sim a\}$, and the \emph{quotient set} is $A/{\sim}=\{[a]:a\in A\}$.

\begin{theorem}[Fundamental theorem of equivalence relations]
Let $\sim$ be an equivalence relation on $A$. Then
\begin{enumerate}[leftmargin=2em]
\item every $a\in A$ belongs to its own class $[a]$;
\item for $a,b\in A$, either $[a]=[b]$ or $[a]\cap[b]=\varnothing$;
\item consequently the classes form a partition of $A$.
\end{enumerate}
Conversely, every partition $\mathcal P$ of $A$ defines an equivalence relation by $a\sim b$ iff $a$ and $b$ lie in the same block of $\mathcal P$, and its classes are exactly the blocks.
\end{theorem}
\begin{proof}
(1) follows from reflexivity. For (2), suppose $c\in[a]\cap[b]$, so $c\sim a$ and $c\sim b$. If $x\in[a]$ then $x\sim a$; by symmetry $a\sim c$; by transitivity $x\sim c$ and then $x\sim b$, so $x\in[b]$. Hence $[a]\subseteq[b]$, and the reverse inclusion follows by exchanging the roles of $a$ and $b$. (3) The classes are nonempty by (1), cover $A$ by (1), and are pairwise disjoint or equal by (2). For the converse, reflexivity and symmetry are immediate, and transitivity holds because the blocks are disjoint: if $a,b$ share a block and $b,c$ share a block, both blocks contain $b$ and therefore coincide.
\end{proof}

This theorem underlies all the number-system constructions of Part~II: integers and rationals are equivalence classes of pairs. It also explains why ``spike sorting'' is mathematically a partition problem: an assignment of each spike to exactly one unit is the same as an equivalence relation ``was fired by the same unit as.''

\section{Closures of relations}
A relation that fails to be transitive can be enlarged to a transitive one. The \emph{transitive closure} $R^+$ of $R\subseteq A\times A$ is the smallest transitive relation containing $R$. Concretely, $aR^+b$ iff there is a finite chain $a=x_0Rx_1R\cdots Rx_k=b$ with $k\ge1$. The \emph{reflexive--symmetric--transitive closure} of $R$ is the smallest equivalence relation containing $R$; its classes are the \emph{connected components} of the undirected graph whose edges are the pairs in $R$.

\begin{proposition}
The intersection of any nonempty family of transitive relations on $A$ is transitive. Hence the transitive closure exists and equals the intersection of all transitive relations containing $R$.
\end{proposition}
\begin{proof}
Let $(T_i)_{i\in I}$ be transitive and $T=\bigcap_iT_i$. If $aTb$ and $bTc$, then for each $i$ we have $aT_ib$ and $bT_ic$, hence $aT_ic$. Thus $aTc$. The family of transitive relations containing $R$ is nonempty, since $A\times A$ belongs to it, so its intersection is transitive, contains $R$, and is contained in every transitive relation containing $R$.
\end{proof}

The worked example on correlation thresholds above can now be stated precisely: clustering by connected components replaces the threshold relation by its equivalence closure, which may merge neurons that are not directly correlated at all. Long ``chains'' of moderately correlated units can join two unrelated groups.

\section{Orders}
A \emph{partial order} $\preceq$ on $A$ is reflexive, antisymmetric ($a\preceq b$ and $b\preceq a$ imply $a=b$), and transitive. It is a \emph{total order} if, in addition, any two elements are comparable. The usual $\le$ on numbers is total; inclusion $\subseteq$ on the power set of a set with at least two elements is only partial, since $\{u\}$ and $\{v\}$ are incomparable. A \emph{strict} order $<$ is irreflexive and transitive; from a partial order one obtains a strict one by $a<b\iff a\preceq b\land a\ne b$.

Orders appear in data whenever there is ranking: spike times are totally ordered, but the relation ``neuron $u$ tends to fire before neuron $v$'' estimated from cross-correlograms need not even be transitive. Sequence-detection methods that assume a consistent total order of activation must test, not presuppose, that the estimated precedence relation has that structure.

\begin{workedexample}[Modular phase]
On the real line define $x\sim y$ iff $x-y=2\pi k$ for some integer $k$. This is an equivalence relation: $x-x=0$; if $x-y=2\pi k$ then $y-x=2\pi(-k)$; and $x-y=2\pi k$, $y-z=2\pi m$ give $x-z=2\pi(k+m)$. The quotient $\mathbb R/{\sim}$ is the circle of phases. Each class has exactly one representative in $(-\pi,\pi]$, which is the ``wrapped'' phase returned by \texttt{angle}. Arithmetic on representatives is well defined for addition (adding phases) but not for ordering: ``phase $\alpha$ is less than phase $\beta$'' depends on the representative chosen.
\end{workedexample}
\begin{exercises}
\item List $A\times B$ for $A=\{a,b\}$, $B=\{0,1\}$, and explain why $A\times B\ne B\times A$ generally.
\item Construct a reflexive symmetric relation that is not transitive.
\item Prove the inverse of an equivalence relation is the same relation.
\item Interpret a directed graph edge as a relation; explain why its direction does not prove causal direction.
\item Show that ``$a\equiv b \pmod 3$'' is an equivalence relation on $\mathbb Z$ and list its classes.
\item Compute the transitive closure of $R=\{(1,2),(2,3),(3,4)\}$ on $\{1,2,3,4\}$.
\item Prove that the inverse of a partial order is a partial order.
\item Give a precedence relation among three neurons that is not transitive, and explain why no total order is consistent with it.
\item Let $A$ be a finite set with a partition into $k$ blocks. How many equivalence relations on $A$ have exactly these classes?
\end{exercises}

\chapter{Functions: inputs, outputs, and transformations}\label{ch:functions}
\begin{readergoal}Distinguish relations from functions; work with domain, codomain, image, inverse, and composition.\end{readergoal}
\section{A function is a special relation}
A function $f:A\to B$ assigns to each $a\in A$ exactly one value $f(a)\in B$. It may assign the same output to many inputs. The set $A$ is its \emph{domain}, $B$ its \emph{codomain}, and $f(A)=\{f(a):a\in A\}$ its image or range. A relation that assigns two distinct outputs to the same input is not a function. These distinctions matter when saying ``the neural signal determines the behavior'': an observed relation between activity and behavior is not automatically a deterministic function.

A function is \emph{injective} when different inputs have different outputs, \emph{surjective} when every point in the codomain occurs as an output, and \emph{bijective} when both conditions hold. Only a bijection $f:A\to B$ has a genuine inverse function $f^{-1}:B\to A$. Every relation has an inverse relation, but an arbitrary function does not have an inverse function.

\section{Compositions and transformations}
Given $f:A\to B$ and $g:B\to C$, their composition is $g\circ f:A\to C$, where $(g\circ f)(a)=g(f(a))$. Composition generally does not commute. In electrophysiology, one might take raw voltage $x$, apply a filter $F$, and then a nonlinear amplitude transform $H$: $H(F(x))$. Reversing the order may be meaningless or produce a different result.

\section{Sequences are functions}
A sequence $(x_n)_{n\in\mathbb N}$ is a function from an index set of natural numbers into a set of values. A finite sampled signal is a function $x:\{0,\ldots,N-1\}\to\mathbb R$. A continuously indexed signal is a function $x:I\to\mathbb R$ for some time interval $I$. A multichannel signal is $x:I\to\mathbb R^C$. The difference between continuous and sampled time is not a notational accident: sampling introduces an observation process and potentially destroys information.

\section{Functions of functions}
A \emph{functional} maps functions to numbers. For example, energy over an interval can be defined as $E[x]=\int_I |x(t)|^2\,dt$. A \emph{linear operator} maps functions to functions while respecting addition and scalar multiplication. Filtering, differentiation, and Fourier transformation are operators under suitable assumptions. These ideas will become precise after we construct number systems, vector spaces, and limits.

\begin{workedexample}[A map that loses information]
The square map $q:\mathbb R\to\mathbb R$, $q(x)=x^2$, is not injective: $q(-2)=q(2)=4$. Its inverse relation associates $4$ with both $2$ and $-2$. Restricting the domain to $[0,\infty)$ and codomain to $[0,\infty)$ makes $q$ bijective and its inverse is the nonnegative square root. Signal power similarly discards phase and sign information: it cannot reconstruct the original signal by itself.
\end{workedexample}
% --- expanded edition (October 2026) ---
\section{Images and preimages of sets}
For $f:A\to B$ and $S\subseteq A$, the \emph{image} of $S$ is $f(S)=\{f(a):a\in S\}$; for $T\subseteq B$, the \emph{preimage} is $f^{-1}(T)=\{a\in A:f(a)\in T\}$. The preimage is defined for every function, whether or not $f$ is invertible; the notation $f^{-1}(T)$ does not presuppose an inverse function.

\begin{proposition}[Preimages respect set operations]
For $f:A\to B$ and subsets $T,T_i\subseteq B$:
\[
f^{-1}\Bigl(\bigcup_iT_i\Bigr)=\bigcup_if^{-1}(T_i),\quad
f^{-1}\Bigl(\bigcap_iT_i\Bigr)=\bigcap_if^{-1}(T_i),\quad
f^{-1}(B\setminus T)=A\setminus f^{-1}(T).
\]
Images respect unions, $f(\bigcup_iS_i)=\bigcup_if(S_i)$, but in general only $f(S_1\cap S_2)\subseteq f(S_1)\cap f(S_2)$.
\end{proposition}
\begin{proof}
$a\in f^{-1}(\bigcap_iT_i)$ iff $f(a)\in T_i$ for every $i$ iff $a\in f^{-1}(T_i)$ for every $i$. The union and complement cases are identical in form. For images and intersections, if $b=f(a)$ with $a\in S_1\cap S_2$ then $b\in f(S_1)$ and $b\in f(S_2)$. Equality can fail: with $f(x)=x^2$, $S_1=\{-1\}$, $S_2=\{1\}$, we have $f(S_1\cap S_2)=\varnothing$ but $f(S_1)\cap f(S_2)=\{1\}$.
\end{proof}
This asymmetry explains why measurability (Part~V) is defined through preimages: a function is measurable when preimages of measurable sets are measurable, precisely because preimages commute with all set operations.

\section{Left and right inverses}
A function $g:B\to A$ is a \emph{left inverse} of $f:A\to B$ if $g\circ f=\mathrm{id}_A$, and a \emph{right inverse} if $f\circ g=\mathrm{id}_B$.

\begin{theorem}
Let $A$ be nonempty and $f:A\to B$.
\begin{enumerate}[leftmargin=2em]
\item $f$ has a left inverse if and only if $f$ is injective.
\item $f$ has a right inverse if and only if $f$ is surjective (using the axiom of choice for infinite $B$).
\item If $f$ has both a left inverse $g$ and a right inverse $h$, then $g=h$, and this common function is the unique inverse $f^{-1}$.
\end{enumerate}
\end{theorem}
\begin{proof}
(1) If $g\circ f=\mathrm{id}$ and $f(a)=f(a')$, then $a=g(f(a))=g(f(a'))=a'$. Conversely, if $f$ is injective, fix $a_0\in A$ and define $g(b)$ to be the unique $a$ with $f(a)=b$ when $b\in f(A)$, and $a_0$ otherwise. (2) If $f\circ h=\mathrm{id}$, every $b=f(h(b))$ lies in the image. Conversely, if $f$ is surjective, each preimage $f^{-1}(\{b\})$ is nonempty; choosing one element from each defines $h$. (3) $g=g\circ\mathrm{id}_B=g\circ(f\circ h)=(g\circ f)\circ h=\mathrm{id}_A\circ h=h$.
\end{proof}

A left inverse is a \emph{decoder} that recovers the input whenever the output came from the encoder. A right inverse is a \emph{generator} that produces, for each target output, some input achieving it. In neural decoding, a stimulus-to-response map that is not injective (two stimuli evoke the same response) has no perfect decoder; in neural control, a response space larger than the set of achievable responses has no right inverse.

\section{Composition, injectivity, and surjectivity}
\begin{proposition}
Let $f:A\to B$ and $g:B\to C$.
(i) If $f$ and $g$ are injective, so is $g\circ f$.
(ii) If $f$ and $g$ are surjective, so is $g\circ f$.
(iii) If $g\circ f$ is injective, then $f$ is injective.
(iv) If $g\circ f$ is surjective, then $g$ is surjective.
\end{proposition}
\begin{proof}
(i) $g(f(a))=g(f(a'))$ gives $f(a)=f(a')$ and then $a=a'$. (ii) For $c\in C$ choose $b$ with $g(b)=c$, then $a$ with $f(a)=b$. (iii) If $f(a)=f(a')$ then $g(f(a))=g(f(a'))$, so $a=a'$. (iv) Every $c=g(f(a))$ is a value of $g$.
\end{proof}
Statement (iii) has an important consequence for processing pipelines: if any early stage $f$ discards information (is not injective), no later stage $g$ can restore it. Rectification, thresholding, binning, and taking power are all non-injective stages.

\section{Cardinality and countability}
Sets $A$ and $B$ have the \emph{same cardinality}, $|A|=|B|$, if there is a bijection between them. A set is \emph{countable} if it is finite or has the cardinality of $\mathbb N_0$. The integers are countable via $0,1,-1,2,-2,\ldots$, and so are the rationals and every finite product of countable sets.

\begin{theorem}[Cantor]
For every set $A$ there is no surjection $A\to\mathcal P(A)$. In particular $|A|<|\mathcal P(A)|$.
\end{theorem}
\begin{proof}
Let $f:A\to\mathcal P(A)$ and set $D=\{a\in A:a\notin f(a)\}$. If $D=f(d)$ for some $d$, then $d\in D\iff d\notin f(d)=D$, a contradiction. Hence $D$ is not in the image of $f$.
\end{proof}
Applied to $A=\mathbb N_0$, Cantor's theorem shows that the set of infinite binary sequences (equivalently, the set of all possible infinite spike rasters at fixed resolution) is uncountable. This is one reason why statements about ``all possible spike trains'' require measure theory rather than counting.
\begin{exercises}
\item Decide whether $f(x)=x^3$ from $\mathbb R$ to $\mathbb R$ is bijective.
\item State domain and codomain for a histogram of spike counts in 10-ms bins.
\item Find $f\circ g$ and $g\circ f$ for $f(x)=x+1$ and $g(x)=x^2$.
\item Explain why a spectrogram generally does not contain sufficient information to uniquely recover a waveform without additional phase assumptions.
\item Prove that $f(f^{-1}(T))\subseteq T$ with equality when $f$ is surjective, and $S\subseteq f^{-1}(f(S))$ with equality when $f$ is injective.
\item Give a function with a left inverse but no right inverse, and describe all its left inverses.
\item Show that $g\circ f$ injective does not imply $g$ injective.
\item Construct an explicit bijection $\mathbb N_0\to\mathbb Z$.
\item A spike detector maps a 30\,kHz voltage trace to a list of threshold-crossing times. Identify the domain, codomain, and why the map is not injective.
\end{exercises}

\chapter{Bridges from statements to relations and functions}
\label{ch:bridges-language}
\begin{readergoal}Turn statements in everyday language into small, verifiable arguments about sets, relations, and functions. You need only the preceding language chapters.\end{readergoal}

\noindent\textbf{How to read a bridge.} Each numbered bridge states exactly what is to be proved, spells out the argument, and connects it to a possible representation of recorded data. A small proof is not disposable: the same reasoning returns inside more advanced methods.

\bridgeentry{extensionality}{Set equality is equality of membership}
\begin{proposition}For sets $A,B$, we have $A=B$ if and only if $A\subseteq B$ and $B\subseteq A$.\end{proposition}
\begin{proof}If $A=B$, every element of either set belongs to the other. Conversely, $A\subseteq B$ says that $x\in A$ implies $x\in B$, and $B\subseteq A$ says the reverse. Thus for every object $x$, $x\in A$ exactly when $x\in B$. By extensionality the sets are equal.\end{proof}
\paragraph{Example.} A detector returning timestamps $D$ and a manually verified set $G$ agree exactly when no event belongs only to $D$ and no event belongs only to $G$. Matching counts alone is insufficient.

\bridgeentry{subset-order}{Inclusion is a partial order}
\begin{proposition}The relation $\subseteq$ is reflexive, transitive and antisymmetric on the power set $\mathcal P(U)$.\end{proposition}
\begin{proof}Every element of $A$ belongs to $A$, so $A\subseteq A$. If $A\subseteq B$ and $B\subseteq C$, then any $x\in A$ belongs to $B$ and then $C$, proving transitivity. If $A\subseteq B$ and $B\subseteq A$, Bridge~\ref{bridge:extensionality} gives $A=B$.\end{proof}
\paragraph{Example.} Increasing a permissive spike detection threshold does not necessarily produce nested candidate sets unless the acceptance rule is itself monotone in that threshold; one should first prove monotonicity of the actual rule.

\section{Lean 4 verification: the first two bridges}
\label{sec:lean-bridges-01-02}
\noindent\begin{minipage}[t]{0.44\textwidth}
\textbf{English: BRIDGE-001}

For any sets $A,B$ of elements of a common type, equality is equivalent
to mutual inclusion:
\[A=B\iff(A\subseteq B\land B\subseteq A).\]
If $A=B$, both inclusions follow. Conversely, two-way membership gives
equality by set extensionality.

\medskip\textbf{English: BRIDGE-002}

Set inclusion is reflexive ($A\subseteq A$), transitive
($A\subseteq B\subseteq C\Rightarrow A\subseteq C$), and
antisymmetric ($A\subseteq B\land B\subseteq A\Rightarrow A=B$).
For the final property, reuse BRIDGE-001.

\medskip\textit{Formal correspondence:} \texttt{Set $\alpha$} represents
sets whose members belong to an arbitrary type $\alpha$.
The three order laws are stated explicitly in Lean.
The code is checked separately by the pinned Lean toolchain.
\end{minipage}\hfill
\begin{minipage}[t]{0.53\textwidth}
\textbf{Lean 4 / Mathlib}

\lstinputlisting[language=LeanFour,basicstyle=\ttfamily\tiny,
breaklines=true,columns=fullflexible,keepspaces=true,frame=single]
{../../formal/lean/Con2physFormal/Bridges.lean}
\end{minipage}


\bridgeentry{cartesian}{Cartesian products define possible pairs}
\begin{proposition}For sets $A,B$, the Cartesian product $A\times B$ consists exactly of the ordered pairs $(a,b)$ with $a\in A$ and $b\in B$. If either factor is empty, the product is empty.\end{proposition}
\begin{proof}By definition a pair is in $A\times B$ only when both coordinate conditions hold. When $A=\varnothing$, no first coordinate can satisfy the requirement, and similarly for $B=\varnothing$.\end{proof}
\paragraph{Example.} If $U$ is a set of units and $T$ a recording interval, $U\times T$ describes \emph{possible} unit--time pairs; the observed spike relation is only a subset of this product.

\section{Lean 4 verification: Cartesian products}
\label{sec:lean-bridge-03}
For types $\alpha$ and $\beta$, and sets $A\subseteq\alpha$ and
$B\subseteq\beta$, an ordered pair $(a,b)$ belongs to $A\times B$
precisely when $a\in A$ and $b\in B$. If either set is empty,
the product has no members. These are three distinct Lean theorems,
representing the full statement of Bridge~\ref{bridge:cartesian}.

\lstinputlisting[language=LeanFour,basicstyle=\ttfamily\footnotesize,
breaklines=true,columns=fullflexible,keepspaces=true,frame=single,
firstline=28,lastline=47]
{../../formal/lean/Con2physFormal/Bridges.lean}


\bridgeentry{relation-composition}{Composition of relations is associative}
\begin{proposition}If $R\subseteq A\times B$, $S\subseteq B\times C$, and $T\subseteq C\times D$, then $(T\circ S)\circ R=T\circ(S\circ R)$.\end{proposition}
\begin{proof}A pair $(a,d)$ belongs to the left side exactly when there are $b\in B$ and $c\in C$ such that $aRb$, $bSc$, and $cTd$. The same witnesses $b,c$ characterize the right side. Both relations therefore contain exactly the same pairs.\end{proof}
\paragraph{Example.} A relation between units and channels followed by one between channels and brain regions can be composed. Composition alone does not make either assignment biologically correct.

\section{Lean 4 verification: relational composition}
\label{sec:lean-bridge-04}
For typed relations $R\subseteq X\times Y$, $S\subseteq Y\times Z$
and $T\subseteq Z\times W$, we define \emph{composition} with the
same order as function composition: $\operatorname{relComp}(S,R)$
relates $x$ to $z$ when some $y$ satisfies $xRy$ and $ySz$.
The witnesses for either association are precisely the same $y,z$,
so composition is associative.

\lstinputlisting[language=LeanFour,basicstyle=\ttfamily\footnotesize,
breaklines=true,columns=fullflexible,keepspaces=true,frame=single,
firstline=51,lastline=73]
{../../formal/lean/Con2physFormal/Bridges.lean}


\bridgeentry{equivalence-partition}{Equivalence classes form a partition}
\begin{proposition}If $\sim$ is reflexive, symmetric and transitive on $A$, the classes $[a]=\{b\in A:b\sim a\}$ cover $A$ and any two classes are identical or disjoint.\end{proposition}
\begin{proof}Reflexivity places each $a$ in $[a]$, so the classes cover $A$. If $x\in[a]\cap[b]$, symmetry and transitivity give $a\sim b$. For any $y\sim a$, transitivity gives $y\sim b$, and conversely, hence $[a]=[b]$.\end{proof}
\paragraph{Example.} Grouping timestamps by identical half-open time-bin index defines equivalence classes. Grouping by ``is within 2 ms'' generally fails transitivity and is not an equivalence relation.

\section{Lean 4: equivalence classes form a partition}
\label{sec:lean-bridge-05}
Let $r$ be a reflexive, symmetric, and transitive relation on a
type $\alpha$. We define the class of $a$ by
$[a]=\{b:r(b,a)\}$. Reflexivity proves that every $a$ lies
in its own class, so the classes cover the carrier.

If $x$ lies in both $[a]$ and $[b]$, symmetry followed by
transitivity gives $a\,r\,b$. Consequently every element of
$[a]$ lies in $[b]$, and conversely. Thus intersecting
classes are equal. This establishes the two characteristic
partition properties without assuming that the relation is
equality.

\lstinputlisting[language=LeanFour,basicstyle=\ttfamily\footnotesize,
breaklines=true,columns=fullflexible,keepspaces=true,frame=single,
firstline=73,lastline=97]
{../../formal/lean/Con2physFormal/Bridges.lean}


\bridgeentry{image-preimage}{Images preserve unions; preimages preserve unions and intersections}
\begin{proposition}For $f:A\to B$ and subsets $E,F\subseteq A$, $f(E\cup F)=f(E)\cup f(F)$. For $C,D\subseteq B$, $f^{-1}(C\cap D)=f^{-1}(C)\cap f^{-1}(D)$ and $f^{-1}(C\cup D)=f^{-1}(C)\cup f^{-1}(D)$.\end{proposition}
\begin{proof}A value $y$ belongs to $f(E\cup F)$ exactly when some input in $E$ or $F$ maps to $y$. This is the union of the images. An input $x$ belongs to $f^{-1}(C\cap D)$ exactly when $f(x)$ belongs to both $C$ and $D$. The union case follows by replacing ``both'' with ``at least one''.\end{proof}
\paragraph{Example.} A trial-to-condition mapping permits selecting all trials whose condition lies in a specified set via a preimage. Unlike an inverse function, the preimage is defined even for many-to-one mappings.

\bridgeentry{inverse-functions}{A two-sided inverse is equivalent to bijectivity}
\begin{proposition}A function $f:A\to B$ has a function $g:B\to A$ with $g\circ f=\mathrm{id}_A$ and $f\circ g=\mathrm{id}_B$ if and only if $f$ is injective and surjective.\end{proposition}
\begin{proof}If $f(a_1)=f(a_2)$, applying $g$ gives $a_1=a_2$, so $f$ is injective. For each $b\in B$, $b=f(g(b))$, so it is surjective. Conversely, bijectivity supplies exactly one $a$ for each $b=f(a)$; assigning $g(b)=a$ defines both inverse identities.\end{proof}
\paragraph{Example.} A many-to-one mapping from a continuous waveform to its power spectrum is not invertible: phase information has been discarded.

\bridgeentry{countable-pairs}{Pairs of natural indices can be enumerated}
\begin{proposition}The set $\mathbb N_0\times\mathbb N_0$ is countable.\end{proposition}
\begin{proof}Arrange the pairs along diagonals with fixed sum $s=m+n$: $(0,0)$, then $(0,1),(1,0)$, then $(0,2),(1,1),(2,0)$, and so on. Every pair appears on its finite-sum diagonal. Numbering entries as they are visited gives an enumeration by natural numbers.\end{proof}
\paragraph{Example.} A countable list of possible channel--time-bin indices need not contain every continuous-time event. There is a difference between indexing samples and describing physical time.

\begin{exercises}
\item Construct two sets with equal cardinality that are not equal.
\item Explain why ``within 1 ms'' is not generally transitive.
\item For a many-to-one $f$, exhibit $E,F$ for which $f(E\cap F)\neq f(E)\cap f(F)$.
\end{exercises}
\part{The construction of number systems}
\chapter{Natural numbers, counting, and induction}\label{ch:naturals}
\begin{readergoal}See how counting can be built from basic rules, rather than accepted as an unexplained operation.\end{readergoal}
\section{Counting precedes measurement}
When we say there were seven action potentials, we count discrete events. Counting does not require time units, voltage calibration, or a model of firing probability. The natural numbers are the first formal setting for such statements. We adopt $\mathbb N_0=\{0,1,2,\ldots\}$, reserving $\mathbb N_+=\{1,2,\ldots\}$. Authors differ about including zero, so every technical treatment should make its convention explicit.

\section{Successors and the Peano viewpoint}
Imagine a set $N$ with a distinguished element $0$ and successor operation $S:N\to N$. The basic idea is that zero is not the successor of anything; distinct inputs have distinct successors; and any subset containing zero and closed under the successor operation contains all of $N$. The final clause is the induction principle. These structural properties establish a unique natural-number system up to isomorphism, under the usual axiomatic background \cite{MendelsonNumbers}.

Define addition recursively by $n+0=n$ and $n+S(m)=S(n+m)$. Define multiplication through $n\cdot0=0$ and $n\cdot S(m)=n\cdot m+n$. These are \emph{definitions}, while commutativity, associativity, and distributivity are derived properties. We ordinarily write $n+1$ rather than $S(n)$ after connecting the symbols.

\section{Induction as a way to prove infinite families of claims}
Suppose we wish to prove $\sum_{k=1}^n k=n(n+1)/2$ for all positive integers $n$. The base case $n=1$ holds. Assume it is true for $n$. Then
\[\sum_{k=1}^{n+1}k=\frac{n(n+1)}2+(n+1)=\frac{(n+1)(n+2)}2.\]
Therefore the claim holds for $n+1$; induction establishes every case. Induction is a proof principle, not an empirical inference from observing a finite number of examples.

\section{Finite sets and combinatorics}
Two finite sets have the same cardinality if there exists a bijection between them. The number of length-$n$ binary vectors is $2^n$; the number of ways to choose $k$ distinct objects from $n$ without order is $\binom nk=n!/[k!(n-k)!]$. If 20 channels are inspected in all unordered pairs, the number of pairs is $\binom{20}{2}=190$. This count matters when reporting a family of connectivity tests.

\begin{workedexample}[Discrete time windows]
A recording sampled at $f_s=1000$ samples/s for 10 seconds has $N=f_sT=10000$ samples if we choose the index set $\{0,\ldots,9999\}$. The sample at index $n$ corresponds to $t_n=n/f_s$, so the last recorded time is $9.999$ s, not 10 s. Distinguish the \emph{duration of acquisition} from the \emph{timestamp of its final sample}.
\end{workedexample}
\begin{exercises}
\item Prove by induction that $1+3+\cdots+(2n-1)=n^2$.
\item Count unordered neuron pairs from $N=80$ recorded units.
\item Show that a finite sequence of $N$ entries is a function on $\{0,\ldots,N-1\}$.
\item Explain why mathematical induction and statistical generalization from trials are different operations.
\end{exercises}
\chapter{Integers: extending counting to differences}\label{ch:integers}
\begin{readergoal}Understand negative quantities through equivalence classes and learn what a legitimate algebraic extension must preserve.\end{readergoal}
\section{Why natural numbers are not enough}
A count cannot be negative, but a \emph{difference between counts} can be. We need numbers representing expressions like $3-5$ without insisting that subtraction remain inside the natural numbers. Define a relation on pairs $(a,b)\in\mathbb N_0^2$ by
\[(a,b)\sim(c,d)\quad\Longleftrightarrow\quad a+d=b+c.\]
This is an equivalence relation; its classes represent formal differences $a-b$. For example, $(3,5)\sim(0,2)$ because $3+2=5+0$. We call the set of equivalence classes $\mathbb Z$.

\section{Operations on classes}
Write $[(a,b)]$ for the class of $(a,b)$. Addition can be defined by
\[[(a,b)]+[(c,d)]=[(a+c,b+d)],\]
and multiplication by
\[[(a,b)]\cdot[(c,d)]=[(ac+bd,ad+bc)].\]
A good definition must be \emph{well defined}: choosing different representatives of the same classes must not change the resulting class. This issue appears again in quotient spaces and latent-variable models; an output must not depend on an arbitrary representation of the same underlying object.

\section{Order and absolute value}
Integers inherit addition, multiplication, and a total order. Negative numbers permit signed offsets and differences, but not all integers possess multiplicative inverses within $\mathbb Z$. The absolute value $|n|$ removes sign: $|n|=n$ when $n\ge0$, and $|n|=-n$ otherwise. It satisfies the triangle inequality $|a+b|\le|a|+|b|$.

\section{A neuroscience example}
Spike-count differences $D=N_{\rm post}-N_{\rm pre}$ may be negative although each count is nonnegative. However, raw subtraction is sensitive to the lengths of the counting intervals. If the intervals differ, compare normalized rates $N/T$ rather than interpreting the count difference uncritically. Mathematical legitimacy does not guarantee scientific comparability.

% --- expanded edition (October 2026) ---
\section{Proving that the operations are well defined}
\begin{proposition}
Integer addition and multiplication defined on classes do not depend on the chosen representatives.
\end{proposition}
\begin{proof}
Suppose $(a,b)\sim(a',b')$ and $(c,d)\sim(c',d')$, so $a+b'=b+a'$ and $c+d'=d+c'$. For addition we need $(a+c,b+d)\sim(a'+c',b'+d')$, i.e.\ $a+c+b'+d'=b+d+a'+c'$, which is the sum of the two hypotheses. For multiplication we need
\[
ac+bd+a'd'+b'c'=ad+bc+a'c'+b'd'.
\]
It suffices to change one representative at a time. Replacing $(a,b)$ by $(a',b')$ with $c,d$ fixed, the required identity $ac+bd+a'd+b'c=ad+bc+a'c+b'd$ rearranges to $c(a+b')+d(b+a')=c(b+a')+d(a+b')$, which holds because $a+b'=b+a'$. The second representative is handled symmetrically.
\end{proof}

The map $n\mapsto[(n,0)]$ embeds $\mathbb N_0$ into $\mathbb Z$ injectively and preserves addition and multiplication. Every class $[(a,b)]$ equals either $[(n,0)]$ or $[(0,n)]$ for a unique $n$, which we write $n$ or $-n$. In this sense $\mathbb Z$ \emph{contains} $\mathbb N_0$, even though, strictly, the new numbers are sets of pairs.

\section{The ring axioms}
With these operations, $\mathbb Z$ is a \emph{commutative ring with identity}: addition is associative and commutative with identity $0$ and additive inverses; multiplication is associative and commutative with identity $1$; and multiplication distributes over addition. The ring has no zero divisors: $ab=0$ implies $a=0$ or $b=0$. Consequently cancellation holds: if $ac=bc$ and $c\ne0$, then $a=b$.

\begin{proposition}
In any ring, $0\cdot a=0$ and $(-1)a=-a$. Hence $(-a)(-b)=ab$.
\end{proposition}
\begin{proof}
$0\cdot a=(0+0)a=0\cdot a+0\cdot a$; adding $-(0\cdot a)$ to both sides gives $0=0\cdot a$. Next $a+(-1)a=(1+(-1))a=0\cdot a=0$, so $(-1)a$ is the additive inverse of $a$. Finally $(-a)(-b)=(-1)(-1)ab$, and $(-1)(-1)=-(-1)=1$.
\end{proof}
The rule ``a negative times a negative is positive'' is therefore not a convention but a consequence of distributivity. Any arithmetic in which distributivity holds must satisfy it.

\section{Divisibility and the greatest common divisor}
For integers $a,b$ not both zero, $\gcd(a,b)$ is the largest nonnegative common divisor (equivalently, the positive generator of the ideal $a\\mathbb Z+b\\mathbb Z$).

\begin{theorem}[B\'ezout]
For integers $a,b$ not both zero there exist integers $x,y$ with $ax+by=\gcd(a,b)$.
\end{theorem}
\begin{proof}
Let $I=\{ax+by:x,y\in\mathbb Z\}$, and let $g$ be its smallest positive element (by well ordering). Divide $a$ by $g$: $a=qg+r$ with $0\le r<g$. Then $r=a-qg\in I$, so minimality forces $r=0$; thus $g\mid a$, and likewise $g\mid b$. Any common divisor of $a$ and $b$ divides every element of $I$, in particular $g$. Hence $g=\gcd(a,b)$.
\end{proof}
The Euclidean algorithm computes the gcd by repeated division, using $\gcd(a,b)=\gcd(b,a\bmod b)$. Gcds arise when combining signals with different sampling rates: resampling from 30\,kHz to 1.25\,kHz uses the rational factor $1250/30000=1/24$, and polyphase resampling routines reduce such ratios by their gcd.

\begin{workedexample}[Signed lags]
In a cross-correlogram between spike trains $u$ and $v$, the lag $\tau=t_v-t_u$ takes integer values in units of bins. Negative lags mean $v$ preceded $u$. The set of lags in a window $\{-K,\ldots,K\}$ has $2K+1$ elements, and reversing the roles of $u$ and $v$ maps lag $\tau$ to $-\tau$. A peak at negative lag is thus not a different phenomenon from a peak at positive lag with the neurons swapped; the sign records a convention of which train is the reference.
\end{workedexample}
\begin{exercises}
\item Verify symmetry and transitivity of the equivalence relation on ordered pairs.
\item Show that $[(a,b)]+[(b,a)]$ is the integer zero.
\item Explain why subtraction of two nonnegative power estimates need not be nonnegative.
\item Prove that $\mathbb Z$ has no zero divisors, assuming the corresponding fact for $\mathbb N_0$.
\item Use the Euclidean algorithm to compute $\gcd(30000,1250)$ and find B\'ezout coefficients.
\item Prove that if $c\ne0$ and $ac=bc$ in $\mathbb Z$, then $a=b$.
\item Define $\le$ on $\mathbb Z$ by $[(a,b)]\le[(c,d)]$ iff $a+d\le b+c$, and prove that it is well defined.
\end{exercises}

\chapter{Rational numbers, ratios, and fields}\label{ch:rationals}
\begin{readergoal}Construct fractions and learn why units, denominators, and field axioms matter in quantitative science.\end{readergoal}
\section{Fractions as equivalence classes}
Take integer pairs $(a,b)$ with $b\ne0$. Say $(a,b)\sim(c,d)$ when $ad=bc$. The resulting equivalence class represents the fraction $a/b$. Ratios like $1/2=2/4$ are different representations of the same rational number. Addition and multiplication take their familiar forms
\[\frac ab+\frac cd=\frac{ad+bc}{bd},\qquad \frac ab\cdot\frac cd=\frac{ac}{bd}.\]
Rationals form the field $\mathbb Q$: nonzero values possess multiplicative inverses, and the field operations obey associativity, commutativity, and distributivity. Division by zero is undefined: a rate with zero exposure duration is not an extreme number but an invalid operation.

\section{What does a ratio mean physically?}
Suppose a unit emits 75 spikes over 15 seconds. Its time-averaged rate is $75/15=5$ spikes/s. The numerator and denominator have different roles and units. When a software routine returns a mean across observations, ask what each observation represents and whether every observation has the same weight. The mean of within-trial rates need not equal total count divided by total exposure if trial durations differ.

\section{Rationals are dense but incomplete}
Between any two rational numbers $a<b$, their midpoint $(a+b)/2$ is another rational. Yet the equation $x^2=2$ has no rational solution. To show this, suppose $x=p/q$ is in lowest terms and $p^2=2q^2$. Then $p$ must be even; substituting $p=2k$ shows $q$ also even, contradicting coprimality. Thus the rational numbers have gaps from the viewpoint of limits and geometry.

\begin{workedexample}[Units and normalization]
Compare $\mathrm{spikes}/\mathrm{bin}$ with $\mathrm{spikes}/\mathrm{s}$. If the bin width is $\Delta t$ seconds, the rate estimate is $r=n/\Delta t$ in spikes/s. Changing $\Delta t$ changes the variance and resolution of the estimate even if the underlying experiment is unchanged. This anticipates the bias-variance trade-off in spike-rate estimation.
\end{workedexample}
\begin{exercises}
\item Explain what fails if $b=0$ is allowed in a fraction.
\item Calculate the pooled firing rate across a 5-s trial with 20 spikes and a 15-s trial with 30 spikes; compare with the unweighted average of the two trial rates.
\item Prove that between two unequal rationals there is a rational number.
\end{exercises}
\chapter{The real number system, order, and completeness}\label{ch:reals}
\begin{readergoal}Understand why the real line is more than rational arithmetic, and why continuity and convergence depend on completeness.\end{readergoal}
\section{The real line as an ordered field}
We wish to represent distances, times, and voltages without rational-number gaps. The real numbers $\mathbb R$ are an \emph{ordered field} satisfying a further completeness property. One characterization is the \emph{least-upper-bound property}: every nonempty subset of $\mathbb R$ bounded above has a least upper bound. If $A=\{x\in\mathbb R:x^2<2,\ x\ge0\}$, then $\sup A=\sqrt2$. This value cannot be supplied by rational numbers alone.

A real number may be developed rigorously using Dedekind cuts or equivalence classes of Cauchy sequences of rational numbers; either construction requires proving that the resulting operations and order are well defined. We will use the completeness axiom without adopting one construction as primary. For a fuller foundational treatment, consult \cite{MendelsonNumbers}.

\section{Bounds, supremum, and infimum}
A number $M$ is an upper bound for $A$ if $x\le M$ for every $x\in A$. The supremum is the smallest such bound; it need not belong to the set. For the open interval $(0,1)$, the supremum is 1, although 1 is not a member. The infimum is defined dually. A \emph{maximum} must be a member of the set and greater than or equal to every other member; distinguish maximum from supremum.

\section{Distances and the triangle inequality}
The absolute value defines distance on $\mathbb R$ by $d(x,y)=|x-y|$. Distance is nonnegative, symmetric, zero precisely for equal points, and obeys $d(x,z)\le d(x,y)+d(y,z)$. These four properties define a \emph{metric} on a more general space. A tolerance criterion for a numerical solution is meaningful only with a specified distance, scale, and units.

\section{Interval notation and physical uncertainty}
The interval $[a,b]$ includes endpoints, whereas $(a,b)$ excludes them. An experimental measurement can be modeled as an uncertain value lying in some interval, but the interval should not be mistaken for a statistical confidence interval without an inference procedure. A confidence interval concerns the long-run coverage of an estimator under a model; it is not simply an arbitrary range surrounding a measurement.

\begin{workedexample}[The value of completeness]
Take a sequence of rational approximations $1,1.4,1.41,1.414,\ldots$ to $\sqrt2$. Each term can be rational while its limit is not. A space of numbers meant to support convergence of such sequences must contain the limiting value. The same idea underlies convergence of numerical algorithms for signal processing.
\end{workedexample}
\begin{exercises}
\item Find the supremum, infimum, minimum, and maximum, when defined, for $A=(0,1]$.
\item Prove the reverse triangle inequality $\big||x|-|y|\big|\le |x-y|$.
\item Explain why a bounded measurement record does not imply its generating stochastic process has bounded support.
\end{exercises}
\chapter{Complex numbers: amplitude and phase}\label{ch:complex}
\begin{readergoal}Build the complex plane from real coordinates and understand Euler's formula before meeting Fourier and Hilbert transforms.\end{readergoal}
\section{Pairs with new arithmetic}
A complex number $z=a+ib$ has real part $a$ and imaginary part $b$, where $i^2=-1$. We may regard $z$ as the ordered pair $(a,b)$ of reals, with addition $(a,b)+(c,d)=(a+c,b+d)$ and multiplication $(a,b)(c,d)=(ac-bd,ad+bc)$. The conjugate is $\bar z=a-ib$ and $|z|=\sqrt{z\bar z}=\sqrt{a^2+b^2}$. For $z\ne0$, $z^{-1}=\bar z/|z|^2$.

\section{A plane with angles}
Write $z=r(\cos\theta+i\sin\theta)$ for modulus $r\ge0$ and argument $\theta$. Angles differing by integer multiples of $2\pi$ represent the same direction. Euler's formula $e^{i\theta}=\cos\theta+i\sin\theta$ gives the polar representation $z=re^{i\theta}$. Multiplication adds arguments and multiplies moduli:
\[(r_1e^{i\theta_1})(r_2e^{i\theta_2})=r_1r_2e^{i(\theta_1+\theta_2)}.\]
The periodicity of phase will require circular, rather than ordinary linear, statistics.

\section{Why complex numbers are indispensable for signals}
A sine wave can be written $x(t)=A\cos(2\pi ft+\phi)=\Re\{Ae^{i(2\pi ft+\phi)}\}$. Multiplication by $e^{i\omega t}$ shifts the angular coordinate at a constant rate. The Fourier transform treats a signal as a superposition of such rotating complex exponentials. The complex coefficient encodes two quantities: magnitude and phase. Discarding one generally loses information about the waveform.

\begin{workedexample}[Phase wraps]
The angles $179^\circ$ and $-179^\circ$ differ numerically by $358^\circ$ but lie just $2^\circ$ apart on a circle. Their arithmetic mean is $0^\circ$, which points in the opposite direction. A circular mean instead uses the argument of $e^{i179^\circ}+e^{-i179^\circ}$ and correctly points near $180^\circ$.
\end{workedexample}
\begin{exercises}
\item Compute $(1+i)^2$ and $|1+i|$.
\item Show $|zw|=|z||w|$ using complex conjugation.
\item Express $\cos\theta\cos\phi$ as a combination of complex exponentials and simplify.
\item Explain why phase averages need not be meaningful when oscillatory amplitude is close to zero.
\end{exercises}
\chapter{Bridges between counting and number systems}\label{ch:bridges-numbers}
\begin{readergoal}Understand why integers, rational numbers, real numbers, and complex numbers have the algebraic properties used later in measurement and numerical analysis.\end{readergoal}

\bridgeentry{addition-induction}{Associativity of addition from recursion}
\begin{proposition}If addition on $\mathbb N_0$ is defined by $a+0=a$ and $a+S(b)=S(a+b)$, then $(a+b)+c=a+(b+c)$.\end{proposition}
\begin{proof}Induct on $c$. For $c=0$ both sides equal $a+b$. Assume the identity for $c$. Then $(a+b)+S(c)=S((a+b)+c)=S(a+(b+c))=a+S(b+c)=a+(b+S(c))$.\end{proof}
\paragraph{Example.} Accumulating three disjoint spike-count windows is independent of how we group the additions. It is not independent of whether the windows overlap.

\section{Lean 4: addition associativity by induction}
\label{sec:lean-bridge-09}
The recursive equations $a+0=a$ and $a+S(c)=S(a+c)$ lead to
associativity by induction on the third argument. The base case
simplifies by addition with zero. The successor step rewrites
addition at a successor, then invokes the induction hypothesis.
The formal proof mirrors the English argument in
Bridge~\ref{bridge:addition-induction} rather than merely citing
a library theorem of associativity.

\lstinputlisting[language=LeanFour,basicstyle=\ttfamily\footnotesize,
breaklines=true,columns=fullflexible,keepspaces=true,frame=single]
{../../formal/lean/Con2physFormal/BridgesNumbers.lean}


\bridgeentry{integer-pairs}{Constructing integers from pairs of naturals}
\begin{proposition}For $(a,b),(c,d)\in\mathbb N_0^2$, define $(a,b)\sim(c,d)$ when $a+d=b+c$. This is an equivalence relation. Addition $[(a,b)]+[(c,d)]=[(a+c,b+d)]$ is independent of representatives.\end{proposition}
\begin{proof}Reflexivity and symmetry follow from equality. If $a+d=b+c$ and $c+f=d+e$, adding equations and canceling $c+d$ yields $a+f=b+e$, so transitivity holds. If $a+b'=b+a'$ and $c+d'=d+c'$, adding gives $(a+c)+(b'+d')=(b+d)+(a'+c')$, proving representative-independent addition.\end{proof}
\paragraph{Example.} The signed time difference between events is modeled by an integer when timestamps lie on a sample grid; negative offsets have a meaning not available in the natural numbers.

\bridgeentry{rational-pairs}{Rational arithmetic is independent of representation}
\begin{proposition}For integers $a,c$ and nonzero integers $b,d$, set $(a,b)\sim(c,d)$ when $ad=bc$. Rational addition $[(a,b)]+[(c,d)]=[(ad+bc,bd)]$ is well-defined.\end{proposition}
\begin{proof}If $ab'=a'b$ and $cd'=c'd$, then after multiplying by $dd'$ and $bb'$, respectively, one obtains $(ad+bc)b'd'=(a'd'+b'c')bd$. Thus equivalent input fractions yield equivalent output fractions. The nonzero-denominator condition ensures the construction remains in its domain.\end{proof}
\paragraph{Example.} $1/2$ and $500/1000$ represent the same rational sampling interval; changing the fraction's written form should not change computed sampling times in exact arithmetic.

\bridgeentry{rational-density}{There is a rational between distinct reals}
\begin{proposition}If $a<b$ are real, there exists $q\in\mathbb Q$ with $a<q<b$.\end{proposition}
\begin{proof}Choose $n\in\mathbb N$ with $n(b-a)>1$ using the Archimedean property (proved next). Choose an integer $m$ with $na<m\leq na+1$. Since $m\leq na+1<nb$, the rational $q=m/n$ is strictly between $a$ and $b$.\end{proof}
\paragraph{Example.} Any ideal real-valued threshold can be approximated by rational thresholds, but a recorded voltage has finite instrument resolution and may not distinguish them.

\bridgeentry{archimedean}{Archimedean property from completeness}
\begin{proposition}For every real $x$ there is a natural number $n>x$.\end{proposition}
\begin{proof}Otherwise $\mathbb N_0$ would be nonempty and bounded above. Let $s=\sup\mathbb N_0$. Then $s-1$ is not an upper bound, so some $n>s-1$ belongs to $\mathbb N_0$. Its successor $n+1>s$ also belongs to $\mathbb N_0$, contradicting that $s$ is an upper bound.\end{proof}
\paragraph{Example.} Given any finite desired time resolution $\Delta>0$, a sufficiently large integer sampling-rate multiple can make $1/n<\Delta$ in an ideal mathematical model. Hardware bandwidth is a separate question.

\bridgeentry{supremum-epsilon}{A supremum can be approached from below}
\begin{proposition}If a nonempty set $A\subset\mathbb R$ is bounded above and $s=\sup A$, then for every $\varepsilon>0$ some $a\in A$ satisfies $s-\varepsilon<a\leq s$.\end{proposition}
\begin{proof}If no such $a$ existed, every element of $A$ would be at most $s-\varepsilon$. Then $s-\varepsilon$ would be an upper bound smaller than the least upper bound $s$, a contradiction.\end{proof}
\paragraph{Example.} The maximal attainable value of an optimization criterion can be a supremum without being reached by any tested finite set of hyperparameters.

\bridgeentry{floor-bins}{Floor partitions nonnegative time into half-open bins}
\begin{proposition}For $\Delta>0$ and $t\geq0$, there is a unique $k\in\mathbb N_0$ with $k\Delta\leq t<(k+1)\Delta$, namely $k=\lfloor t/\Delta\rfloor$.\end{proposition}
\begin{proof}The integer-part property gives $k\leq t/\Delta<k+1$. Multiplying by positive $\Delta$ yields the interval condition. Two distinct integers cannot satisfy the inequalities because their half-open intervals do not overlap.\end{proof}
\paragraph{Example.} The timestamp $t=0.100$ s belongs to the second $0.1$-s bin $[0.1,0.2)$ if timestamps and boundaries are exact. Floating-point computations require explicit endpoint handling.

\bridgeentry{complex-modulus}{Conjugation makes complex magnitude multiplicative}
\begin{proposition}For $z,w\in\mathbb C$, $\overline{zw}=\bar z\bar w$ and $|zw|=|z||w|$.\end{proposition}
\begin{proof}Expanding $(a+ib)(c+id)$ and reversing the sign of the imaginary part proves the conjugation identity. Since $|z|^2=z\bar z$, we have $|zw|^2=zw\overline{zw}=(z\bar z)(w\bar w)=|z|^2|w|^2$. Taking nonnegative square roots establishes the result.\end{proof}
\paragraph{Example.} Complex Fourier coefficients encode both amplitude and phase; replacing every coefficient by its modulus loses phase information even though magnitudes multiply neatly.

\begin{exercises}
\item Prove cancellation in the natural numbers using induction.
\item Exhibit two distinct integer pairs representing the same integer.
\item Show that the union of two adjacent half-open time bins is again a half-open interval.
\end{exercises}
\part{Calculus and analysis}
\chapter{Sequences, limits, and continuity}\label{ch:limits}
\begin{readergoal}Understand the precision behind ``approaches'', ``converges'', and ``continuous'' using explicit tolerances.\end{readergoal}
\section{An infinite list is a function}
The sequence $x_n=1/n$ approaches zero as $n$ grows. But what should ``approaches'' mean? We say $x_n\to L$ if, for every tolerance $\varepsilon>0$, there exists an index $N_\varepsilon$ such that $|x_n-L|<\varepsilon$ whenever $n\ge N_\varepsilon$. The order of the quantifiers matters. We must succeed for \emph{every} tolerance, although the required index may depend on the chosen tolerance.

For $x_n=1/n$, choose $N_\varepsilon>1/\varepsilon$. Then $n\ge N_\varepsilon$ implies $|1/n|<\varepsilon$. Convergence is not the same as attaining the limit; no term in this sequence equals zero. Nor is convergence established by drawing the first few points. An infinite claim requires an argument that controls all sufficiently late terms.

\section{Cauchy sequences and completeness}
A sequence is \emph{Cauchy} if its sufficiently late terms become arbitrarily close to one another: $\forall\varepsilon>0$ there exists $N$ such that $|x_m-x_n|<\varepsilon$ for all $m,n\ge N$. In a complete metric space every Cauchy sequence converges to a point of the space. The rational approximations to $\sqrt2$ form a Cauchy sequence of rationals but fail to converge \emph{within} the rationals. This shows why completeness is essential in analysis.

\section{Functions of a real variable}
We say $f(x)$ tends to $L$ as $x\to a$ if for every $\varepsilon>0$ we can find $\delta>0$ such that $0<|x-a|<\delta$ implies $|f(x)-L|<\varepsilon$. A function is continuous at $a$ when its limit equals its value at that point. Uniform continuity strengthens the statement by demanding one $\delta$ for a given $\varepsilon$ that works across the entire domain.

\section{Why this matters for continuous-time recordings}
A continuous-time model $x(t)$ does not mean its actual neuronal source changes smoothly; it means the mathematical function has a value for each real-valued time in its domain. A sampled waveform records only certain time points. Interpolating between them requires assumptions about smoothness or bandwidth. Similarly, increasing the sampling frequency is not a mathematical proof that a measurement device approximates the underlying field arbitrarily well: instrument bandwidth and noise also matter.

\begin{workedexample}[Finite-difference derivatives]
For a smooth $x(t)$, the quantity $[x(t+h)-x(t)]/h$ may approach a derivative as $h\to0$. In recorded data, however, taking ever smaller $h$ may amplify measurement noise. Convergence of an ideal mathematical expression and consistency of a noisy numerical estimator are separate questions.
\end{workedexample}
% --- expanded edition (October 2026) ---
\section{Algebra of limits}
Proofs from the $\varepsilon$--$N$ definition become routine once a few general rules are established.

\begin{proposition}[Uniqueness and boundedness]
A convergent sequence has exactly one limit and is bounded.
\end{proposition}
\begin{proof}
If $x_n\to L$ and $x_n\to L'$ with $L\ne L'$, take $\varepsilon=|L-L'|/2$. For large $n$ both $|x_n-L|<\varepsilon$ and $|x_n-L'|<\varepsilon$, so $|L-L'|<2\varepsilon=|L-L'|$, a contradiction. For boundedness, choose $N$ with $|x_n-L|<1$ for $n\ge N$; then $|x_n|\le\max\{|x_1|,\ldots,|x_{N-1}|,|L|+1\}$ for all $n$.
\end{proof}

\begin{theorem}[Limit laws]
If $x_n\to a$ and $y_n\to b$, then $x_n+y_n\to a+b$, $x_ny_n\to ab$, and, if $b\ne0$, $x_n/y_n\to a/b$ (for $n$ large enough that $y_n\ne0$).
\end{theorem}
\begin{proof}
Sum: given $\varepsilon$, choose $N$ so that $|x_n-a|<\varepsilon/2$ and $|y_n-b|<\varepsilon/2$ for $n\ge N$, and use the triangle inequality. Product: write
\[
x_ny_n-ab=x_n(y_n-b)+b(x_n-a).
\]
Since $(x_n)$ is bounded by some $M$, $|x_ny_n-ab|\le M|y_n-b|+|b||x_n-a|$, which can be made smaller than $\varepsilon$. Quotient: it suffices to show $1/y_n\to1/b$. For large $n$, $|y_n|>|b|/2$, so $|1/y_n-1/b|=|b-y_n|/(|y_n||b|)\le2|y_n-b|/|b|^2\to0$.
\end{proof}

\begin{proposition}[Squeeze]
If $x_n\le y_n\le z_n$ for all large $n$ and $x_n\to L$, $z_n\to L$, then $y_n\to L$.
\end{proposition}
\begin{proof}
For large $n$, $L-\varepsilon<x_n\le y_n\le z_n<L+\varepsilon$.
\end{proof}

\section{Monotone sequences and the Bolzano--Weierstrass theorem}
\begin{theorem}[Monotone convergence]
A nondecreasing sequence that is bounded above converges to its supremum.
\end{theorem}
\begin{proof}
Let $s=\sup_nx_n$, which exists by completeness. Given $\varepsilon>0$, the approximation property gives $N$ with $x_N>s-\varepsilon$. For $n\ge N$, monotonicity gives $s-\varepsilon<x_N\le x_n\le s$.
\end{proof}

\begin{theorem}[Bolzano--Weierstrass]
Every bounded real sequence has a convergent subsequence.
\end{theorem}
\begin{proof}
Call an index $m$ a \emph{peak} if $x_m\ge x_n$ for all $n>m$. If there are infinitely many peaks $m_1<m_2<\cdots$, then $x_{m_1}\ge x_{m_2}\ge\cdots$ is a bounded nonincreasing subsequence. Otherwise, after the last peak every index $n_1$ has a later $n_2$ with $x_{n_2}>x_{n_1}$, and repeating yields an increasing subsequence. In both cases, monotone convergence applies.
\end{proof}

\begin{theorem}[Cauchy criterion]
A real sequence converges if and only if it is Cauchy.
\end{theorem}
\begin{proof}
If $x_n\to L$, then $|x_m-x_n|\le|x_m-L|+|L-x_n|<\varepsilon$ for large $m,n$. Conversely, a Cauchy sequence is bounded (as in the boundedness proof), so by Bolzano--Weierstrass a subsequence converges to some $L$. Given $\varepsilon$, choose $N$ with $|x_m-x_n|<\varepsilon/2$ for $m,n\ge N$ and a subsequence index $n_k\ge N$ with $|x_{n_k}-L|<\varepsilon/2$; then $|x_n-L|<\varepsilon$ for $n\ge N$.
\end{proof}
The Cauchy criterion is what iterative numerical methods actually check: a stopping rule such as ``stop when successive iterates differ by less than $10^{-8}$'' is a finite surrogate of the Cauchy condition. It is only a surrogate. The harmonic partial sums $H_n=\sum_{k\le n}1/k$ have successive differences $1/n\to0$ but diverge, so small steps alone do not prove convergence.

\section{Series}
A series $\sum_{k=1}^\infty a_k$ converges if its partial sums $S_n=\sum_{k\le n}a_k$ converge. The geometric series $\sum_{k\ge0}r^k$ converges to $1/(1-r)$ iff $|r|<1$, since $S_n=(1-r^{n+1})/(1-r)$. A series converges \emph{absolutely} if $\sum|a_k|<\infty$; absolute convergence implies convergence (by the Cauchy criterion and the triangle inequality) and allows rearrangement without changing the sum. The alternating harmonic series $\sum(-1)^{k+1}/k$ converges, but not absolutely, and can be rearranged to converge to any prescribed real number. The comparison test states that $0\le a_k\le b_k$ and $\sum b_k<\infty$ imply $\sum a_k<\infty$.

Absolute summability is the stability condition for discrete filters (Chapter~\ref{ch:fourier}): if $\sum_k|h_k|<\infty$ and $|x_n|\le M$, then $|\sum_kh_kx_{n-k}|\le M\sum_k|h_k|$.

\section{Continuity on closed intervals}
\begin{theorem}[Intermediate value theorem]
If $f:[a,b]\to\mathbb R$ is continuous and $f(a)<c<f(b)$, there is $x\in(a,b)$ with $f(x)=c$.
\end{theorem}
\begin{proof}
Let $x=\sup\{t\in[a,b]:f(t)<c\}$. If $f(x)<c$, continuity keeps $f<c$ slightly to the right of $x$, contradicting the supremum; $x<b$ because $f(b)>c$. If $f(x)>c$, continuity keeps $f>c$ on an interval $(x-\delta,x]$, so $x-\delta$ is a smaller upper bound, a contradiction. Hence $f(x)=c$.
\end{proof}

\begin{theorem}[Extreme value theorem]
A continuous function on a closed bounded interval is bounded and attains its maximum and minimum.
\end{theorem}
\begin{proof}
If $f$ were unbounded, there would be $x_n$ with $|f(x_n)|>n$. A subsequence converges to some $x\in[a,b]$, and continuity forces $f(x_{n_k})\to f(x)$, contradicting unboundedness. Let $M=\sup f$; choose $x_n$ with $f(x_n)>M-1/n$, extract a convergent subsequence with limit $x^\ast$, and conclude $f(x^\ast)=M$ by continuity.
\end{proof}
Both theorems fail on open intervals or for discontinuous functions: $f(x)=1/x$ on $(0,1]$ is unbounded, and a step function skips intermediate values. A threshold-crossing detector implicitly relies on the intermediate value theorem for the underlying continuous signal; on samples, the crossing time must be interpolated.

\begin{workedexample}[Bisection as a constructive proof]
To find a zero of a continuous $f$ with $f(a)<0<f(b)$, evaluate $f$ at the midpoint and keep the half-interval on which the sign changes. After $n$ steps the interval has length $(b-a)2^{-n}$ and, by the nested interval property, the endpoints converge to a zero. To locate a threshold crossing between two samples to $1\,\mu$s when the samples are $33\,\mu$s apart requires $\lceil\log_2 33\rceil=6$ bisection steps on an interpolated signal. The accuracy refers to the interpolant, not to the unobserved voltage.
\end{workedexample}
\begin{exercises}
\item Prove $2/n\to0$ from the definition.
\item Give a sequence that is bounded but does not converge.
\item Determine whether the step function $f(x)=0$ for $x<0$, $f(x)=1$ for $x\ge0$ is continuous at zero.
\item Explain the conceptual difference between a limit as $n\to\infty$ and a sample size $n$ that remains finite.
\item Prove from the definition that $(n+1)/(2n+3)\to1/2$.
\item Show that $x_n=(1+1/n)^n$ is increasing and bounded above by $3$, and conclude that it converges.
\item Prove that $\sum_{k\ge1}1/k^2$ converges by comparison with $\sum 1/(k(k-1))$ for $k\ge2$.
\item Show that a polynomial of odd degree has a real root.
\item Give a continuous function on $(0,1)$ that is bounded but does not attain its supremum.
\end{exercises}

\chapter{Derivatives, integrals, and differential equations}\label{ch:calculus}
\begin{readergoal}Build rate of change and accumulation from limits and relate them to state variables and observed neural signals.\end{readergoal}
\section{A derivative is a limit of ratios}
For a differentiable function $f$, the derivative at $t$ is
\[f'(t)=\lim_{h\to0}\frac{f(t+h)-f(t)}h.\]
Its units are output units divided by input units: differentiating voltage with respect to time gives volts per second. A derivative expresses an instantaneous rate of change under a continuous model; a finite difference merely estimates this rate using discrete observations. The distinction becomes serious for noisy measurements.

\section{Integration as accumulation}
A Riemann integral $\int_a^b f(t)\,dt$ is a limit of sums of function values multiplied by small interval widths. If $r(t)$ is a firing intensity in spikes/s, its integral over $[a,b]$ has units of expected spikes, and under a conditional-intensity point-process model it represents integrated intensity, not the realized count. If $p(t)$ is power in an appropriate density scale, integrating it over time yields an accumulated quantity with corresponding units.

The fundamental theorem of calculus connects these operations: if $F(t)=\int_a^t f(s)\,ds$ and $f$ is continuous, then $F'(t)=f(t)$. Thus differentiation and integration are linked through limiting constructions.

\section{Differential equations: rules for evolving states}
A differential equation specifies how a variable changes, not necessarily its value at every time. Consider
\[\tau\frac{dx}{dt}=-x+u(t).\]
With $u(t)=u_0$ constant, the solution is $x(t)=u_0+[x(0)-u_0]e^{-t/\tau}$. The time constant $\tau>0$ controls relaxation. Without an initial condition, many trajectories satisfy the same differential equation. This is a model for first-order filtering and a building block for membrane and firing-rate dynamics.

\section{A distinction between models and data}
Recorded spikes are events and need not be represented by a differentiable function. The firing-rate function is an inferred or modeled quantity. Writing $\dot r(t)$ for the derivative of an estimated rate assumes a degree of smoothness that the raw point process does not possess. We will state these choices explicitly when we develop state-space models.

% --- expanded edition (October 2026) ---
\section{Rules of differentiation}
\begin{proposition}[Differentiability implies continuity]
If $f$ is differentiable at $t$, it is continuous at $t$.
\end{proposition}
\begin{proof}
$f(t+h)-f(t)=h\cdot\frac{f(t+h)-f(t)}h\to0\cdot f'(t)=0$ as $h\to0$.
\end{proof}
The converse is false: $|t|$ is continuous but not differentiable at $0$.

\begin{theorem}[Product and chain rules]
If $f,g$ are differentiable at $t$, then $(fg)'(t)=f'(t)g(t)+f(t)g'(t)$. If $g$ is differentiable at $t$ and $f$ at $g(t)$, then $(f\circ g)'(t)=f'(g(t))g'(t)$.
\end{theorem}
\begin{proof}
For the product,
\[
\frac{f(t+h)g(t+h)-f(t)g(t)}h=\frac{f(t+h)-f(t)}hg(t+h)+f(t)\frac{g(t+h)-g(t)}h,
\]
and $g(t+h)\to g(t)$ by continuity. For the chain rule, define $\phi(k)=[f(g(t)+k)-f(g(t))]/k$ for $k\ne0$ and $\phi(0)=f'(g(t))$; $\phi$ is continuous at $0$. With $k=g(t+h)-g(t)$, we have $f(g(t+h))-f(g(t))=\phi(k)k$ in all cases, so the difference quotient equals $\phi(k)\,[g(t+h)-g(t)]/h\to f'(g(t))g'(t)$.
\end{proof}

\section{The mean value theorem}
\begin{theorem}[Rolle and mean value]
Let $f$ be continuous on $[a,b]$ and differentiable on $(a,b)$. If $f(a)=f(b)$, there is $c\in(a,b)$ with $f'(c)=0$. In general there is $c\in(a,b)$ with
\[
f'(c)=\frac{f(b)-f(a)}{b-a}.
\]
\end{theorem}
\begin{proof}
By the extreme value theorem $f$ attains a maximum and a minimum on $[a,b]$. If both occur at endpoints, $f$ is constant and $f'\equiv0$. Otherwise an extremum occurs at an interior $c$, where the one-sided difference quotients have opposite signs, forcing $f'(c)=0$. For the general case apply Rolle's theorem to $g(t)=f(t)-\frac{f(b)-f(a)}{b-a}(t-a)$.
\end{proof}
Consequences: a function with $f'\equiv0$ on an interval is constant; $f'>0$ implies $f$ is strictly increasing; and if $|f'|\le K$ then $|f(x)-f(y)|\le K|x-y|$ (a \emph{Lipschitz} bound). The last statement bounds how fast a signal can change given a bound on its slope, which links bandwidth to sample-to-sample variation.

\section{Taylor's theorem and the error of finite differences}
\begin{theorem}[Taylor with Lagrange remainder]
If $f$ is $n+1$ times differentiable on an interval containing $t$ and $t+h$, then
\[
f(t+h)=\sum_{k=0}^n\frac{f^{(k)}(t)}{k!}h^k+\frac{f^{(n+1)}(\xi)}{(n+1)!}h^{n+1}
\]
for some $\xi$ between $t$ and $t+h$.
\end{theorem}
The proof applies Rolle's theorem repeatedly to the difference between $f$ and its Taylor polynomial. Two consequences matter for sampled signals. The forward difference satisfies
\[
\frac{f(t+h)-f(t)}h=f'(t)+\tfrac h2f''(\xi),
\]
an $O(h)$ error, whereas the central difference $[f(t+h)-f(t-h)]/(2h)=f'(t)+O(h^2)$. If each sample carries independent noise of standard deviation $\sigma$, the forward difference adds noise of standard deviation $\sqrt2\sigma/h$. The total error $\approx\frac h2|f''|+\sqrt2\sigma/h$ is minimized at $h^\ast=\sqrt{2\sqrt2\sigma/|f''|}$: shrinking $h$ below this value makes the derivative estimate \emph{worse}.

\section{The exponential function}
Define $\exp(x)=\sum_{k\ge0}x^k/k!$, which converges absolutely for every real or complex $x$ by comparison with a geometric series once $k>2|x|$. Termwise multiplication of absolutely convergent series and the binomial theorem give $\exp(x+y)=\exp(x)\exp(y)$, and termwise differentiation (justified by uniform convergence on bounded sets) gives $\exp'=\exp$. Writing $e^x=\exp(x)$ and splitting the series for $e^{i\theta}$ into even and odd powers yields
\[
e^{i\theta}=\sum_{m\ge0}\frac{(-1)^m\theta^{2m}}{(2m)!}+i\sum_{m\ge0}\frac{(-1)^m\theta^{2m+1}}{(2m+1)!}=\cos\theta+i\sin\theta,
\]
which proves Euler's formula once cosine and sine are defined by these series.

\section{Linear first-order equations in general}
\begin{theorem}
Let $\tau>0$ and $u$ be continuous. The unique solution of $\tau\dot x=-x+u(t)$ with $x(0)=x_0$ is
\[
x(t)=x_0e^{-t/\tau}+\frac1\tau\int_0^te^{-(t-s)/\tau}u(s)\,ds.
\]
\end{theorem}
\begin{proof}
Multiply by the integrating factor $e^{t/\tau}/\tau$: $\frac d{dt}\bigl(e^{t/\tau}x\bigr)=e^{t/\tau}u(t)/\tau$. Integrate from $0$ to $t$ and multiply by $e^{-t/\tau}$. If $x,y$ are two solutions, $w=x-y$ satisfies $\tau\dot w=-w$, $w(0)=0$, so $\frac d{dt}(e^{t/\tau}w)=0$ and $w\equiv0$.
\end{proof}
The solution is a \emph{convolution} of the input with the exponential kernel $k(t)=\tau^{-1}e^{-t/\tau}\mathbf 1_{t\ge0}$. A leaky integrator, an RC membrane, and an exponential synaptic filter are the same mathematical object. Its step response reaches $1-e^{-1}\approx63\%$ of its final value after one time constant.

\begin{workedexample}[Membrane time constant from a current step]
A passive membrane obeys $C\dot V=-(V-E_L)/R+I(t)$, i.e.\ $\tau\dot V=-(V-E_L)+RI$ with $\tau=RC$. For a step $I(t)=I_0$ at $t=0$ from rest, $V(t)=E_L+RI_0(1-e^{-t/\tau})$. With $R=100\,\mathrm{M\Omega}$ and $C=200$\,pF, $\tau=20$\,ms and a 50\,pA step produces a 5\,mV steady depolarization. Fitting $\tau$ from data presumes the passive model; active conductances, space clamp errors, and series resistance change the curve.
\end{workedexample}
\begin{exercises}
\item Differentiate $f(t)=e^{-t/\tau}$ and confirm the result has units of inverse time.
\item Solve $dx/dt=-ax$ for $a>0$ and a specified initial condition.
\item Explain why a spike-count cumulative function is not ordinarily differentiable at spike times.
\item Compute $\int_0^T 3\,dt$ and interpret it for a rate of 3 spikes/s.
\item Use the mean value theorem to prove $|\sin x-\sin y|\le|x-y|$.
\item Show that the central difference has error $O(h^2)$ using Taylor's theorem to third order.
\item Solve $\tau\dot x=-x+\sin(\omega t)$ with $x(0)=0$ and find the steady-state amplitude and phase lag.
\item Show that $f(t)=t^2\sin(1/t)$ (with $f(0)=0$) is differentiable at $0$ but $f'$ is not continuous there.
\item Derive the optimal step $h^\ast$ for the noisy forward difference.
\end{exercises}

\chapter{Several variables, optimization, and geometry}\label{ch:multi}
\begin{readergoal}Generalize derivatives to gradients and Jacobians and understand curvature, critical points, and constrained optimization.\end{readergoal}
\section{The gradient is a local direction of change}
For $f:\mathbb R^n\to\mathbb R$, the gradient is $\nabla f=(\partial f/\partial x_1,\ldots,\partial f/\partial x_n)^\top$. Its dot product with a direction $v$ gives the directional derivative $D_vf=\nabla f\cdot v$ for a differentiable function. In a neighborhood of $x$, $f(x+h)\approx f(x)+\nabla f(x)^\top h$. The gradient depends on the units and coordinate scaling of the input variables.

A vector-valued map $F:\mathbb R^n\to\mathbb R^m$ has Jacobian $J_{ij}=\partial F_i/\partial x_j$. Near a point, $F(x+h)\approx F(x)+J_F(x)h$. The chain rule for $G\circ F$ is $J_{G\circ F}(x)=J_G(F(x))J_F(x)$, provided the dimensions match.

\section{Second derivatives and the Hessian}
The Hessian of a twice differentiable scalar function is the matrix of second partial derivatives $H_{ij}=\partial^2f/(\partial x_i\partial x_j)$. At a stationary point $\nabla f=0$, a positive-definite Hessian gives a strict local minimum; an indefinite Hessian indicates a saddle. The quadratic approximation
\[f(x+h)\approx f(x)+\nabla f(x)^\top h+\tfrac12 h^\top H_f(x)h\]
provides a local account of curvature and motivates Newton-type optimization.

\section{Optimization as fitting a model}
Least squares seeks $\hat\theta=\arg\min_\theta\sum_i[y_i-f_\theta(x_i)]^2$. The objective is a chosen measure of mismatch, not a guarantee that the parameter has a physiological meaning. Regularization changes the optimization problem, e.g.
\[\hat\theta=\arg\min_\theta\left\{\|y-X\theta\|_2^2+\lambda\|\theta\|_2^2\right\}.\]
The penalty limits parameter magnitude and can improve stability, but also introduces bias. Constraints such as $\theta\ge0$ are meaningful only if supported by the model and units.

\section{The manifold idea starts locally}
A smooth curve in $\mathbb R^3$ may be one-dimensional although it occupies a three-dimensional ambient space. Near a regular point, its variation can be approximated by a tangent direction. General smooth manifolds are spaces locally modeled on $\mathbb R^d$; the ambient dimension and intrinsic dimension need not coincide. PCA identifies a linear subspace, but whether the data genuinely concentrate around a nonlinear manifold is an additional question.

% --- expanded edition (October 2026) ---
\section{Differentiability in several variables}
Partial derivatives alone do not make a function well behaved. The function $f(x,y)=xy/(x^2+y^2)$ with $f(0,0)=0$ has both partial derivatives equal to zero at the origin, yet is not even continuous there (along $y=x$ it equals $1/2$). The correct notion is \emph{(Fr\'echet) differentiability}: $f:\mathbb R^n\to\mathbb R$ is differentiable at $x$ if there is a vector $g$ with
\[
\lim_{h\to0}\frac{|f(x+h)-f(x)-g^\top h|}{\|h\|}=0,
\]
in which case $g=\nabla f(x)$. A useful sufficient condition is that all partial derivatives exist near $x$ and are continuous at $x$. When second partials are continuous, mixed partials commute (Schwarz--Clairaut), so the Hessian is symmetric.

\begin{proposition}[Steepest ascent]
If $f$ is differentiable at $x$ and $\nabla f(x)\ne0$, then among unit vectors $v$ the directional derivative $D_vf(x)$ is maximized by $v=\nabla f(x)/\|\nabla f(x)\|$, with maximal value $\|\nabla f(x)\|$.
\end{proposition}
\begin{proof}
By the Cauchy--Schwarz inequality, $D_vf=\nabla f^\top v\le\|\nabla f\|\|v\|=\|\nabla f\|$, with equality iff $v$ is a positive multiple of $\nabla f$.
\end{proof}

\section{Convexity}
A set $C$ is convex if it contains the segment between any two of its points. A function $f$ on a convex set is \emph{convex} if
\[
f(\lambda x+(1-\lambda)y)\le\lambda f(x)+(1-\lambda)f(y)\qquad(0\le\lambda\le1).
\]
\begin{theorem}
Let $f$ be twice continuously differentiable on an open convex set. Then $f$ is convex iff its Hessian is positive semidefinite everywhere. If $f$ is convex and differentiable, every stationary point is a global minimizer.
\end{theorem}
\begin{proof}[Proof of the second statement]
For differentiable convex $f$, the convexity inequality rearranges to $[f(x+\lambda(y-x))-f(x)]/\lambda\le f(y)-f(x)$; letting $\lambda\downarrow0$ gives the \emph{supporting hyperplane inequality} $f(y)\ge f(x)+\nabla f(x)^\top(y-x)$. If $\nabla f(x)=0$, then $f(y)\ge f(x)$ for every $y$.
\end{proof}
Convexity is why some fitting problems are reliable and others are not. Least squares, ridge regression, and the Poisson GLM with canonical log link (Chapter~\ref{ch:glm}) have convex negative log-likelihoods, so any local optimizer that converges finds the global optimum. Fitting a mixture model, a neural network, or a tensor decomposition is nonconvex; different initializations may give different answers, and agreement across restarts is evidence, not proof, of a global optimum.

\section{Constrained optimization and Lagrange multipliers}
\begin{theorem}[Lagrange multipliers, one constraint]
Let $f,g$ be continuously differentiable and suppose $x^\ast$ is a local extremum of $f$ subject to $g(x)=0$, with $\nabla g(x^\ast)\ne0$. Then $\nabla f(x^\ast)=\lambda\nabla g(x^\ast)$ for some scalar $\lambda$.
\end{theorem}
The geometric idea: on the constraint surface, $f$ cannot increase to first order in any tangent direction, so $\nabla f$ has no component tangent to the surface and must be parallel to its normal $\nabla g$. The PCA derivation in Chapter~\ref{ch:pca} is exactly this theorem with $f(v)=v^\top Sv$ and $g(v)=v^\top v-1$. The multiplier has a meaning: it is the rate at which the optimal value changes as the constraint level changes. In PCA it equals the captured variance.

\section{Gradient descent and Newton's method}
To minimize $f$, \emph{gradient descent} iterates $x_{k+1}=x_k-\eta\nabla f(x_k)$. For the quadratic $f(x)=\tfrac12x^\top Ax$ with $A$ symmetric positive definite and eigenvalues $0<\mu\le\cdots\le L$, the iteration is $x_{k+1}=(I-\eta A)x_k$, which converges iff $0<\eta<2/L$; the best fixed step $\eta=2/(\mu+L)$ contracts errors by the factor $(\kappa-1)/(\kappa+1)$ per step, with condition number $\kappa=L/\mu$. Badly scaled variables (e.g.\ one regressor in volts and another in microvolts) produce large $\kappa$ and slow convergence; standardizing regressors is therefore a numerical as well as a statistical decision.

\emph{Newton's method} uses the quadratic model: $x_{k+1}=x_k-H_f(x_k)^{-1}\nabla f(x_k)$. On a quadratic it converges in one step; near a nondegenerate minimum it converges quadratically (the number of correct digits roughly doubles each step). Its cost is forming and solving with the Hessian. Iteratively reweighted least squares for GLMs is Newton's method applied to the log-likelihood.

\begin{workedexample}[A constrained variance problem]
Maximize $f(x,y)=3x^2+2xy+3y^2$ on the unit circle $x^2+y^2=1$. Lagrange's condition gives $\begin{pmatrix}3&1\\1&3\end{pmatrix}\begin{pmatrix}x\\y\end{pmatrix}=\lambda\begin{pmatrix}x\\y\end{pmatrix}$, an eigenvalue problem with eigenvalues $4$ (direction $(1,1)/\sqrt2$) and $2$ (direction $(1,-1)/\sqrt2$). The maximum is $4$, the minimum is $2$. If $f$ is the variance of a projected two-channel signal, the first direction is the common-mode component and the second is the differential (bipolar) component.
\end{workedexample}
\begin{exercises}
\item Compute the gradient and Hessian of $f(x,y)=x^2+2y^2$.
\item Use the chain rule to differentiate $g(t)=\sin(x(t)^2)$.
\item Compare L2 and L1 penalties conceptually in regression.
\item Give an example of a curved one-dimensional set in $\mathbb R^2$.
\item Show that $f(x,y)=xy/(x^2+y^2)$, $f(0,0)=0$, is discontinuous at the origin.
\item Prove that a sum of convex functions is convex, and that $\|x\|_2^2$ is convex.
\item Use Lagrange multipliers to minimize $x^2+y^2$ subject to $x+2y=5$.
\item For $f(x)=\tfrac12(x_1^2+100x_2^2)$, compute the condition number and the best fixed-step gradient-descent contraction factor.
\item Perform two Newton steps for $f(x)=x^4/4-x$ starting at $x_0=2$.
\end{exercises}

\chapter{Integration, function spaces, and transforms}\label{ch:functionspaces}
\begin{readergoal}Prepare the mathematical ground for Hilbert spaces, Fourier coefficients, energy, and stochastic-process covariance.\end{readergoal}
\section{From vectors of numbers to functions}
The Euclidean inner product $\sum_{i=1}^n x_iy_i$ measures agreement between two finite vectors. For sufficiently integrable real functions on an interval, the analogous quantity is
\[\langle f,g\rangle=\int_a^b f(t)g(t)\,dt.\]
For complex functions we conjugate the second or first argument according to convention; here we use $\langle f,g\rangle=\int_a^b f(t)\overline{g(t)}\,dt$, linear in the first argument. The induced norm is $\|f\|_2=(\int |f|^2)^{1/2}$. Signals with finite squared integral belong to the space $L^2$, after identifying functions equal almost everywhere.

\section{Orthogonality and approximation}
Two functions are orthogonal when their inner product is zero. If $e_1,\ldots,e_k$ are orthonormal, the projection of $f$ onto their span is $\sum_j\langle f,e_j\rangle e_j$. The error is orthogonal to the span: this is the same geometry later used in least squares, PCA, and finite-basis spectral expansions.

\section{Measure and probability: what does ``almost everywhere'' mean?}
A measure assigns nonnegative sizes to suitable sets. In the familiar one-dimensional setting, interval length is the basic example. The Lebesgue integral extends Riemann integration and provides a stable theory for limits of integrable functions. ``Almost everywhere'' means except on a set of measure zero. Two functions that differ at finitely many time points can represent the same $L^2$ element, even though pointwise spike events at those instants could be scientifically significant. This reminds us that the choice of mathematical object must fit the observation type.

\section{Generalized signals and impulses}
The ideal Dirac delta $\delta(t-t_0)$ is not an ordinary $L^2$ function. In generalized-function notation its defining property is $\int f(t)\delta(t-t_0)\,dt=f(t_0)$ for suitable test functions. A spike train can be represented formally as $s(t)=\sum_j\delta(t-t_j)$, but practical recordings contain finite-time waveforms and discrete timestamps. Both descriptions are useful, and neither should be confused with the other.

\begin{workedexample}[Projection onto one sinusoid]
On $[0,T]$, let $e(t)=\sqrt{2/T}\cos(2\pi kt/T)$ for a nonzero integer $k$. Its squared integral is one. The coefficient $a=\int_0^T x(t)e(t)\,dt$ determines the best approximation $ae(t)$ to $x(t)$ within this one-dimensional subspace under squared-integral error. This is the geometric core of Fourier analysis.
\end{workedexample}
\begin{exercises}
\item Verify that $\int_0^T\sin(2\pi nt/T)\cos(2\pi mt/T)\,dt=0$ for positive integers $m,n$.
\item Explain why delta functions require care in ordinary function spaces.
\item State the difference between orthogonality and statistical independence.
\end{exercises}
\chapter{Bridges from number systems to calculus}\label{ch:bridges-analysis}
\begin{readergoal}Learn to justify limits, continuity, differentiation, and integration rather than treating calculus rules as unexplained computational recipes.\end{readergoal}

\bridgeentry{limit-unique}{A convergent sequence has only one limit}
\begin{proposition}If $x_n\to a$ and $x_n\to b$ in $\mathbb R$, then $a=b$.\end{proposition}
\begin{proof}Suppose $a\neq b$. Set $\varepsilon=|a-b|/3$. For all sufficiently large $n$, both $|x_n-a|<\varepsilon$ and $|x_n-b|<\varepsilon$. The triangle inequality yields $|a-b|\leq |a-x_n|+|x_n-b|<2|a-b|/3$, a contradiction.\end{proof}
\paragraph{Example.} Two numerical routines approximating the same well-defined limit may have different finite-sample outputs; a proven convergent sequence cannot have two distinct limiting values.

\bridgeentry{limit-arithmetic}{Limits respect addition and multiplication}
\begin{proposition}If real sequences $x_n\to a$ and $y_n\to b$, then $x_n+y_n\to a+b$ and $x_ny_n\to ab$.\end{proposition}
\begin{proof}For sums, $|(x_n+y_n)-(a+b)|\leq |x_n-a|+|y_n-b|\to0$. For products, write $x_ny_n-ab=(x_n-a)y_n+a(y_n-b)$. A convergent sequence is bounded, so the first term tends to zero and the second does too.\end{proof}
\paragraph{Example.} If a mean and a calibration gain both converge in a controlled limit, their product converges. It does not follow that multiplying two biased finite-sample estimators produces an unbiased estimate.

\bridgeentry{monotone-convergence}{Bounded monotone real sequences converge}
\begin{proposition}Every nondecreasing real sequence bounded above converges to the supremum of its values.\end{proposition}
\begin{proof}Let $s=\sup\{x_n:n\geq0\}$. For any $\varepsilon>0$, the supremum approximation property provides $N$ with $s-\varepsilon<x_N\leq s$. Monotonicity then gives $s-\varepsilon<x_n\leq s$ for all $n\geq N$, which is convergence to $s$.\end{proof}
\paragraph{Example.} A running maximum of a bounded score converges mathematically, but the limit may depend on the candidate set and need not estimate true predictive accuracy.

\bridgeentry{continuity-composition}{Composition preserves continuity}
\begin{proposition}If $f$ is continuous at $a$ and $g$ is continuous at $f(a)$, then $g\circ f$ is continuous at $a$.\end{proposition}
\begin{proof}Given an output tolerance $\varepsilon$, continuity of $g$ supplies a tolerance $\eta$ on its input. Continuity of $f$ supplies $\delta$ such that $|x-a|<\delta$ implies $|f(x)-f(a)|<\eta$. The latter implies $|g(f(x))-g(f(a))|<\varepsilon$.\end{proof}
\paragraph{Example.} A continuously varying model amplitude followed by a continuous nonlinear transfer function remains continuous. Thresholding is discontinuous and can violate this conclusion.

\bridgeentry{intermediate-value}{A continuous curve cannot jump over an intermediate value}
\begin{theorem}[Intermediate value theorem]If $f$ is continuous on $[a,b]$ and $f(a)<c<f(b)$, then some $t\in(a,b)$ satisfies $f(t)=c$.\end{theorem}
\begin{proof}Let $S=\{x\in[a,b]:f(x)\leq c\}$ and $s=\sup S$. Continuity at $s$ rules out $f(s)<c$ (then a nearby point to the right would still lie in $S$), and rules out $f(s)>c$ (then nearby points to the left could not approach $s$ from $S$). Hence $f(s)=c$.\end{proof}
\paragraph{Example.} If a truly continuous filtered trace changes sign, it must cross zero. A pair of digital samples with opposite signs only implies a crossing if a continuous interpolation model is justified.

\bridgeentry{extreme-value}{A continuous function has extrema on a compact interval}
\begin{theorem}[Extreme value theorem]If $f:[a,b]\to\mathbb R$ is continuous, it attains both a minimum and a maximum.\end{theorem}
\begin{proof}Continuity on the compact interval makes $f([a,b])$ compact: for any convergent sequence of images, extract a convergent subsequence of preimages and use continuity to identify its limit. The image is thus bounded and closed, so its supremum and infimum belong to it.\end{proof}
\paragraph{Example.} A continuous fit score on a closed, bounded parameter interval has an optimum. Open intervals, discontinuities, or unconstrained coefficients require separate existence arguments.

\bridgeentry{derivative-rules}{Sum and product differentiation rules}
\begin{proposition}For functions differentiable at $t$, $(f+g)'(t)=f'(t)+g'(t)$ and $(fg)'(t)=f'(t)g(t)+f(t)g'(t)$.\end{proposition}
\begin{proof}Linearity follows by splitting the difference quotient. For the product, add and subtract $f(t+h)g(t)$ in the numerator:
\[
\frac{f(t+h)g(t+h)-f(t)g(t)}h=
 f(t+h)\frac{g(t+h)-g(t)}h+
 g(t)\frac{f(t+h)-f(t)}h.
\]
Differentiability implies continuity of $f$, so the limit yields the formula.\end{proof}
\paragraph{Example.} Differentiating an ideal amplitude-modulated oscillation contains both envelope change and carrier change, explaining why temporal derivatives mix these mechanisms.

\bridgeentry{chain-rule}{The chain rule follows from local linear approximation}
\begin{theorem}[Chain rule]If $f$ is differentiable at $t$ and $g$ at $f(t)$, then $(g\circ f)'(t)=g'(f(t))f'(t)$.\end{theorem}
\begin{proof}Write $g(y+u)=g(y)+g'(y)u+r(u)u$ where $r(u)\to0$ as $u\to0$, with $r(0)=0$. Substitute $y=f(t)$ and $u=f(t+h)-f(t)$; divide by $h$. The first term converges to $g'(f(t))f'(t)$ and the remainder to zero because $u/h$ stays bounded.\end{proof}
\paragraph{Example.} For a fitted intensity $\lambda(t)=\exp(\eta(t))$, we have $\lambda'(t)=\lambda(t)\eta'(t)$, provided $\eta$ is differentiable.

\bridgeentry{taylor-first}{Differentiability provides a first-order remainder}
\begin{proposition}A function $f$ differentiable at $a$ satisfies $f(a+h)=f(a)+f'(a)h+o(h)$ as $h\to0$.\end{proposition}
\begin{proof}Set $r(h)=[f(a+h)-f(a)]/h-f'(a)$ for nonzero $h$ and $r(0)=0$. Differentiability says $r(h)\to0$. Rearranging gives $f(a+h)=f(a)+f'(a)h+h r(h)$, and $h r(h)=o(h)$.\end{proof}
\paragraph{Example.} A local linear approximation of a nonlinear observation model is accurate only near its expansion point; the remainder is not globally negligible.

\bridgeentry{fundamental-calculus}{Integrating a rate recovers accumulated change}
\begin{theorem}[Fundamental theorem of calculus]If $f$ is continuous on $[a,b]$ and $F(t)=\int_a^t f(s)\,ds$, then $F'(t)=f(t)$ at interior points.\end{theorem}
\begin{proof}For $h\neq0$ with $t+h\in[a,b]$, $[F(t+h)-F(t)]/h$ is the average of $f$ over the short interval from $t$ to $t+h$. By continuity, this average differs from $f(t)$ by at most $\sup_{s\text{ between }t,t+h}|f(s)-f(t)|$, which tends to zero.\end{proof}
\paragraph{Example.} Integrating a conditional spike intensity gives an expected count only under an appropriate point-process model; the calculus theorem alone does not establish that the fitted intensity is correct.

\begin{exercises}
\item Prove continuity of $x\mapsto x^2$ using limit arithmetic.
\item Explain which hypothesis fails for the Heaviside step function in the intermediate value theorem.
\item Simulate first-difference derivatives of a noisy sine wave at increasing sample rates.
\end{exercises}
\part{Linear algebra and numerical models}
\chapter{Vectors, matrices, and linear mappings}\label{ch:linear}
\begin{readergoal}Learn the grammar of population matrices before using covariance, singular values, or latent factors.\end{readergoal}
\section{Arrays are not merely spreadsheets}
A vector is an element of a vector space: an object that can be added to another vector and scaled while obeying specific algebraic rules. A matrix $A\in\mathbb R^{m\times n}$ represents a linear mapping $A:\mathbb R^n\to\mathbb R^m$. Its $j$th column is the image of the $j$th standard basis vector. The product $Ax$ forms linear combinations of these columns. Dimensions are part of the meaning: $AB$ is defined only when the number of columns of $A$ equals the number of rows of $B$.

For a population of $N$ neurons sampled at $T$ time points, set $X\in\mathbb R^{T\times N}$, with rows indexing time and columns indexing neurons. Then $Xv$ is a time series obtained by weighting neurons with $v\in\mathbb R^N$. If a recording is stored in channel-by-time orientation, as often happens in MATLAB, it must be transposed before blindly applying a formula whose convention is time-by-neuron.

\section{Rank, kernels, and invertibility}
The rank of $A$ is the dimension of its image or column space. Its kernel $\ker A=\{x:Ax=0\}$ contains directions invisible to the mapping. Rank-nullity says $\dim\ker A+\rank A=n$. A square matrix is invertible precisely when its kernel contains only zero. A measurement operator with a nontrivial kernel cannot uniquely reconstruct arbitrary inputs: different sources differ by kernel vectors while producing identical observations.

\section{Inner products and projections}
The Euclidean inner product is $\langle x,y\rangle=x^\top y$, with length $\|x\|_2=\sqrt{x^\top x}$. For a matrix $U$ whose columns are orthonormal, $U^\top U=I$, the orthogonal projection onto their span is $P=UU^\top$. Check $P^2=P$ and $P^\top=P$. The residual $(I-P)x$ is orthogonal to the selected subspace.

\section{Least squares as orthogonal projection}
If $X$ has full column rank, minimizing $\|y-X\beta\|_2^2$ yields the normal equations $X^\top X\hat\beta=X^\top y$ and solution $\hat\beta=(X^\top X)^{-1}X^\top y$. In practice, QR or SVD is usually numerically safer than explicitly forming the inverse of $X^\top X$; squaring the condition number can magnify numerical problems.

\begin{workedexample}[Which axis contains the data?]
A matrix of 30 channels and 1,000,000 samples has shape $30\times10^6$. An operation computing covariances among channels should treat the 30 channels as variables, not the million sample positions. A mistaken transpose can yield a million-by-million covariance matrix or an out-of-memory failure. Mathematical notation doubles as a memory-design contract.
\end{workedexample}
\begin{exercises}
\item Verify the projection identities for a unit vector $u$ and $P=uu^\top$.
\item Find the kernel of $A=\begin{bmatrix}1&1\\2&2\end{bmatrix}$.
\item Derive the normal equations by differentiating the squared residual norm.
\item Explain why rank deficiency compromises parameter uniqueness.
\end{exercises}
\chapter{Eigenvalues, singular values, and principal components}\label{ch:pca}
\begin{readergoal}Understand the spectrum of a matrix, its geometry, and the exact optimization solved by PCA.\end{readergoal}
\section{Eigenvectors as preserved directions}
For a square matrix $A$, an eigenvector $v\ne0$ satisfies $Av=\lambda v$. The scalar $\lambda$ describes the action on that direction. A real symmetric matrix admits an orthonormal basis of real eigenvectors and a factorization $A=Q\Lambda Q^\top$. A covariance matrix is symmetric positive semidefinite, so its eigenvalues are nonnegative.

\section{Singular value decomposition}
Every real matrix $X\in\mathbb R^{T\times N}$ admits an SVD $X=U\Sigma V^\top$ with nonnegative singular values in $\Sigma$ and orthonormal singular vectors. This construction exists even for rectangular or rank-deficient matrices. The right singular vectors describe directions in variable space; the left singular vectors describe directions in observation space. Truncating to the first $k$ components gives the best rank-$k$ approximation in Frobenius norm (Eckart--Young theorem), provided singular values are ordered descending.

\section{Deriving PCA}
Let $X_c$ denote the column-centered observations and $S=X_c^\top X_c/(T-1)$. A projected observation onto a unit vector $v$ has sample variance $v^\top Sv$. The maximum-variance direction solves
\[v_1=\arg\max_{\|v\|=1}v^\top Sv.\]
Using the Lagrangian $\mathcal L(v,\lambda)=v^\top Sv-\lambda(v^\top v-1)$ gives $Sv=\lambda v$. The maximizing direction is an eigenvector of the largest eigenvalue. Subsequent components maximize variance subject to orthogonality. If $X_c=U\Sigma V^\top$, then $\lambda_j=\sigma_j^2/(T-1)$.

\section{What PCA does not promise}
PCA maximizes sample variance, not class separation, behavioral relevance, temporal predictability, or a causal mechanism. Large variance can come from artifacts, gain fluctuations, or unmodeled experimental factors. Cross-validation, null models, and careful centering are essential. As \cite{CunninghamYu2014,Churchland2012} discuss, structured population dynamics may require objectives different from variance maximization.

\begin{workedexample}[An effective dimension]
If $\lambda_1,\ldots,\lambda_N$ are nonnegative variances, the participation ratio $d_{\rm PR}=(\sum_i\lambda_i)^2/\sum_i\lambda_i^2$ equals $N$ when all eigenvalues are equal and equals one when only one is positive. It is a continuous summary of spectral concentration, not an integer count of ``real'' sources.\end{workedexample}
\begin{exercises}
\item Derive $\lambda_j=\sigma_j^2/(T-1)$ for centered observations.
\item Compare participation ratios for $(10,0,0)$ and $(4,3,3)$.
\item Explain why replacing trials by a trial average can remove single-trial fluctuations.
\item Give a case where the leading principal component tracks recording artifact rather than neural activity.
\end{exercises}
\chapter{Tensors, reduced-rank models, and inverse problems}\label{ch:inverse}
\begin{readergoal}Move from matrix factorization to multivariate measurements, and learn how identifiable models differ from underdetermined ones.\end{readergoal}
\section{More than two indices}
An experiment with neurons $n$, time $t$, and trial $r$ may be stored as a third-order tensor $\mathcal X_{ntr}$. Its modes represent different axes of variation. A CANDECOMP/PARAFAC (CP) approximation is
\[\mathcal X_{ntr}\approx\sum_{k=1}^K a_{nk}b_{tk}c_{rk}.\]
Interpretation requires confronting scaling and permutation ambiguities: multiplying one factor by $\alpha$ and dividing another by $\alpha$ leaves the product unchanged. Additional uniqueness conditions depend on the factors and rank. The mathematical pattern is broader than one physiological explanation.

\section{Reduced-rank regression}
For predictor matrix $X\in\mathbb R^{T\times p}$ and response matrix $Y\in\mathbb R^{T\times q}$, ordinary multivariate regression fits $Y\approx XB$. Reduced-rank regression requires $\rank(B)\le r$, controlling the dimension of predictive communication. Calling the resulting dimensions ``communication subspaces'' is an interpretation built on the predictive model, not proof that specific axons mediate the influence.

\section{The forward model is not the inverse model}
Let measurements obey $y=Hx+\varepsilon$, where $x$ is a source, $H$ a measurement operator, and $\varepsilon$ noise. The \emph{forward problem} computes $y$ from a chosen source and measurement physics. The \emph{inverse problem} estimates $x$ from $y$. If $H$ has a nontrivial kernel, $x$ and $x+z$ generate the same noiseless measurement whenever $Hz=0$. This is an identifiability failure, not a limitation remedied merely by better numerical precision.

For extracellular electric potentials, source geometry and volume conduction affect $H$, and several current patterns may explain the same electrode measurements. Biophysical treatments emphasize the interpretation of forward models and the assumptions behind inverse methods \cite{Buzsaki2012,Einevoll2013}. Regularization can choose one plausible solution, but it cannot magically restore missing independent information.

\section{Regularized reconstruction}
One possible estimator is $\hat x=\arg\min_x\{\|y-Hx\|_2^2+\lambda\|Lx\|_2^2\}$. If the normal equations yield an invertible matrix, $\hat x=(H^\top H+\lambda L^\top L)^{-1}H^\top y$. The operator $L$ expresses a prior preference (e.g., smoothness), while $\lambda$ sets its influence. The result answers a conditional question: which sources fit the observations well while obeying the chosen regularity assumption?

\begin{exercises}
\item Exhibit two distinct $x\in\mathbb R^2$ satisfying $[1\ 1]x=2$.
\item Explain why a CP decomposition's factors may be arbitrarily rescaled.
\item Describe what you would need to claim that a reduced-rank relationship reflects direct anatomical connectivity.
\end{exercises}
\chapter{Bridges from vectors to neural population geometry}\label{ch:bridges-linear}
\begin{readergoal}Derive basic linear-algebra tools from definitions and see exactly how they lead to PCA, projections, and regression.\end{readergoal}

\bridgeentry{span-basis}{Linear independence makes coordinate representations unique}
\begin{proposition}If $v_1,\ldots,v_k$ are linearly independent vectors, every vector in their span has exactly one representation $x=\sum_{j=1}^k a_jv_j$.\end{proposition}
\begin{proof}Existence follows from the definition of span. If $x=\sum a_jv_j=\sum b_jv_j$, subtracting gives $\sum(a_j-b_j)v_j=0$. Linear independence forces $a_j-b_j=0$ for each $j$, so coefficients are unique.\end{proof}
\paragraph{Example.} A latent subspace can have unique coordinates after fixing a basis, but another rotated basis describes the same subspace. ``Component 1'' need not be a unique biological source.

\bridgeentry{rank-nullity}{Dimensions split into visible and invisible directions}
\begin{theorem}[Rank--nullity]For a linear map $A:V\to W$ with $\dim V=n<\infty$, $\dim\ker A+\dim\operatorname{im}A=n$.\end{theorem}
\begin{proof}Choose a basis $u_1,\ldots,u_r$ of $\ker A$ and extend it to a basis of $V$ with $v_1,\ldots,v_{n-r}$. Their images $Av_1,\ldots,Av_{n-r}$ span $\operatorname{im}A$. If a linear combination of these images vanishes, the corresponding combination of $v_j$ lies in the kernel and must have zero coefficients by independence of the full basis. Thus the images form a basis of dimension $n-r$.\end{proof}
\paragraph{Example.} A lead-field map with nonzero nullspace cannot uniquely recover all neural current distributions from extracellular electrodes.

\bridgeentry{inner-cs}{Cauchy--Schwarz bounds the inner product}
\begin{theorem}[Cauchy--Schwarz]For a real inner-product space, $|\langle x,y\rangle|\leq\|x\|\|y\|$.\end{theorem}
\begin{proof}For $y\neq0$, nonnegativity of $\|x-\alpha y\|^2$ with $\alpha=\langle x,y\rangle/\|y\|^2$ gives $\|x\|^2-\langle x,y\rangle^2/\|y\|^2\geq0$. Rearranging and taking square roots proves the claim. If $y=0$ both sides vanish.\end{proof}
\paragraph{Example.} Applying this inequality to centered random variables implies $|\rho|\leq1$ when their variances are positive.

\bridgeentry{norm-triangle}{The triangle inequality follows from Cauchy--Schwarz}
\begin{proposition}For the norm $\|x\|=\sqrt{\langle x,x\rangle}$, $\|x+y\|\leq\|x\|+\|y\|$.\end{proposition}
\begin{proof}Expand $\|x+y\|^2=\|x\|^2+2\langle x,y\rangle+\|y\|^2$. The previous bridge bounds the middle term by $2\|x\|\|y\|$, making the sum no larger than $(\|x\|+\|y\|)^2$. Nonnegative square roots preserve the inequality.\end{proof}
\paragraph{Example.} Reconstruction error from two perturbations is at most the sum of their norm bounds; that inequality says nothing about whether the model is physiologically meaningful.

\bridgeentry{pythagoras}{Orthogonal components add energy, not amplitude}
\begin{proposition}If $\langle x,y\rangle=0$, then $\|x+y\|^2=\|x\|^2+\|y\|^2$.\end{proposition}
\begin{proof}The general squared-norm expansion has the cross term $2\langle x,y\rangle$. Orthogonality makes this term zero.\end{proof}
\paragraph{Example.} Orthogonal Fourier components add squared energy over a complete ideal window. Raw amplitudes are not generally additive as magnitudes.

\bridgeentry{gram-schmidt}{Independent vectors can be orthogonalized}
\begin{theorem}[Gram--Schmidt procedure]Every finite linearly independent family $v_1,\ldots,v_k$ in an inner-product space spans an orthonormal family $q_1,\ldots,q_k$ with matching partial spans.\end{theorem}
\begin{proof}Set $u_1=v_1$ and $q_1=u_1/\|u_1\|$. Recursively let $u_j=v_j-\sum_{i<j}\langle v_j,q_i\rangle q_i$. This vector is perpendicular to preceding $q_i$ by cancellation of inner products, and it cannot be zero without making $v_j$ a combination of preceding vectors. Normalize $q_j=u_j/\|u_j\|$ and proceed inductively.\end{proof}
\paragraph{Example.} Orthonormal bases permit stable projections, but finite-precision implementations often prefer Householder QR over classical Gram--Schmidt.

\bridgeentry{orthogonal-projection}{Projection solves constrained least squares}
\begin{proposition}If columns of $Q\in\mathbb R^{n\times k}$ are orthonormal, the closest vector to $y$ in $\operatorname{im}Q$ is $QQ^\top y$.\end{proposition}
\begin{proof}Write $p=QQ^\top y$ and $r=y-p$. Then $Q^\top r=Q^\top y-Q^\top QQ^\top y=0$. For any $z\in\operatorname{im}Q$, $y-z=r+(p-z)$ is a sum of orthogonal vectors, so $\|y-z\|^2=\|r\|^2+\|p-z\|^2\geq\|r\|^2$. The minimum occurs at $z=p$.\end{proof}
\paragraph{Example.} Linear neural decoding fits a projection onto a predictor subspace. Validation must still separate training and held-out trials.

\bridgeentry{covariance-psd}{Covariance matrices are positive semidefinite}
\begin{proposition}For a random vector $X\in\mathbb R^n$ with finite second moments, its covariance matrix $C$ is symmetric positive semidefinite: $v^\top Cv\geq0$ for all $v$.\end{proposition}
\begin{proof}Symmetry follows from $\Cov(X_i,X_j)=\Cov(X_j,X_i)$. Expanding the variance of the scalar $v^\top X$ yields $v^\top Cv=\Var(v^\top X)\geq0$.\end{proof}
\paragraph{Example.} PCA uses the eigendecomposition of $C$ and nonnegative eigenvalues. A numerical covariance estimate with materially negative eigenvalues calls for an implementation or modeling audit.

\begin{exercises}
\item Find a basis and nullspace for $A=\begin{bmatrix}1&1\\2&2\end{bmatrix}$.
\item Derive the ordinary least-squares normal equations by orthogonality.
\item Explain why PCA axes can change substantially when leading eigenvalues are nearly equal.
\end{exercises}
\part{Probability, statistics, and time}
\chapter{Probability begins with sets of possible outcomes}\label{ch:probability}
\begin{readergoal}Translate uncertainty into an explicit sample space, events, conditional probabilities, and independence assumptions.\end{readergoal}
\section{Experiments and events}
A random experiment has possible outcomes in a \emph{sample space} $\Omega$. An \emph{event} is a subset $A\subseteq\Omega$. A probability measure $P$ assigns numbers to events with $P(A)\ge0$, $P(\Omega)=1$, and countable additivity across disjoint events. This connects the set operations learned in Chapter~\ref{ch:sets} to the calculus of uncertainty. A model of stochastic spike counts has a different sample space from a model of electrode voltage traces; neither should be silently substituted for the other.

\section{Conditional probability}
When $P(B)>0$, define $P(A\mid B)=P(A\cap B)/P(B)$. Rearranging gives the multiplication rule $P(A\cap B)=P(A\mid B)P(B)$. If $A$ and $B$ are independent, $P(A\cap B)=P(A)P(B)$. Independence is a model property, not merely a statement that the observed sample correlation vanishes.

The law of total probability for a partition $(B_j)$ is $P(A)=\sum_jP(A\mid B_j)P(B_j)$. Bayes' theorem follows:
\[P(B_j\mid A)=\frac{P(A\mid B_j)P(B_j)}{\sum_kP(A\mid B_k)P(B_k)}.\]
In Bayesian analysis, a prior over model parameters and a likelihood for observed data yield a posterior. The posterior is conditional on both the likelihood and the prior; a credible interval and a frequentist confidence interval do not have interchangeable interpretations.

\section{Conditioning changes the question}
Two neurons can appear correlated because each varies with movement speed. Conditioning on behavior may reduce this apparent association. But conditioning can also introduce new dependencies, as in collider bias. A scientific claim must specify what was conditioned on, why, and whether the variables were measured accurately. The conditional modeling framework for spike jitter makes this especially explicit \cite{Amarasingham2012}.

\begin{workedexample}[A base-rate effect]
Suppose only 1 percent of candidate events are genuine. A detector identifies 90 percent of genuine events and falsely flags 5 percent of nongenuine events. The probability that a flagged event is genuine is $0.90(0.01)/[0.90(0.01)+0.05(0.99)]\approx0.154$. High sensitivity does not by itself imply high positive predictive value.
\end{workedexample}
\begin{exercises}
\item Derive $P(A\cup B)=P(A)+P(B)-P(A\cap B)$.
\item Give independent events and dependent events from a six-sided die experiment.
\item Explain why a detector's posterior event probability depends on prevalence.
\end{exercises}
\chapter{Random variables, expectation, and statistical moments}\label{ch:moments}
\begin{readergoal}Define distributions and calculate moments before discussing firing-rate variability or moving-window envelope statistics.\end{readergoal}
\section{A random variable is a function}
A random variable $X:\Omega\to\mathbb R$ assigns a numerical value to each possible outcome, subject to a measurability requirement that ensures events such as $\{\omega:X(\omega)\le x\}$ have defined probabilities. Its cumulative distribution function is $F_X(x)=P(X\le x)$. Discrete random variables have probability mass functions; sufficiently regular continuous variables have probability densities. A density is not itself a probability; probabilities arise through integration over regions.

\section{Means, variance, skewness, kurtosis}
The mean is $\mu=E[X]$; variance is $\sigma^2=E[(X-\mu)^2]$; standardized skewness is $\gamma_1=E[(X-\mu)^3]/\sigma^3$; and standardized kurtosis is $\gamma_2=E[(X-\mu)^4]/\sigma^4$. Some authors report excess kurtosis $\gamma_2-3$, which equals zero for a Gaussian distribution. If $\sigma=0$, standardized higher moments are undefined. Higher moments are sensitive to rare large amplitudes and may require large samples for stable estimation.

For a time-indexed envelope $A(t)$, a sliding-window empirical mean is $\hat\mu(t)=\frac1{M}\sum_{j=0}^{M-1}A(t+j\Delta t)$ with analogous central moments. If consecutive windows overlap heavily, the resulting statistics are strongly dependent: sampling an estimate every 100 ms does not create a new independent 10-s experiment at every step.

\section{Covariance and correlation}
Covariance is $\Cov(X,Y)=E[(X-EX)(Y-EY)]$. Its scale depends on units, while Pearson correlation normalizes by standard deviations: $\rho=\Cov(X,Y)/(\sigma_X\sigma_Y)$ for nonzero variances. A zero correlation does not guarantee independence in general. Let $X$ be symmetric around zero and $Y=X^2$; then $\Cov(X,Y)=E[X^3]-E[X]E[X^2]=0$, although $Y$ is determined by $X$.

\section{The law of total variance}
For random variables $X,Z$ with finite variance,
\[\Var(X)=E[\Var(X\mid Z)]+\Var(E[X\mid Z]).\]
The first term measures variation remaining after conditioning; the second measures variation in conditional means. This identity is an essential bridge to decomposing neuronal trial-to-trial variability into point-process and gain contributions \cite{Goris2014}.

\begin{workedexample}[Poisson versus overdispersed counts]
For $N\sim\mathrm{Poisson}(\lambda)$, $E[N]=\Var(N)=\lambda$. If $N\mid\Lambda\sim\mathrm{Poisson}(\Lambda)$ and $\Lambda$ varies, then $\Var(N)=E[\Lambda]+\Var(\Lambda)$. Extra variance can arise from latent fluctuations even if spikes are conditionally Poisson.
\end{workedexample}
\begin{exercises}
\item Calculate the population mean, variance, skewness, and kurtosis of the discrete values $0,1,2$ under equal probabilities.
\item Prove $\Var(aX+b)=a^2\Var(X)$.
\item Explain the difference between window overlap and an actual temporal correlation in the source.
\item Derive the overdispersion identity from total variance.
\end{exercises}
\chapter{Inference, likelihood, and stochastic processes}\label{ch:stochastic}
\begin{readergoal}Separate population parameters from estimators, and connect time-dependent random variables to autocorrelation and stationarity.\end{readergoal}
\section{Data, models, estimators}
An unknown parameter $\theta$ describes a proposed data-generating distribution $p(x\mid\theta)$. An \emph{estimator} $\hat\theta(X_1,\ldots,X_n)$ is a function of random data; its realized value is an estimate. An estimator can be biased, have high variance, or be inconsistent. Mean squared error decomposes as $\mathrm{MSE}(\hat\theta)=\Var(\hat\theta)+[E\hat\theta-\theta]^2$ when these quantities exist. Estimation is therefore not reduced to maximizing an observed fit score.

The likelihood $L(\theta;x)=p(x\mid\theta)$ views the observed data as fixed and the parameter as variable. For independent Poisson counts $n_i$ with intensities $\lambda_i(\theta)$, the log-likelihood is
\[\ell(\theta)=\sum_i\{n_i\log\lambda_i(\theta)-\lambda_i(\theta)-\log(n_i!)\}.\]
The assumption of independent counts must be justified or modified. A large likelihood does not prove the model is correct, and a small $p$-value does not give the probability that the null hypothesis is true.

\section{Time indexed random variables}
A stochastic process is a family $(X_t)_{t\in T}$ of random variables on a common probability space. For weak stationarity, the mean is constant and the covariance between $X_t$ and $X_s$ depends only on their lag $t-s$ (assuming finite second moments). Define autocovariance $C_X(\tau)=\Cov(X_t,X_{t+\tau})$. Its Fourier transform, under appropriate conditions, gives a power spectral density. Stationarity is a property of the model, not something one can prove from one short recording.

\section{Conditional structure and hidden state}
A state-dependent process may have time-varying conditional mean or variance despite an apparently stable recording. A latent state $Z_t$ might represent arousal, behavior, or an abstract computational state. The distinction between $p(X_t\mid Z_t)$ and the marginal $p(X_t)$ explains why pooling states can create spurious correlations and overdispersion. This is a key insight in interpreting correlations \cite{CohenKohn2011,Goris2014}.

\begin{workedexample}[A stochastic oscillator is not just a sine wave]
A pure sinusoid has a well-defined frequency and deterministic phase evolution. A noisy oscillatory process may have a broad spectral peak and time-varying phase. Conversely, a stochastic process with $1/f$-like spectrum can place substantial power within the theta band without containing an independent theta oscillator. The distinction motivates explicit spectral modeling \cite{Donoghue2020,Gerster2022}.
\end{workedexample}
\begin{exercises}
\item Explain what changes when observations are serially correlated but an iid likelihood is used.
\item Derive the mean and variance of an average of $n$ independent variables with equal finite variance.
\item Define weak stationarity in your own words and give a counterexample.
\item Explain why trial averaging can suppress evidence of state switches.
\end{exercises}
\chapter{Bridges from sets to random variables and spike counts}\label{ch:bridges-probability}
\begin{readergoal}Understand random events, independence, averages, covariance, and the law of total variance from elementary probability definitions.\end{readergoal}

\bridgeentry{event-sigma}{Events form a sigma-algebra rather than an arbitrary list}
\begin{proposition}If $\mathcal F$ is a sigma-algebra on $\Omega$, it contains $\varnothing$, and it is closed under finite intersections and countable intersections.\end{proposition}
\begin{proof}By definition $\Omega\in\mathcal F$, complements of members belong to $\mathcal F$, and countable unions of members belong to $\mathcal F$. Thus $\varnothing=\Omega^c$ is in $\mathcal F$. De Morgan's laws turn intersections into complements of unions of complements, so the intersections also belong to $\mathcal F$.\end{proof}
\paragraph{Example.} A predicate such as ``the spike count exceeds five'' must refer to a measurable event before assigning it a probability. In finite discrete experiments all subsets of the sample space can be used.

\bridgeentry{independent-complements}{Independence survives taking complements}
\begin{proposition}If events $A,B$ are independent, then $A^c,B$ and $A^c,B^c$ are independent.\end{proposition}
\begin{proof}Since $B=(A\cap B)\sqcup(A^c\cap B)$, finite additivity gives $\Prob(A^c\cap B)=\Prob(B)-\Prob(A\cap B)=(1-\Prob(A))\Prob(B)$. Hence $A^c$ and $B$ are independent. Repeating the same calculation with complement of $B$ yields independence of both complements.\end{proof}
\paragraph{Example.} Independence of independent-trial events includes their non-occurrences. Recorded spike counts from adjacent bins generally do not automatically satisfy this condition.

\bridgeentry{expectation-linear}{Expectation is linear even for dependent variables}
\begin{proposition}If $X,Y$ are integrable and $a,b\in\mathbb R$, then $\E[aX+bY]=a\E[X]+b\E[Y]$.\end{proposition}
\begin{proof}In a finite experiment, expand the probability-weighted sum $\sum_\omega (aX(\omega)+bY(\omega))\Prob(\{\omega\})$ and distribute multiplication and addition. For integrable variables the corresponding property follows from linearity of the Lebesgue integral. No independence assumption is used.\end{proof}
\paragraph{Example.} The expected combined spike count of two units equals the sum of their expected individual counts even if their activities are strongly correlated.

\bridgeentry{indicators}{Indicator variables encode event counts}
\begin{proposition}For measurable $A$, $\E[\mathbf1_A]=\Prob(A)$. For events $A_1,\ldots,A_n$, the count $N=\sum_{j=1}^n\mathbf1_{A_j}$ satisfies $\E N=\sum_{j=1}^n\Prob(A_j)$.\end{proposition}
\begin{proof}The indicator takes value one on $A$ and zero on $A^c$, so its expectation is $1\cdot\Prob(A)+0\cdot\Prob(A^c)$. Apply expectation linearity to the finite sum.\end{proof}
\paragraph{Example.} A per-bin count is a sum of event indicators. Its variance depends on cross-covariances, so its expectation alone cannot establish Poisson statistics.

\bridgeentry{tower}{Conditional expectation obeys the tower property}
\begin{theorem}[Law of total expectation]If $X$ is integrable and $G$ a conditioning variable, then $\E[\E[X\mid G]]=\E[X]$.\end{theorem}
\begin{proof}For a finite-valued $G$, write $\E[X\mid G=g]$ as a conditional weighted average. Multiplying by $\Prob(G=g)$ and summing over $g$ recovers the sum over every outcome, hence $\E[X]$. The general statement follows from the defining integral identity of conditional expectation with the conditioning event chosen as the entire sample space.\end{proof}
\paragraph{Example.} A latent behavioral state changes conditional firing rate, but averaging over states recovers the unconditional mean if the weighting distribution is correct.

\bridgeentry{variance-total}{Total variance separates within- and between-state effects}
\begin{theorem}[Law of total variance]If $X$ is square-integrable, then $\Var(X)=\E[\Var(X\mid G)]+\Var(\E[X\mid G])$.\end{theorem}
\begin{proof}Use $\Var(X)=\E[X^2]-(\E X)^2$ and $\Var(X\mid G)=\E[X^2\mid G]-(\E[X\mid G])^2$. Take expectations of the latter, apply the tower property to $X^2$, and add $\Var(\E[X\mid G])=\E[(\E[X\mid G])^2]-(\E X)^2$. The intermediate squared-conditional-mean terms cancel.\end{proof}
\paragraph{Example.} Trial-to-trial latent gain can produce spike-count variance above its conditional Poisson component. The decomposition does not say which physiological variable the gain represents.

\bridgeentry{covariance-bilinear}{Covariance is bilinear under finite moments}
\begin{proposition}For square-integrable $X,Y,Z$ and scalars $a,b$, covariance is linear in its first argument:
\[
\Cov(aX+bY,Z)=a\Cov(X,Z)+b\Cov(Y,Z).
\]
\end{proposition}
\begin{proof}By definition, $\Cov(W,Z)=\E[(W-\E W)(Z-\E Z)]$. Insert $W=aX+bY$, use expectation linearity to center it as $a(X-\E X)+b(Y-\E Y)$, distribute, and use expectation linearity once more.\end{proof}
\paragraph{Example.} The variance of a pooled population count includes covariance between neurons. Ignoring correlated noise can misstate the precision of a population statistic.

\bridgeentry{poisson-mixture}{Random Poisson intensity creates overdispersion}
\begin{proposition}If $N\mid G\sim\operatorname{Poisson}(\mu G)$ with $\mu>0$, $G>0$, $\E G=1$, and finite $\Var(G)$, then $\E N=\mu$ and $\Var(N)=\mu+\mu^2\Var(G)$.\end{proposition}
\begin{proof}Conditional Poisson mean and variance are both $\mu G$. Total expectation gives $\E N=\mu\E G=\mu$. Total variance gives $\Var(N)=\E[\mu G]+\Var(\mu G)=\mu+\mu^2\Var(G)$.\end{proof}
\paragraph{Example.} The Fano factor is $\Var(N)/\E N=1+\mu\Var(G)$. It differs from Fano's information-theoretic inequality. An estimator such as \texttt{\detokenize{spike_count_fano.m}} measures a sample statistic; it does not identify $G$ without a latent-variable model.

\begin{exercises}
\item Produce two dependent Bernoulli variables and verify expectation linearity numerically.
\item Derive $\Var(X+Y)=\Var(X)+\Var(Y)+2\Cov(X,Y)$.
\item Compare Poisson and lognormal-gain mixed-Poisson simulations at equal mean counts.
\end{exercises}
\part{Mathematics of electrophysiological signals}
\chapter{From a physical voltage to a sampled sequence}\label{ch:sampling}
\begin{readergoal}Understand the measurement chain, units, sample rates, aliasing, and the epistemic difference between source and signal.\end{readergoal}
\section{What electrodes actually measure}
Extracellular electrodes record potential differences relative to a reference. Their outputs reflect contributions of multiple transmembrane currents, electrode geometry, and conductive tissue. The recorded LFP is not a direct readout of a particular neuron or synapse. A physical forward model connects currents to voltages \cite{Buzsaki2012,Einevoll2013}. Analog amplification, referencing, and filtering occur before the digitized values are stored.

\section{Sampled time is a lattice}
For sampling frequency $f_s$, write $\Delta t=1/f_s$ and $x[n]=x(n\Delta t)$ for integer $n$ in the acquisition index set. The Nyquist frequency is $f_s/2$. A sinusoid at a frequency above Nyquist can produce the same discrete samples as a lower-frequency sinusoid; this is \emph{aliasing}. A pre-sampling low-pass filter is therefore required to suppress frequencies that would fold into the band of interest. Merely applying a digital low-pass filter after sampling cannot undo aliasing already introduced.

\section{The ideal sampling theorem and its assumptions}
If a continuous signal is strictly bandlimited to frequencies $|f|<f_s/2$, its samples determine the original signal through ideal sinc interpolation. Actual brain recordings are neither exactly bandlimited nor free of instrumentation artifacts. The theorem establishes a mathematical possibility under stated conditions, not a promise that every sampled waveform reproduces all physical current sources.

\section{Storage, quantization, and memory mapping}
A discrete recording is represented in finite precision. With $C$ channels, $N$ samples, and $b$ bytes per stored sample, uncompressed size is $C N b$ bytes (ignoring metadata and headers). Memory-mapped storage can support access to selected time blocks without allocating the entire dataset. It does not guarantee fast noncontiguous access: array layout, strides, page faults, and disk bandwidth still matter. Always document sample type, endian convention, orientation, sampling rate, units, and channel metadata.

\begin{workedexample}[A large acquisition]
For 128 channels at 30,000 samples/s over one hour, stored as signed 16-bit integers, the raw size is $128\cdot30000\cdot3600\cdot2=27{,}648{,}000{,}000$ bytes (about 27.65 GB in decimal units). Converting the entire array to 64-bit floats multiplies memory by four. A streaming analysis therefore need not sacrifice mathematical precision in its intermediate accumulators, but should avoid careless full-array expansion.
\end{workedexample}
\begin{exercises}
\item Determine the Nyquist frequency for an LFP sampled at 1250 Hz.
\item Calculate the storage size of 80 channels sampled at 2 kHz for 2 hours using float32.
\item Explain why the reference electrode can affect estimated coherence.
\item Distinguish a physical source model, an amplifier signal, and a floating-point array.
\end{exercises}
\chapter{Fourier analysis, convolution, and digital filtering}\label{ch:fourier}
\begin{readergoal}Derive spectral representations and see why filtering is a temporal operation with potential artifacts.\end{readergoal}
\section{The discrete Fourier transform}
For $N$ discrete samples, define
\[X[k]=\sum_{n=0}^{N-1}x[n]e^{-i2\pi kn/N},\qquad k=0,\ldots,N-1.\]
The inverse is $x[n]=N^{-1}\sum_{k=0}^{N-1}X[k]e^{i2\pi kn/N}$. The frequency index represents $f_k=kf_s/N$ modulo $f_s$. Frequency resolution in the ordinary DFT grid is $\Delta f=f_s/N=1/T$ where $T=N/f_s$. Zero-padding interpolates the frequency grid but does not create new resolving power for nearby components.

Parseval's identity for this convention is $\sum_n|x[n]|^2=N^{-1}\sum_k|X[k]|^2$. One-sided and two-sided PSD conventions differ by factors of two away from DC and Nyquist; implementations must state normalization and units. The PSD estimates power per hertz, not simply squared FFT amplitude.

\section{Convolution and linear time-invariant systems}
A discrete linear time-invariant filter has output $y[n]=\sum_m h[m]x[n-m]$, where $h$ is its impulse response. Under appropriate conditions, $Y(e^{i\omega})=H(e^{i\omega})X(e^{i\omega})$. The transfer function has both magnitude and phase. Its impulse and step responses show how brief transients can be smeared in time or transformed into artificial-looking oscillations \cite{CheveigneNelken2019}.

\section{FIR, IIR, and zero-phase filtering}
A finite impulse-response (FIR) filter has a finite set of coefficients; an infinite impulse-response (IIR) filter includes feedback. Applying a forward-backward filter can cancel net phase shift in the interior of a record but is noncausal and changes the effective magnitude response. Boundaries require special handling; filtering a short isolated segment can produce misleading edge effects. In an online decoder, a filter that uses future samples is not available without added latency.

\section{Cross-correlation and spectral duality}
For a weakly stationary process, the autocovariance $C_{xx}[\ell]=E[(X_n-\mu)(X_{n+\ell}-\mu)]$ is related to the PSD through a Fourier transform (Wiener--Khinchin). This relationship creates a bridge between lag-domain dependencies and frequency-domain power. Oscillatory-looking autocovariance and a spectral peak are related descriptions; neither alone proves the existence of a distinct biological oscillator.

\begin{workedexample}[Moving-average low-pass filter]
Take $y[n]=(x[n]+x[n-1])/2$. Its frequency response is $H(e^{i\omega})=\tfrac12(1+e^{-i\omega})=e^{-i\omega/2}\cos(\omega/2)$. Thus magnitude attenuation and phase delay both emerge from an apparently simple average. At Nyquist ($\omega=\pi$), the response is zero.
\end{workedexample}
\begin{exercises}
\item Calculate the DFT of a constant sequence and identify its DC term.
\item Derive the moving-average frequency response for three samples.
\item Explain the difference between a Fourier coefficient and a PSD estimate.
\item Design a simulation exposing a filter-induced false oscillatory event.
\end{exercises}
\chapter{Spectrograms, multitapers, and periodic-versus-aperiodic spectra}\label{ch:spectral}
\begin{readergoal}Distinguish a spectral estimate from the generating spectrum and avoid inferring oscillations from arbitrary band power.\end{readergoal}
\section{Why spectral estimates fluctuate}
A finite realization contains sampling variability and boundary effects. The periodogram's frequency resolution changes with recording duration, and its variance does not simply vanish by taking a long record without appropriate averaging. Windowing reduces leakage at the price of altered resolution. Welch estimates average periodograms across segments; the segments may overlap, but this does not make them independent replicates.

\section{Spectrograms and the time-frequency trade-off}
The short-time Fourier transform is
\[X(\tau,f)=\int_{-\infty}^{\infty}x(t)w(t-\tau)e^{-i2\pi ft}\,dt,\]
where $w$ is a localized window. Narrow time windows localize fast transients while broadening spectral resolution. Long windows sharpen frequency information but smear event timing. A spectrogram $|X(\tau,f)|^2$ is not an instantaneous power measurement at a mathematical point: it summarizes a neighborhood determined by the window.

\section{Multitaper estimation}
The multitaper method forms spectral estimates using several mutually orthogonal tapers with controlled spectral concentration. Discrete prolate spheroidal sequences are commonly selected for a time-bandwidth product $NW$; the number of useful tapers is often on the order of $2NW-1$. The smoothing bandwidth and record length must be reported. Multitaper methods are central to the Chronux approach \cite{Bokil2010}.

\section{Spectral peaks over the background}
Brain spectra often contain aperiodic structure approximated locally by a $1/f^\beta$ law, along with putative narrowband components. A rise in band power may result from an oscillatory peak, a change in the aperiodic offset, a change in the exponent, or a frequency shift. Models such as FOOOF/specparam parameterize periodic and aperiodic components; IRASA uses irregular resampling to estimate the aperiodic part. These methods have distinct failure modes, especially when peaks are wide or overlap \cite{Donoghue2020,Gerster2022}.

\begin{workedexample}[Two indistinguishable band powers]
Imagine condition A has a 10-Hz Gaussian peak over a low background while condition B has no distinct peak but a sufficiently elevated broadband background. The integrals of their spectral power from 8 to 12 Hz may be equal. The same band-power summary corresponds to different generative accounts; it is therefore not a sufficient statistic for the existence of an oscillator.
\end{workedexample}
% --- expanded edition (October 2026) ---
\section{The periodogram is inconsistent}
For a stationary process sampled at $f_s$ with $N$ samples, the periodogram is
\[
\hat S(f_k)=\frac{1}{Nf_s}\Bigl|\sum_{n=0}^{N-1}x[n]e^{-2\pi ikn/N}\Bigr|^2.
\]
\begin{proposition}
If $x[n]$ is Gaussian white noise with variance $\sigma^2$, then for $0<k<N/2$ the values $\hat S(f_k)$ are independent, each distributed as $\frac{\sigma^2}{f_s}\cdot\frac{\chi^2_2}{2}$, i.e.\ exponential with mean equal to the true spectral density $\sigma^2/f_s$.
\end{proposition}
\begin{proof}[Sketch]
The real and imaginary parts of $X[k]$ are linear combinations of independent Gaussians with orthogonal coefficient vectors $\cos(2\pi kn/N)$ and $\sin(2\pi kn/N)$, each of squared norm $N/2$. They are therefore independent $\mathcal N(0,N\sigma^2/2)$ variables, and $|X[k]|^2$ is $N\sigma^2/2$ times a $\chi^2_2$ variable. Distinct bins use orthogonal coefficient vectors and are independent.
\end{proof}
The standard deviation of an exponential variable equals its mean. Doubling the record length doubles the number of bins but leaves each bin's relative error at $100\%$. This is \emph{inconsistency}: more data give finer resolution, not lower variance. Variance falls only by averaging, across segments (Welch), tapers (multitaper), neighboring frequencies (smoothing), or trials.

\section{Bias, variance, and degrees of freedom}
Averaging $K$ approximately independent periodogram-like estimates gives an estimate distributed approximately as $S(f)\chi^2_{2K}/(2K)$. Its relative standard deviation is $1/\sqrt K$, and a $95\%$ confidence interval for $S(f)$ is
\[
\Bigl[\frac{2K\hat S(f)}{\chi^2_{2K,0.975}},\ \frac{2K\hat S(f)}{\chi^2_{2K,0.025}}\Bigr],
\]
which on a logarithmic (dB) scale has a width independent of $S(f)$. This is one reason to plot spectra in decibels. Overlapping Welch segments are correlated; with $50\%$ overlap and a Hann window the effective number of independent segments is somewhat less than the nominal number.

\section{Slepian tapers}
Among all sequences of length $N$ with unit energy, which one concentrates the most of its spectral energy inside the band $[-W,W]$? The answer is the first discrete prolate spheroidal (Slepian) sequence, the leading eigenvector of the $N\times N$ matrix
\[
A_{mn}=\frac{\sin(2\pi W(m-n))}{\pi(m-n)}\quad(m\ne n),\qquad A_{nn}=2W,
\]
with $W$ in cycles per sample. The eigenvalues $\lambda_0>\lambda_1>\cdots$ are the fractions of energy inside the band. Approximately $2NW$ of them are close to one and the rest fall rapidly to zero; the corresponding eigenvectors are mutually orthogonal. The multitaper estimate averages the $K\approx2NW-1$ well-concentrated eigenspectra, achieving $K$ degrees of freedom (in pairs) with leakage controlled at resolution $2W$. A multitaper analysis therefore has two parameters with physical meaning: the window length $T=N/f_s$ and the half-bandwidth $Wf_s$ in hertz. Their product fixes $K$.

\section{The aperiodic component}
Write a parametric spectral model as
\[
\log_{10}S(f)=b-\log_{10}(k+f^\chi)+\sum_jG_j(f),
\]
with offset $b$, optional ``knee'' $k$, exponent $\chi$, and Gaussian peaks $G_j$ in log-power. When $k=0$, the aperiodic part is a straight line of slope $-\chi$ in log--log coordinates. Two consequences follow. First, a change in $\chi$ rotates the spectrum and changes band power in opposite directions at low and high frequencies, without any oscillation changing. Second, the ratio of two band powers (e.g.\ theta/delta) depends on $\chi$ even when no peak exists. The exponent has been related to the balance of excitation and inhibition in models, but that link is model-dependent and not a theorem about tissue.

\begin{workedexample}[Resolution, variance, and record length]
A 4-s LFP segment sampled at 1000\,Hz has DFT resolution $0.25$\,Hz. With multitaper half-bandwidth $W=1$\,Hz, $NW=4$ and $K=7$ tapers give $14$ degrees of freedom; the 95\% confidence interval spans roughly $-2.7$ to $+4.0$\,dB around the estimate. Doubling $W$ to $2$\,Hz yields $K=15$ and an interval of about $-2.0$ to $+2.5$\,dB, at the cost of blurring peaks narrower than $4$\,Hz. Neither choice is ``correct''; the report must state it.
\end{workedexample}
\begin{exercises}
\item Explain the meaning of time-bandwidth product in multitaper estimation.
\item Simulate a signal whose $1/f$ slope changes but whose oscillatory peak remains fixed.
\item Explain why a spectral peak detector needs a null or baseline definition.
\item Describe two ways spectrogram edge artifacts can mislead interpretation.
\item Show that the periodogram of white noise has the same relative variance for every record length.
\item Compute the 95\% confidence interval in dB for $K=1$ and $K=10$ independent averaged periodograms.
\item For $T=2$\,s and half-bandwidth $3$\,Hz, how many Slepian tapers are well concentrated?
\item Show that a change in aperiodic exponent with fixed value at $f=20$\,Hz increases 4--8\,Hz power and decreases 30--80\,Hz power.
\end{exercises}

\chapter{Analytic signals, circular statistics, and cross-frequency coupling}\label{ch:hilbert}
\begin{readergoal}Connect complex phase to the analytic signal and understand exactly what phase-amplitude coupling measures.\end{readergoal}
\section{The analytic signal}
For a real-valued signal $x(t)$, define its Hilbert transform $\mathcal H\{x\}(t)$ by a principal-value convolution with $1/(\pi t)$, under suitable function-space assumptions. The analytic signal is
\[z(t)=x(t)+i\mathcal H\{x\}(t)=a(t)e^{i\phi(t)},\]
where $a(t)=|z(t)|$ is amplitude and $\phi(t)=\arg z(t)$ is phase. In Fourier space, analytic-signal construction suppresses negative frequencies and doubles positive frequencies (with appropriate DC and Nyquist conventions). These operations acquire a clean oscillatory interpretation for a reasonably narrowband component. Phase can be unstable when instantaneous amplitude is very small.

\section{Filtering comes before interpretation}
Often the signal is band-pass filtered to $x_b(t)$ before Hilbert transformation. The resulting $a_b(t)$ and $\phi_b(t)$ describe the \emph{filtered component}, not a unique property of the raw potential. Choice of passband, filter order, padding, and transitions can change the estimates. For nonsinusoidal waves, harmonics introduced by waveform shape can appear as high-frequency power linked to the low-frequency phase \cite{ColeVoytek2018}.

\section{Circular means and concentration}
For phases $\phi_1,\ldots,\phi_n$, set $R=n^{-1}\sum_j e^{i\phi_j}$. The mean direction is $\arg R$ when $R\ne0$, and the mean resultant length $|R|$ measures concentration between zero and one. Classical tests such as Rayleigh's test target nonuniformity in directions under specified assumptions. They do not automatically correct for temporal autocorrelation or repeated units from the same animal \cite{Berens2009}.

\section{Phase locking and phase-amplitude dependence}
The phase-locking value of phase differences is $\mathrm{PLV}=|n^{-1}\sum_j e^{i(\phi_{1j}-\phi_{2j})}|$. It measures concentration of phase differences, not amplitude modulation. A PAC statistic such as the mean vector length $|n^{-1}\sum_j A_{\mathrm{high}}(t_j)e^{i\phi_{\mathrm{low}}(t_j)}|$ quantifies phase dependence of a high-band amplitude; its scale and biases require normalization and surrogates. The modulation index based on divergence of binned amplitude distributions is a different statistic. These quantities should not be conflated.

\section{The interpretation trap}
A significant PAC value need not represent two physiologically distinct oscillators interacting. Common nonstationarity, sharp transients, harmonics, or filtering artifacts can generate cross-frequency relationships \cite{Aru2015}. Good practice checks waveform shape, aperiodic structure, trial boundaries, and meaningful surrogate null models. Early PAC observations in human cortex motivate the question, but they do not eliminate the general identifiability issues \cite{Canolty2006}.

\begin{workedexample}[One nonsinusoidal oscillator]
Consider the following waveform:
\[x(t)=\cos(2\pi f_0 t)+0.5\cos(4\pi f_0 t).\]
Its higher-frequency harmonic is phase-determined by the same fundamental oscillator. Filtering separately near $f_0$ and $2f_0$ produces components whose phase relationship is stable. It would be misleading to infer two interacting independent biological oscillators from this constructed example.
\end{workedexample}
% --- expanded edition (October 2026) ---
\section{The Hilbert transform in the frequency domain}
With the convention $\hat x(f)=\int x(t)e^{-2\pi ift}dt$, the Hilbert transform acts as the multiplier
\[
\widehat{\mathcal Hx}(f)=-i\operatorname{sgn}(f)\,\hat x(f).
\]
It shifts every positive-frequency component by $-\pi/2$ and every negative-frequency component by $+\pi/2$, leaving magnitudes unchanged. Hence $\mathcal H\cos=\sin$ and $\mathcal H\sin=-\cos$, the analytic signal of $\cos(2\pi f_0t)$ is $e^{2\pi if_0t}$, and $\widehat z(f)=2\hat x(f)\mathbf 1_{f>0}+\hat x(0)\mathbf 1_{f=0}$. Because the multiplier has modulus one, $\mathcal H$ preserves energy and $\mathcal H(\mathcal Hx)=-x$ for signals without a DC component.

\section{When is amplitude-phase decomposition meaningful?}
Any real signal has an analytic representation $a(t)e^{i\phi(t)}$, so the decomposition always exists. The question is whether $a$ and $\phi$ have the intended interpretation as ``envelope'' and ``phase of a rhythm.'' Bedrosian's theorem gives a sufficient condition: if $a(t)$ has spectrum in $[-B,B]$ and the carrier frequency exceeds $B$, then the analytic signal of $a(t)\cos(2\pi f_0t+\theta)$ is $a(t)e^{i(2\pi f_0t+\theta)}$, so the envelope is recovered exactly. In words: the amplitude must vary slowly compared with the oscillation. A theta-band filter of width 4\,Hz centred at 8\,Hz approximately satisfies this; a 1--20\,Hz ``broadband'' filter does not, and its ``phase'' mixes several rhythms.

The \emph{instantaneous frequency} $\frac1{2\pi}\phi'(t)$ is a derivative of phase and inherits noise amplification. When $a(t)$ passes near zero the phase can jump by nearly $\pi$, producing large positive or negative instantaneous frequencies that have no physiological meaning. Such points should be masked using an amplitude threshold.

\section{Distributions on the circle}
The \emph{von Mises} distribution with mean direction $\mu$ and concentration $\kappa\ge0$ has density
\[
p(\theta)=\frac{e^{\kappa\cos(\theta-\mu)}}{2\pi I_0(\kappa)},
\]
where $I_0$ is a modified Bessel function. It plays the role of the normal distribution on the circle: $\kappa=0$ is uniform, and large $\kappa$ approximates $\mathcal N(\mu,1/\kappa)$. Its expected resultant length is $A(\kappa)=I_1(\kappa)/I_0(\kappa)$.

\begin{proposition}[Bias of the resultant length]
If $\phi_1,\ldots,\phi_n$ are iid uniform on the circle, then $\E\bigl[|R|^2\bigr]=1/n$ for $R=\frac1n\sum_je^{i\phi_j}$.
\end{proposition}
\begin{proof}
$|R|^2=\frac1{n^2}\sum_{j,k}e^{i(\phi_j-\phi_k)}$. The $n$ diagonal terms equal $1$; for $j\ne k$, $\E[e^{i(\phi_j-\phi_k)}]=\E[e^{i\phi_j}]\E[e^{-i\phi_k}]=0$.
\end{proof}
Hence $|R|\approx1/\sqrt n$ for pure noise: 25 spikes give an expected PLV near $0.2$ without any locking. Comparisons of PLV across conditions with different spike counts are biased. The \emph{pairwise phase consistency} $\mathrm{PPC}=\frac{2}{n(n-1)}\sum_{j<k}\cos(\phi_j-\phi_k)$ removes the diagonal terms and is an unbiased estimator of the squared population resultant length, $\mathrm{PPC}=(n|R|^2-1)/(n-1)$. The Rayleigh test uses $Z=n|R|^2$, which is approximately exponential with mean one under uniformity, so $P(Z>z)\approx e^{-z}$ for large $n$.

\section{Two PAC estimators and their assumptions}
Let $\phi(t)$ be a low-frequency phase and $A(t)$ a high-frequency amplitude.
\begin{enumerate}[leftmargin=2em]
\item The \emph{mean vector length} $\mathrm{MVL}=|\frac1T\sum_tA(t)e^{i\phi(t)}|$ depends on the overall scale of $A$; it is usually normalized by surrogates.
\item The \emph{modulation index} bins $\phi$ into $J$ bins, computes the mean amplitude $\bar A_j$ in each, normalizes to a distribution $p_j=\bar A_j/\sum_k\bar A_k$, and reports $\mathrm{MI}=D_{\mathrm{KL}}(p\,\|\,u)/\log J$, where $D_{\mathrm{KL}}(p\,\|\,u)=\log J-H(p)$ is the KL divergence from the uniform distribution. It is invariant to scaling $A$ but depends on $J$ and on the number of cycles.
\end{enumerate}
Both estimators are positive under the null and grow with recording length differently, so they must be compared with surrogates that preserve the autocorrelation of each signal (e.g.\ cutting $A$ at a random point and swapping the two pieces), not with zero.

\begin{workedexample}[Envelope of two close sinusoids]
Let $x(t)=\cos(2\pi\cdot 8t)+\cos(2\pi\cdot 9t)=2\cos(\pi t)\cos(2\pi\cdot8.5t)$. Its analytic signal is $e^{2\pi i8t}+e^{2\pi i9t}$, with amplitude $|2\cos(\pi t)|$, a 1\,Hz beat, and phase advancing at 8.5\,Hz but jumping by $\pi$ every second when the amplitude passes through zero. An envelope fluctuation and a ``phase reset'' appear although two perfectly stationary rhythms are present. Amplitude modulation and phase discontinuities in a band-limited signal are therefore not, by themselves, evidence of transient events.
\end{workedexample}
\begin{exercises}
\item Show why the arithmetic mean of angles near $\pm\pi$ is unreliable.
\item Distinguish PLV from PAC with explicit input variables.
\item Propose a surrogate method and specify exactly what structure it preserves.
\item Explain why zero-phase filters cannot be interpreted as real-time causal processing.
\item Show that $\mathcal H\{\cos(2\pi f_0t)\}=\sin(2\pi f_0t)$ for $f_0>0$ using the frequency-domain multiplier.
\item Compute the expected PLV for 100 uniformly distributed spike phases, and compare with the value $0.2$ for 25 phases.
\item Prove the identity $\mathrm{PPC}=(n|R|^2-1)/(n-1)$.
\item Show that the modulation index is zero exactly when the binned amplitude distribution is uniform.
\item In the beat example, find the times at which the phase jumps and the size of the jump.
\end{exercises}

\chapter{Cross-spectra, coherence, and Granger prediction}\label{ch:connectivity}
\begin{readergoal}Derive dependence measures for multivariate signals and understand why prediction is not intervention.\end{readergoal}
\section{From autocovariance to cross-covariance}
For weakly stationary processes $X_t$ and $Y_t$, define $C_{xy}(\tau)=E[(X_t-\mu_x)(Y_{t+\tau}-\mu_y)]$. Its Fourier transform is the cross-spectral density $S_{xy}(f)$. The autospectra $S_{xx},S_{yy}$ describe each process; the cross-spectrum is complex and contains magnitude and phase information about their lagged relationship.

Magnitude-squared coherence is
\[\mathcal C_{xy}(f)=\frac{|S_{xy}(f)|^2}{S_{xx}(f)S_{yy}(f)},\]
when the denominator is positive. For valid spectral-density matrices, $0\le\mathcal C_{xy}(f)\le1$. Coherence is non-directional. Magnitude-only summaries discard phase and cannot by themselves indicate whether one observed signal drives another.

\section{Common input and reference contamination}
Suppose $x=s+n_x$ and $y=s+n_y$ share a signal $s$ while $n_x,n_y$ are independent noise processes. Then $x$ and $y$ can exhibit substantial coherence even though no direct interaction between their unique sources exists. Similarly, a reference electrode that carries electrical activity can insert a common component into many recordings. Bipolar referencing or spatial derivatives may reduce some effects but introduce others; one must evaluate the measurement geometry and signal-to-noise ratio \cite{Bastos2016,CohenTsuchiya2018}.

\section{Autoregressive models and direction of prediction}
A vector autoregressive model of order $p$ for $z_t\in\mathbb R^d$ is
\[z_t=\sum_{k=1}^p A_kz_{t-k}+\varepsilon_t.\]
Consider predicting $Y_t$ from its own past versus predicting it from both its own past and the past of $X_t$. A classical Gaussian linear Granger statistic compares residual variances:
\[F_{X\to Y}=\log\frac{\sigma^2_{\rm restricted}}{\sigma^2_{\rm full}}.\]
When the full model predicts better, the statistic is positive. The arrow denotes an improvement in prediction under the candidate model, not proof that manipulating $X$ would change $Y$. Latent common causes, instantaneous mixing, nonstationarity, model order, and noise can alter apparent direction. Spectral Granger extensions decompose predictive relationships by frequency under additional factorization assumptions \cite{Cekic2018}.

\section{Transfer entropy and conditional dependence}
Transfer entropy is a conditional information measure, often written $I(Y_t;X_{t-1}^{(p)}\mid Y_{t-1}^{(q)})$. It compares how informative the past of $X$ is about $Y_t$ once the past of $Y$ is known. Nonlinear estimators can be useful but require much data and careful treatment of bias. Prediction-based asymmetry is not a universal causal-identification guarantee.

\begin{workedexample}[A shared driver]
Let $s_t$ be a stochastic signal, $x_t=s_t+\eta_t$, and $y_t=s_{t-1}+\xi_t$ with independent observation noises. The lag creates directional predictability from $x$ to $y$, but the actual source is $s$, not a connection from the observed $x$ channel to the observed $y$ channel. This is a constructive counterexample to an unrestricted causal interpretation.
\end{workedexample}
\begin{exercises}
\item State why coherence is unchanged by a fixed nonzero scaling of one signal.
\item Distinguish phase locking, coherence, and Granger prediction.
\item Construct an example of common reference that raises coherence without coupling.
\item Describe checks for VAR residual autocorrelation and stability.
\end{exercises}
\chapter{Extracellular potentials and source reconstruction}\label{ch:biophysics}
\begin{readergoal}Connect physical current sources, volume conduction, referencing, and the limitations of current-source-density analysis.\end{readergoal}
\section{Potential is a field}
An extracellular electric potential $\phi(\mathbf r,t)$ is a scalar field. The electric field is $\mathbf E=-\nabla\phi$. Electrode voltages are differences in potential; changing the reference affects the data. In an idealized homogeneous isotropic quasistatic conductor, a point current source contributes potential proportional to $I/(4\pi\sigma r)$, where $\sigma$ is conductivity and $r$ the distance to the source. This approximation illustrates spatial spreading but does not accurately describe every layered or anisotropic brain tissue.

\section{Current source density}
Conservation of charge and Ohm's law under a quasistatic conductor model lead to relationships between current source density and spatial derivatives of potential. Under constant scalar conductivity, the extracellular current-source-density convention can be written $\mathrm{CSD}=-\sigma\nabla^2\phi$ (the sign depends on the source/sink convention). A laminar one-dimensional approximation replaces the Laplacian by the second depth derivative; a discrete central difference is
\[\frac{\partial^2\phi}{\partial z^2}\Big|_{z_j}\approx\frac{\phi_{j+1}-2\phi_j+\phi_{j-1}}{(\Delta z)^2}.\]
This formula assumes uniform spacing and sufficiently smooth fields. Differentiation amplifies high-spatial-frequency noise, often making smoothing or regularization necessary.

\section{Assumptions should travel with the result}
The physics depends on conductivities, boundaries, geometry, electrode positions, and reference schemes. A depth reversal in LFP phase may be consistent with a laminar generator but does not alone locate a unique set of synaptic currents. Modeling and validation with synthetic sources help distinguish plausible mechanisms \cite{Buzsaki2012,Einevoll2013}.

% --- expanded edition (October 2026) ---
\section{From Maxwell's equations to the quasistatic approximation}
In a conductor with conductivity $\sigma$ and permittivity $\varepsilon$, the current density is $\mathbf J=\sigma\mathbf E+\mathbf J^s$, where $\mathbf J^s$ is the transmembrane source current density. Charge conservation reads $\nabla\cdot\mathbf J_{\rm total}=-\partial\rho/\partial t$. For brain tissue at frequencies below a few kilohertz, three approximations are standard: (i) capacitive currents are small compared with ohmic ones, because $\omega\varepsilon/\sigma\ll1$; (ii) magnetic induction is negligible, so $\mathbf E=-\nabla\phi$; and (iii) propagation delays are negligible over centimetres. Under these \emph{quasistatic} assumptions, conservation of current becomes
\[
\nabla\cdot(\sigma\nabla\phi)=\nabla\cdot\mathbf J^s=:I_{\rm CSD},
\]
a Poisson equation relating the potential to the current source density. When $\sigma$ is constant, $\sigma\nabla^2\phi=I_{\rm CSD}$, and the CSD convention of the earlier section follows (up to the sign convention for sources and sinks). Because time enters only as a parameter, each instant of the potential is determined by the instantaneous sources; filtering by tissue is, in this approximation, frequency independent.

\section{The point-source solution and superposition}
\begin{proposition}
In an infinite homogeneous medium, the potential of a point current source $I$ at the origin is $\phi(\mathbf r)=I/(4\pi\sigma|\mathbf r|)$.
\end{proposition}
\begin{proof}
By symmetry the current flows radially and the total current through a sphere of radius $r$ is $I$, so $|\mathbf J|=I/(4\pi r^2)$. Ohm's law gives $|\mathbf E|=I/(4\pi\sigma r^2)$, and integrating $-\partial\phi/\partial r=|\mathbf E|$ with $\phi(\infty)=0$ yields the claim.
\end{proof}
Since the Poisson equation is linear, potentials of several sources add. Every neuron has as much current leaving as entering, so the simplest realistic source is a \emph{dipole}: $+I$ at $\mathbf r_+$ and $-I$ at $\mathbf r_-$ separated by $d$. Far away ($r\gg d$),
\[
\phi(\mathbf r)\approx\frac{I d\cos\theta}{4\pi\sigma r^2},
\]
where $\theta$ is the angle between the dipole axis and $\mathbf r$. Dipole potentials decay as $r^{-2}$, faster than monopoles, and change sign across the plane perpendicular to the dipole: this is the polarity reversal seen across cortical layers. Spatially aligned dipoles (pyramidal cells with parallel apical dendrites) add constructively; randomly oriented ones (many interneurons) largely cancel. The LFP is therefore weighted towards geometrically aligned populations, not towards the most active ones.

\section{Uniqueness and nonuniqueness}
\begin{theorem}[Uniqueness for the forward problem]
Let $\Omega$ be a bounded domain with smooth boundary, $\sigma>0$ smooth, and sources given. Then the solution of $\nabla\cdot(\sigma\nabla\phi)=I_{\rm CSD}$ with prescribed boundary values (or prescribed normal current and a fixed reference) is unique.
\end{theorem}
\begin{proof}
The difference $w$ of two solutions satisfies $\nabla\cdot(\sigma\nabla w)=0$ with zero boundary data. Multiply by $w$ and integrate by parts: $\int_\Omega\sigma|\nabla w|^2=0$, so $\nabla w=0$ and $w$ is constant, hence zero (or fixed by the reference).
\end{proof}
The \emph{inverse} problem has no such uniqueness. Helmholtz already observed that a current distribution inside a volume conductor can be altered by any ``silent'' source configuration, one that produces zero potential outside a region, without changing the measurements made outside it. A closed shell of dipoles pointing radially (a double layer of constant strength) is an explicit example. Any source estimate from extracellular potentials therefore rests on constraints (anatomy, sparsity, smoothness, a fixed number of dipoles) that are not determined by the data.

\section{Inverse CSD and the role of boundaries}
The standard three-point second difference estimates the CSD only at interior contacts and assumes infinite, laterally homogeneous sheets of current. When the activated region has finite lateral extent, the second difference reports spurious sources at the edges of the probe. \emph{Inverse} CSD methods instead assume a parametric source distribution (e.g.\ discs of finite radius $R$ at each contact depth), compute the forward matrix $F$ from sources to potentials, and estimate the CSD as $F^{-1}\phi$, possibly regularized. Kernel CSD performs the same task with a smoothness prior in a reproducing-kernel Hilbert space. Each method is the pseudoinverse or Tikhonov estimator of Chapter~\ref{ch:inverse} with a different forward model and prior.

\begin{workedexample}[How fast does a spike's field decay?]
Model the extracellular action potential of a soma near the electrode as a monopole-dominated near field and a dipole far field. With $\sigma=0.3$\,S/m and a peak net current of $1$\,nA, the monopole formula gives $\phi\approx1\times10^{-9}/(4\pi\cdot0.3\cdot r)$\,V; at $r=50\,\mu$m this is about $5.3\,\mu$V per nanoampere, and spike amplitudes of tens to hundreds of microvolts imply currents of several to tens of nanoamperes. Beyond about $100$--$150\,\mu$m the dipolar $r^{-2}$ decay dominates and the amplitude falls below typical noise levels, which bounds the radius from which a single electrode can isolate units.
\end{workedexample}
\begin{exercises}
\item Derive the units of $\sigma\nabla^2\phi$ in SI units.
\item Explain why a second spatial derivative is sensitive to electrode noise.
\item State two conditions needed for a one-dimensional laminar CSD approximation.
\item Derive the far-field dipole formula by expanding the two point-source potentials to first order in $d/r$.
\item Check that $\omega\varepsilon/\sigma\ll1$ at $1$\,kHz for $\sigma=0.3$\,S/m and $\varepsilon=10^{5}\varepsilon_0$ (with $\varepsilon_0\approx8.85\times10^{-12}$\,F/m).
\item Show that adding a constant to the potential leaves the CSD unchanged, and explain what this implies about referencing.
\item Explain why randomly oriented dipoles contribute little to the LFP but may contribute to multi-unit activity.
\end{exercises}

\chapter{Bridges from finite sequences to spectra}\label{ch:bridges-signals}
\begin{readergoal}Establish why discrete Fourier coefficients, filtering, power, phase, and aliasing behave as they do before treating recorded LFP samples as continuous-time signals.\end{readergoal}

\bridgeentry{fourier-orthogonality}{Complex exponentials form an orthogonal finite family}
\begin{proposition}For integers $k,\ell\in\{0,\ldots,N-1\}$, $N\geq1$,
\[\sum_{n=0}^{N-1}e^{2\pi i(k-\ell)n/N}=\begin{cases}N,&k=\ell,\\0,&k\ne\ell.\end{cases}\]\end{proposition}
\begin{proof}For $k=\ell$, every term is one. Otherwise set $r=e^{2\pi i(k-\ell)/N}$; then $r^N=1$ but $r\ne1$. The finite geometric sum equals $(1-r^N)/(1-r)=0$.\end{proof}
\paragraph{Example.} This finite orthogonality is the algebraic reason a DFT separates integer-grid Fourier bins; an off-grid sinusoid leaks across bins under a finite window.

\bridgeentry{dft-inverse}{The DFT admits an exact finite inverse}
\begin{proposition}For $x_0,\ldots,x_{N-1}\in\mathbb C$, set $X_k=\sum_{n=0}^{N-1}x_ne^{-2\pi i kn/N}$. Then
\[x_n=\frac1N\sum_{k=0}^{N-1}X_ke^{2\pi i kn/N}.\]\end{proposition}
\begin{proof}Insert the definition of $X_k$ into the inverse expression and exchange two \emph{finite} sums. The coefficient of $x_m$ is $N^{-1}\sum_k e^{2\pi i k(n-m)/N}$, which equals one for $m=n$ and zero otherwise by the previous bridge. Only $x_n$ remains.\end{proof}
\paragraph{Example.} An inverse DFT exactly recovers the original finite samples up to numerical roundoff. It does not prove recovery of the unknown continuous biological signal between sample times.

\bridgeentry{parseval-dft}{Finite Fourier energy has an exact accounting identity}
\begin{proposition}With the DFT normalization above, $\sum_n|x_n|^2=\frac1N\sum_k|X_k|^2$.\end{proposition}
\begin{proof}Insert the inverse DFT into $\sum_n x_n\bar x_n$. Expand the finite sums and group terms by pairs $(k,\ell)$. The inner sum over $n$ vanishes whenever $k\ne\ell$ by orthogonality and equals $N$ when $k=\ell$, leaving $N^{-1}\sum_k|X_k|^2$.\end{proof}
\paragraph{Example.} Periodogram normalization should respect energy units and sampling interval. ``Power'' and ``power spectral density'' differ by frequency-bin widths.

\bridgeentry{circular-convolution}{Circular convolution becomes multiplication under the DFT}
\begin{proposition}If $(x\circledast h)_n=\sum_{m=0}^{N-1}x_mh_{(n-m)\bmod N}$, then $\operatorname{DFT}(x\circledast h)_k=X_kH_k$.\end{proposition}
\begin{proof}Substitute the circular sum into the DFT, exchange finite sums, and for each $m$ set $j=(n-m)\bmod N$. Periodicity gives $e^{-2\pi ikn/N}=e^{-2\pi ikm/N}e^{-2\pi ikj/N}$. The two sums factor into $X_k$ and $H_k$.\end{proof}
\paragraph{Example.} Multiplying two equal-length FFTs implements \emph{circular} convolution. To obtain linear finite convolution one must pad sufficiently or handle boundaries separately; otherwise opposite edges mix.

\bridgeentry{shift-phase}{Time shifts rotate Fourier phase}
\begin{proposition}For a circular shift $y_n=x_{(n-n_0)\bmod N}$, the DFT is $Y_k=e^{-2\pi ikn_0/N}X_k$.\end{proposition}
\begin{proof}In the DFT sum set $m=(n-n_0)\bmod N$. Periodicity preserves the summation domain and extracts the multiplier $e^{-2\pi ikn_0/N}$. The multiplier has modulus one, so magnitudes are preserved and phases change.\end{proof}
\paragraph{Example.} Two copies of the same oscillation with a constant temporal shift can have identical PSD yet different cross-spectral phase. PSD alone cannot recover relative timing.

\bridgeentry{aliasing-equivalence}{Uniform samples cannot distinguish aliased sinusoids}
\begin{proposition}For sample rate $F_s>0$, every $f\in\mathbb R$ and integer $m$ satisfy
\[e^{2\pi i (f+mF_s)n/F_s}=e^{2\pi i fn/F_s}\quad\text{for all integers }n.\]\end{proposition}
\begin{proof}The ratio of the two exponentials equals $e^{2\pi i mn}=1$ for every integer $n$. Hence their sample sequences coincide exactly, although their continuous frequencies differ.\end{proof}
\paragraph{Example.} A 90-Hz tone sampled at 100 Hz can appear at a frequency of 10 Hz after folding; no digital filter applied afterward can establish which continuous sinusoid generated the samples. Acquisition requires appropriate analog anti-aliasing.

\begin{exercises}
\item Work out the DFT and inverse DFT of $[1,0,0,0]$ by hand.
\item Numerically compare circular convolution with linear convolution for a short impulse near an array boundary.
\item Demonstrate DFT magnitude invariance under circular shifts and the corresponding phase change.
\end{exercises}
\part{Point processes and neuronal populations}
\chapter{Spike trains as point processes}\label{ch:points}
\begin{readergoal}Model spike timestamps without pretending they are uniformly sampled voltages.\end{readergoal}
\section{Events rather than amplitudes}
A spike train is a locally finite collection of event times $\{t_1,\ldots,t_K\}$. Its counting process is $N(t)=\#\{j:t_j\le t\}$. It is nondecreasing, integer valued, and jumps at event times. In generalized-function notation $dN(t)/dt=\sum_j\delta(t-t_j)$, but the counting process is the more elementary construction.

\section{Poisson models and their scope}
For a homogeneous Poisson process with rate $\lambda>0$, counts in disjoint intervals are independent and $N(b)-N(a)\sim\mathrm{Poisson}(\lambda(b-a))$. The waiting times are exponentially distributed, and the memoryless model has no refractory period. Real neurons often violate some of these assumptions. A time-inhomogeneous Poisson model admits $\lambda(t)$, with integrated intensity $\int_a^b\lambda(s)ds$ over an interval.

A conditional intensity $\lambda(t\mid\mathcal H_t)$ is a limiting event rate given the observed history $\mathcal H_t$. Roughly,
\[\lambda(t\mid\mathcal H_t)=\lim_{\Delta\downarrow0}\frac{P(N(t+\Delta)-N(t)=1\mid\mathcal H_t)}{\Delta},\]
when the limit is defined and multiple jumps are negligible at small scales. History dependence permits refractory behavior, bursting, and coupling covariates.

\section{Binning and the loss of exact timing}
Dividing time into bins of width $\Delta$ creates counts $n_k=N((k+1)\Delta)-N(k\Delta)$. This may simplify statistics but loses information about positions inside bins. Narrower bins increase temporal precision yet yield sparse counts; wider bins stabilize rate estimates but obscure precise coincidence. A precise synchrony claim therefore needs a time-resolution convention and an explicit null model.

\begin{workedexample}[Poisson probability of a silent window]
A 10-spikes/s homogeneous Poisson process has expected count $\lambda\Delta=0.1$ in 10 ms. The chance of observing zero spikes is $e^{-0.1}\approx0.905$. Frequent silent bins are compatible with a nonzero underlying intensity.
\end{workedexample}
% --- expanded edition (October 2026) ---
\section{Constructing the homogeneous Poisson process}
Three descriptions of the homogeneous Poisson process with rate $\lambda$ are equivalent.
\begin{enumerate}[leftmargin=2em]
\item \emph{Counting description}: $N(0)=0$; counts in disjoint intervals are independent; $N(t+s)-N(t)\sim\mathrm{Poisson}(\lambda s)$.
\item \emph{Interval description}: inter-event intervals $\tau_1,\tau_2,\ldots$ are iid exponential with mean $1/\lambda$, and $t_k=\tau_1+\cdots+\tau_k$.
\item \emph{Infinitesimal description}: $P(N(t+h)-N(t)=1)=\lambda h+o(h)$ and $P(N(t+h)-N(t)\ge2)=o(h)$, with independent increments.
\end{enumerate}
\begin{proposition}
Under the counting description, the first event time is exponential with rate $\lambda$.
\end{proposition}
\begin{proof}
$P(t_1>t)=P(N(t)=0)=e^{-\lambda t}$.
\end{proof}
The interval description gives the simplest simulation: draw $\tau_k=-\log(U_k)/\lambda$ for iid uniform $U_k$ and accumulate. The \emph{conditional uniformity} property is also useful: given $N(T)=n$, the $n$ event times are distributed as the order statistics of $n$ iid uniform points on $[0,T]$. This property underlies spike-time randomization nulls; it holds for homogeneous Poisson processes and fails for any process with refractoriness or bursting.

\section{Inhomogeneous processes, thinning, and time rescaling}
For an inhomogeneous Poisson process with intensity $\lambda(t)$, define the cumulative intensity $\Lambda(t)=\int_0^t\lambda(s)\,ds$. Counts in disjoint intervals are independent and $N(b)-N(a)\sim\mathrm{Poisson}(\Lambda(b)-\Lambda(a))$.

\begin{theorem}[Time rescaling for Poisson processes]
If $t_1<t_2<\cdots$ are the events of an inhomogeneous Poisson process with continuous positive intensity, then $\Lambda(t_1),\Lambda(t_2),\ldots$ are the events of a homogeneous Poisson process with rate one. Equivalently, the rescaled intervals $\Lambda(t_k)-\Lambda(t_{k-1})$ are iid exponential with mean one.
\end{theorem}
\begin{proof}
The map $t\mapsto\Lambda(t)$ is continuous and strictly increasing. The number of rescaled events in $(\Lambda(a),\Lambda(b)]$ equals $N(b)-N(a)$, which is Poisson with mean $\Lambda(b)-\Lambda(a)$, the length of the rescaled interval; independence over disjoint intervals is preserved.
\end{proof}
The general theorem replaces $\lambda(t)$ by the conditional intensity, so that any correctly specified history-dependent model also maps spikes to a unit-rate Poisson process. Testing whether the rescaled intervals look exponential (via $u_k=1-e^{-(\Lambda(t_k)-\Lambda(t_{k-1}))}$, which should be uniform) is the basis of KS plots for spike-train models.

\emph{Thinning} simulates an inhomogeneous process: generate a homogeneous process at rate $\lambda_{\max}\ge\sup_t\lambda(t)$ and keep each event at time $t$ independently with probability $\lambda(t)/\lambda_{\max}$. The kept events form a Poisson process of intensity $\lambda(t)$.

\section{Renewal processes and interval statistics}
A \emph{renewal process} has iid inter-event intervals with a general distribution on $(0,\infty)$. It can express refractoriness (a gap near zero) and regularity, but not dependence between successive intervals. Two summary statistics are standard:
\begin{itemize}[leftmargin=2em]
\item the \emph{coefficient of variation} $\mathrm{CV}=\mathrm{sd}(\tau)/\E(\tau)$, equal to $1$ for the Poisson process, below $1$ for regular firing, and above $1$ for bursty firing;
\item the \emph{Fano factor} $F(T)=\Var N(T)/\E N(T)$ of counts in windows of length $T$.
\end{itemize}
For a stationary renewal process, $F(T)\to\mathrm{CV}^2$ as $T\to\infty$. For finite $T$, both are biased by rate fluctuations across a session, so the local measure $\mathrm{CV}_2=\frac{2|\tau_{k+1}-\tau_k|}{\tau_{k+1}+\tau_k}$, averaged over $k$, is often preferred: it compares adjacent intervals and is insensitive to slow rate drift. A gamma-distributed renewal process with shape $\kappa$ has $\mathrm{CV}=1/\sqrt\kappa$.

\section{Second-order statistics: autocorrelograms and the Bartlett spectrum}
For a stationary point process with rate $\lambda$, the \emph{conditional intensity given an event at $0$}, $\lambda\,g(\tau)$, describes how likely another event is at lag $\tau$. Its histogram estimate is the autocorrelogram. For a homogeneous Poisson process, $g\equiv1$: the autocorrelogram is flat. A dip near zero reflects refractoriness; a peak at $\pm125$\,ms reflects 8\,Hz rhythmicity. The Fourier transform of the covariance density gives the point-process spectrum (Bartlett spectrum), which tends to $\lambda$ at high frequencies for any simple process; a spike train is therefore never ``white'' or ``band-limited'' in the way a continuous signal can be.

\begin{workedexample}[Refractoriness changes the counts]
A neuron firing at 20\,spikes/s with an absolute refractory period of $d=3$\,ms can be modeled as a dead-time-modified Poisson process: intervals are $d+$Exponential$(\mu)$ with $d+1/\mu=50$\,ms. Then $\mathrm{CV}=(1/\mu)/(d+1/\mu)=47/50=0.94$, and the long-window Fano factor is $0.94^2\approx0.88$. Refractoriness alone makes counts sub-Poisson, so a measured Fano factor slightly below one is not evidence of a precise temporal code.
\end{workedexample}
\begin{exercises}
\item Derive the probability of at least one spike in a window under a homogeneous Poisson model.
\item Explain why spike sorting uncertainty can affect autocorrelograms.
\item Distinguish an observed spike count from conditional intensity.
\item Give a point-process model that excludes spikes for a fixed period after each event.
\item Show that the sum of two independent Poisson processes with rates $\lambda_1,\lambda_2$ is a Poisson process with rate $\lambda_1+\lambda_2$.
\item Prove that thinning a homogeneous process with constant retention probability $p$ gives a Poisson process with rate $p\lambda$.
\item Compute $\mathrm{CV}$ for a gamma renewal process with shape $\kappa=4$.
\item Explain why the time-rescaling theorem requires the intensity model to use only past information.
\end{exercises}

\chapter{Spike-train GLMs, resampling, and conditional hypotheses}\label{ch:glm}
\begin{readergoal}Build spike prediction models and valid null tests from conditional probabilities instead of selecting statistics first.\end{readergoal}
\section{A Poisson generalized linear model}
A GLM relates the expected count to covariates through a link function. For binned counts $n_k$ over exposure $\Delta_k$, one specification is
\[n_k\mid x_k\sim\mathrm{Poisson}(\mu_k),\qquad\log\mu_k=\log\Delta_k+\beta_0+x_k^\top\beta.\]
The offset $\log\Delta_k$ ensures that changing bin exposure has the correct interpretation in rate units. Under this model, $\lambda_k=\exp(\beta_0+x_k^\top\beta)$ is the conditional event rate. For sufficiently fine bins a Bernoulli model with appropriate link can also be used, but the likelihood assumptions differ.

\section{History filters and coupling}
A point-process GLM can augment predictors with spike-history bases or other units' recent spikes. Negative recent-self-history weights may express refractoriness; positive interactions may express predictive coupling. The estimated coefficients depend on basis choice, observation artifacts, unobserved common inputs, and regularization. Coupling filters are model parameters and not direct measurements of synaptic transmission.

\section{Model assessment rather than fit alone}
Inspect residual or time-rescaled-event structure, held-out log-likelihood, calibration, stability across bin widths, and sensitivity to covariates. Time-rescaling transforms integrated intensities between consecutive events; a correctly specified conditional-intensity model yields rescaled intervals with a known reference distribution. Failure suggests misspecification but does not identify the correct alternative automatically.

\section{Jitter and conditional resampling}
A resampling null model may fix the spike count inside each jitter window and randomize event locations within specified allowed regions. Exactness of conditional tests depends on the null distribution used for resampling. Generic ``jitter the spikes'' is not a statistical procedure until its preserved quantities, randomized quantities, boundary handling, and sampling distribution are fully specified. Conditional modeling clarifies the hypotheses supported by interval jitter and related tests \cite{Amarasingham2012}.

If $T_{\rm obs}$ is a test statistic and $T_1^*,\ldots,T_B^*$ are exchangeable null replicates under the test construction, a common Monte Carlo tail estimate is
\[\hat p=\frac{1+\sum_{b=1}^B\mathbf1\{T_b^*\ge T_{\rm obs}\}}{B+1}.\]
The plus-one convention avoids zero estimated $p$-values, but validity also requires the simulation to represent the intended null.

\begin{workedexample}[What is preserved?]
Suppose one window has three spikes. A conditional jitter method preserving this count cannot test whether the process increased its number of spikes in that window; it can test aspects of fine timing within the window. Rejection must be interpreted relative to preserved coarse structure, not as evidence for a particular synapse.
\end{workedexample}
\begin{exercises}
\item Derive the score of the Poisson log-likelihood with respect to $\beta$.
\item State the null hypothesis and exact conditioning in a 20-ms interval jitter test.
\item Explain why randomizing every timestamp independently across a session may destroy behavioral modulation.
\item Design a split-trial validation plan that avoids training and testing on overlapping windows from the same trial.
\end{exercises}
\chapter{Neuronal covariance and the geometry of population responses}\label{ch:population}
\begin{readergoal}Distinguish signal correlation, noise correlation, factor structure, and the information represented in population covariance.\end{readergoal}
\section{One neuron versus a population}
Let $r(t)=(r_1(t),\ldots,r_N(t))^\top$ represent firing-rate estimates for $N$ units. A time series $r(t)$ traces a curve through $\mathbb R^N$. Neurons' responses may change together because of shared inputs, latent states, task drive, or interactions. The covariance matrix captures second-order cofluctuation but does not identify which explanation generated it.

In repeated-stimulus experiments, \emph{signal correlation} measures similarity in mean tuning across stimuli; \emph{noise correlation} measures trial-to-trial cofluctuation after controlling the stimulus or condition. These are distinct statistics, and their estimates depend on temporal window, firing rate, spike sorting, and brain state \cite{CohenKohn2011}.

\section{Shared gain and overdispersion}
A modulated Poisson count model can be written $N\mid G\sim\mathrm{Poisson}(G\mu)$, where $G>0$ is a random gain with $E[G]=1$. Then $E[N]=\mu$ and
\[\Var(N)=\mu+\mu^2\Var(G).\]
The second term represents fluctuations in the latent gain; the decomposition follows directly from the law of total variance. For two neurons with gains $G_1,G_2$, shared gain fluctuations can generate positive count covariance. This does not require a direct synaptic link between their spike generators \cite{Goris2014}.

\section{Latent factors and factor analysis}
A linear observation model $r_t=Cz_t+d+\varepsilon_t$ represents high-dimensional responses through lower-dimensional latent $z_t$. Factor analysis separates low-rank shared covariance from modeled observation noise. GPFA adds smooth temporal covariance priors to latent trajectories. Identifying the effective latent dimensionality requires validation and attention to scaling, stationarity, and neuron subsampling \cite{CunninghamYu2014,Gao2017}.

\section{Random-matrix caution}
For a high-dimensional sample covariance, random eigenvalue fluctuations can be large even when the true covariance is isotropic. In a simple independent-Gaussian null with feature-to-sample ratio $q$, the Mar\v{c}enko--Pastur bulk has support $\sigma^2(1\pm\sqrt q)^2$ under the standard asymptotic assumptions. This is not a universal biological noise threshold; correlated trials, heterogeneous variances, and temporal smoothing violate simple null assumptions.

\begin{exercises}
\item Derive the gain-modulated variance formula using total variance.
\item Explain how simultaneous recording differs from assembling pseudo-populations across sessions.
\item Give a counterexample where noise correlation increases without a change in the direct connection between neurons.
\item Explain why a large first eigenvalue is not automatically a genuine latent neural mode.
\end{exercises}
\chapter{State-space dynamics, neural manifolds, and rotations}\label{ch:dynamics}
\begin{readergoal}Connect dynamical evolution, latent inference, and geometric structure while keeping descriptive and mechanistic claims separate.\end{readergoal}
\section{A state is more than a snapshot}
A discrete linear dynamical system has $z_{t+1}=Az_t+Bu_t+w_t$ and observations $y_t=Cz_t+v_t$. The hidden state $z_t$ evolves according to a transition rule while observations mix it with noise. If the model is linear Gaussian, Kalman filtering estimates the posterior state sequentially. Smoothing uses future as well as past observations and is not equivalent to an online estimate.

\section{Continuous dynamics and stability}
A continuous linear system obeys $\dot z=Az+Bu(t)$. The eigenvalues of $A$ govern local behavior: negative real parts give asymptotic decay for unforced linear systems, while complex eigenvalues generate damped rotational components. More generally $\dot z=f(z,u,t)$ can represent nonlinear dynamics. Stable equilibria and attractors are properties of the chosen model; inferring them from noisy, partial observations requires care.

\section{jPCA and rotational population trajectories}
Principal-component analysis selects directions of large variance, not necessarily directions displaying dynamics. The jPCA approach first reduces dimensionality, then fits an approximately skew-symmetric linear flow for the derivatives of the reduced trajectories, highlighting rotations \cite{Churchland2012}. The existence of a rotational projection is a description of population activity under chosen preprocessing; attributing that rotation to a particular recurrent mechanism requires additional evidence.

\section{Neural manifolds and coordinate nonuniqueness}
A low-dimensional manifold is a set locally parameterized by a smaller number of coordinates than its ambient dimension. If observations obey $y=g(z)$ with $z\in\mathbb R^d$, many coordinate systems for $z$ can describe the same observed trajectories. Thus latent axes may not have unique biological labels. A fitted low-dimensional model can be useful for decoding even when its coordinates are not uniquely identifiable. Reviews of movement-related manifolds emphasize dynamical and task constraints \cite{Gallego2017}.

\section{Nonlinear latent models}
Sequential autoencoders and recurrent neural networks can infer nonlinear latent dynamics from single trials, under training objectives and architecture assumptions. Models such as LFADS use latent generative dynamics for population spike data. Flexible models can also overfit. Appropriate controls include held-out likelihood, behavior decoding, simulation-based identifiability tests, and stability across recording sessions.

\begin{workedexample}[A two-dimensional rotation]
Let $A=\begin{bmatrix}0&-\omega\\\omega&0\end{bmatrix}$. Then $\dot z=Az$ has solutions $z_1(t)=R\cos(\omega t+\phi)$ and $z_2(t)=R\sin(\omega t+\phi)$. The trajectory is a circle of radius $R$. A change of orthonormal coordinates rotates the picture but leaves its essential frequency unchanged.
\end{workedexample}
\begin{exercises}
\item Determine eigenvalues of the rotation matrix above.
\item Describe a dynamical system with a stable fixed point.
\item Explain why PCA dimensions and dynamical dimensions need not coincide.
\item Name two empirical checks required before claiming a neural manifold controls behavior.
\end{exercises}
\part{Scientific inference and reproducible methods}
\chapter{Hypothesis testing, resampling, and uncertainty}\label{ch:inference}
\begin{readergoal}Understand what a test proves, how uncertainty is quantified, and why a null model must preserve the right structure.\end{readergoal}
\section{What is a hypothesis test?}
A test begins by specifying a null family $H_0$ of possible data-generating mechanisms and a test statistic $T(D)$. The \emph{null distribution} describes the variation of $T$ under that family. A $p$-value is the probability, under the stipulated null model, of observing a statistic at least as extreme as the realized one according to the prespecified rule. It is not the posterior probability that the null is true, nor is it the probability the result occurred ``by chance'' without further qualification.

\section{Type I error, power, and effect size}
A level-$\alpha$ test controls the chance of rejecting a true null, provided its assumptions hold. Power is the rejection probability under a specified alternative. Statistical significance alone says little about effect magnitude. A coherence difference that is tiny but precisely estimated may be statistically significant, while a large effect in a small dataset may remain uncertain. Report the statistic, units, effect size, sampling unit, interval estimate, and model assumptions.

\section{Resampling methods}
A bootstrap approximates an estimator's sampling variability by resampling observational units from the data under specified conditions. The units might be animals or independent trials; sampling every LFP time point independently typically destroys serial dependence. A permutation test requires an exchangeability justification under the null, often by randomizing treatment labels under a randomized design. Surrogate time shifts can test certain synchrony claims only if their shift mechanism preserves the appropriate slow structure while breaking the tested fine-time association. These methods are not interchangeable.

\section{Hierarchical data and pseudoreplication}
Samples nested within a trial, trials within sessions, and sessions within animals are rarely independent. Treating every overlapping 100-ms analysis window as a fresh biological replicate inflates nominal sample size. A hierarchical model, cluster bootstrap, or analysis at the appropriate experimental unit can address some dependencies. Repeated statistical tests across channel pairs, frequency bins, or time bins require a multiple-comparison strategy adapted to the intended claim.

\section{Bias, generalization, and validation}
Preprocessing fitted on all the data before cross-validation can leak information from test folds into training. Selection of frequency bands after inspecting the outcome similarly biases later inference unless validated separately. Hold out whole animals or sessions when the scientific goal is cross-animal generalization. Predictive success on held-out data supports a model's useful regularities, not necessarily the literal truth of its latent biological explanation.

\begin{workedexample}[An exact conditional question]
To test whether two cells exhibit synchrony below a 5-ms scale beyond fixed 20-ms window spike counts, define windows, fixed spike counts, a synchrony statistic, and a resampling mechanism that randomizes fine timing conditional on the fixed counts. Any rejection should refer to the particular 20-ms coarse information held fixed, as emphasized in \cite{Amarasingham2012}.
\end{workedexample}
% --- expanded edition (October 2026) ---
\section{The Neyman--Pearson lemma}
\begin{theorem}[Neyman--Pearson]
For testing a simple null $H_0:X\sim p_0$ against a simple alternative $H_1:X\sim p_1$, the test that rejects when the likelihood ratio $p_1(X)/p_0(X)$ exceeds a threshold $c$ (chosen so that the size is $\alpha$) has the greatest power among all tests of size at most $\alpha$.
\end{theorem}
\begin{proof}
Let $\phi^\ast=\mathbf 1\{p_1>cp_0\}$ and let $\phi$ be any test with $\E_0\phi\le\alpha=\E_0\phi^\ast$. Pointwise, $(\phi^\ast-\phi)(p_1-cp_0)\ge0$, since $\phi^\ast=1$ exactly where $p_1-cp_0>0$. Integrating gives $\E_1\phi^\ast-\E_1\phi\ge c(\E_0\phi^\ast-\E_0\phi)\ge0$.
\end{proof}
The lemma explains why likelihood-ratio statistics are the default: they are optimal for the simplest problem and asymptotically well behaved more generally. Under regularity, Wilks' theorem states that $2\log$ of the maximized likelihood ratio between nested models differing by $q$ parameters is approximately $\chi^2_q$ under the null. Nested GLMs (e.g.\ with and without a coupling term) are compared this way; non-nested models require cross-validation or information criteria instead.

\section{Confidence intervals and their duality with tests}
A $(1-\alpha)$ confidence set $C(X)$ satisfies $P_\theta(\theta\in C(X))\ge1-\alpha$ for every $\theta$. Given a family of level-$\alpha$ tests of $H_0:\theta=\theta_0$, the set of $\theta_0$ not rejected is a $(1-\alpha)$ confidence set, and conversely. The statement concerns the \emph{procedure}: before the data are seen, the random set covers the fixed parameter with the stated probability. A realized interval either contains $\theta$ or does not.

\section{Multiple comparisons}
For $m$ tests, let $V$ be the number of false rejections and $R$ the number of rejections. Two error rates are common:
\begin{itemize}[leftmargin=2em]
\item the \emph{family-wise error rate} $\mathrm{FWER}=P(V\ge1)$, controlled by Bonferroni (reject $p_i\le\alpha/m$) or the uniformly more powerful Holm step-down procedure;
\item the \emph{false discovery rate} $\mathrm{FDR}=\E[V/\max(R,1)]$, controlled by the Benjamini--Hochberg procedure.
\end{itemize}
\begin{theorem}[Benjamini--Hochberg]
Order the $p$-values $p_{(1)}\le\cdots\le p_{(m)}$, let $k$ be the largest $i$ with $p_{(i)}\le i\alpha/m$, and reject the hypotheses with the $k$ smallest $p$-values. If the null $p$-values are independent of each other and of the non-null $p$-values, then $\mathrm{FDR}\le m_0\alpha/m\le\alpha$, where $m_0$ is the number of true nulls.
\end{theorem}
Under arbitrary dependence, replacing $\alpha$ by $\alpha/\sum_{i\le m}1/i$ (Benjamini--Yekutieli) restores control. For maps over time, frequency, or channels, \emph{cluster-based permutation tests} control the FWER for the statement ``there is some effect somewhere'' by comparing the largest observed cluster statistic with its permutation distribution. They do not license claims about the exact extent or location of a cluster.

\section{The bootstrap}
Given iid data $X_1,\ldots,X_n$ and an estimator $\hat\theta$, the nonparametric bootstrap draws $n$ observations with replacement, recomputes $\hat\theta^\ast$, and repeats $B$ times. The spread of $\hat\theta^\ast$ estimates the sampling spread of $\hat\theta$. The percentile interval $[\hat\theta^\ast_{\alpha/2},\hat\theta^\ast_{1-\alpha/2}]$ is simple; bias-corrected and accelerated (BCa) intervals improve coverage for skewed estimators. The bootstrap fails, or needs modification, for extremes (e.g.\ the maximum), for very small samples, and for dependent data, where whole blocks of consecutive observations (block bootstrap) or whole clusters (cluster bootstrap) must be resampled.

\begin{workedexample}[Thirty electrodes, one hundred frequencies]
A study tests coherence change at $30$ channel pairs and $100$ frequency bins, $3000$ tests. With uncorrected $\alpha=0.05$ and no true effect, about $150$ false positives are expected. Bonferroni requires $p\le1.7\times10^{-5}$; Benjamini--Hochberg adapts to the number of true effects; a cluster test over frequency within each pair exploits the fact that true spectral effects span neighboring bins. The three answers address different claims: ``no false positive anywhere'', ``few false positives among the reported ones'', and ``some effect exists in this pair''.
\end{workedexample}
\begin{exercises}
\item State the difference between an effect-size estimate and a hypothesis-test decision.
\item Propose a hierarchical resampling unit for a dataset containing animals, sessions, trials, and time windows.
\item Explain why using the test set to choose a frequency band biases held-out performance.
\item Describe a null distribution that would be inappropriate for testing phase-amplitude coupling under strong signal nonstationarity.
\item Apply Holm and Benjamini--Hochberg at $\alpha=0.05$ to the $p$-values $0.001,0.008,0.012,0.04,0.2$.
\item Derive the likelihood-ratio test for a Poisson rate $\lambda=\lambda_0$ versus $\lambda=\lambda_1>\lambda_0$, and show it rejects for large counts.
\item Explain why bootstrapping individual time points of an autocorrelated LFP envelope underestimates uncertainty.
\item Show that the set of $\theta_0$ not rejected by level-$\alpha$ tests is a $(1-\alpha)$ confidence set.
\end{exercises}

\chapter{Numerical analysis and reproducible scientific computation}\label{ch:numerical}
\begin{readergoal}Translate mathematical equations into algorithms without overlooking finite precision, memory layout, and validation.\end{readergoal}
\section{A formula is not yet an algorithm}
A mathematical expression can define an exact quantity while computer code approximates it. Floating-point arithmetic rounds values and is not always associative: $(a+b)+c$ and $a+(b+c)$ can differ. An apparently exact bin count computed by $\lfloor T/\Delta\rfloor$ can be off by one if $T/\Delta$ is represented just below an integer. Implementations must state endpoint conventions and test boundary cases rather than rely on a universal ad hoc epsilon.

\section{Conditioning and algorithmic stability}
A well-posed problem has a solution that exists, is unique, and depends continuously on the data in an appropriate sense. An ill-conditioned problem may magnify small perturbations even under perfect implementation. A numerically unstable algorithm may magnify errors beyond what conditioning requires. Computing least-squares coefficients via QR or SVD is typically more stable than directly inverting $X^\top X$. Regularization stabilizes some inverse problems by choosing a constrained or penalized solution, at the cost of explicit modeling assumptions.

\section{Memory and complexity}
A dense sample covariance among $N$ neurons requires $O(N^2)$ storage, while a full matrix decomposition can be computationally expensive for large $N$. When sample count $T$ is much smaller than $N$, methods operating on a $T\times T$ Gram matrix can be more efficient. Signal operations can frequently be evaluated chunk by chunk, although filtering chunks requires state continuity or overlap handling. Memory mapping avoids loading full arrays but does not eliminate I/O costs.

\section{Testing the mathematical contract}
Good numerical tests use known synthetic inputs: a sinusoid with known frequency for PSD estimation; a complex exponential for phase operations; a Poisson simulation for count statistics; a low-rank signal with controlled noise for PCA; and a known VAR for Granger analysis. Property-based tests can verify constraints such as covariance symmetry, nonnegative variance, $\mathrm{PLV}\in[0,1]$, and consistency of concatenated block processing with full-array processing within a stated tolerance.

\section{Scientific metadata as part of the result}
An estimate needs sampling rate, time origin, measurement units, reference configuration, filtering parameters, electrode and trial identifiers, quality masks, random seeds, and software versions. Metadata are part of the mathematical map between the array and its real-world interpretation. A floating-point matrix detached from its coordinates can be mathematically valid yet scientifically unusable.

\begin{exercises}
\item Explain why zero-padding an FFT does not improve physical spectral resolution.
\item Design a test for chunked FIR filtering that detects boundary discontinuities.
\item List the metadata necessary to interpret a Hilbert-envelope moment series from 10-second windows.
\end{exercises}
\chapter{A complete path from first definitions to real datasets}\label{ch:roadmap}
\begin{readergoal}Practice the complete mathematical chain from a statement in English to a reproducible conclusion, with simulations preceding claims about physiology.\end{readergoal}
\section{Project I: two signals that look coupled but are not}
Start from the English claim: ``two channels have high coherence because the associated neurons communicate.'' Translate the measurable component to a statistic $\mathcal C_{xy}(f)$ and identify the unsupported causal clause. Construct independent Gaussian noises $n_x,n_y$ and a common oscillatory source $s$, then define $x=s+n_x$ and $y=s+n_y$. Calculate coherence, add measurement noise, and vary reference contamination. Compare estimates with and without the common source. Show that coherence may be high without direct communication. This project is motivated by \cite{Bastos2016,CohenTsuchiya2018}.

\section{Project II: what a theta peak actually means}
Simulate a process with a $1/f$-like background and a genuine narrowband theta oscillator. Next simulate only the aperiodic background, vary its exponent, and compute band power in both cases. Apply periodogram, Welch, and multitaper estimators; then compare parametric aperiodic fitting and an IRASA-type separation. Quantify false positives and biases across simulated repetitions before trying the methods on LFP. Here the hypotheses must distinguish a local peak \emph{over background} from integrated power over a fixed canonical band \cite{Donoghue2020,Gerster2022}.

\section{Project III: a manufactured phase-amplitude relationship}
Simulate a nonsinusoidal slow oscillator with harmonic content, and separately simulate a high-frequency oscillator whose amplitude truly depends on slow phase. Filter both into low and high bands, compute analytic signals, and compare PAC measures. Repeat with sharp transient waveforms and phase surrogates that preserve the important nuisance structure. Identify which simulated condition actually contains two distinct interacting generators. The validation target is the generation mechanism, not merely a significant PAC statistic \cite{Aru2015,ColeVoytek2018}.

\section{Project IV: spikes, jitter, and conditional inference}
Simulate independent Poisson spike trains sharing a slow time-dependent firing-rate modulation. A naive cross-correlogram may show coincidence peaks from rate cofluctuations. Define a conditional interval-jitter hypothesis holding coarse counts fixed and re-evaluate fine-time synchrony. Add precisely coordinated spikes and measure power across timescales and window widths. The lesson is to connect every resampling decision to an explicit null hypothesis \cite{Amarasingham2012}.

\section{Project V: latent neuronal dynamics}
Simulate $z_{t+1}=Az_t+w_t$ for a two-dimensional rotating latent state, then observe $N$ neurons through $r_t=Cz_t+\varepsilon_t$. Compare PCA, factor analysis, and a fitted linear state-space model. Vary the number of observed neurons, number of trials, observation noise, and dimensionality. Track held-out prediction and recovery of the latent subspace up to rotational ambiguity; avoid claiming exact recovery of nonidentifiable coordinate axes \cite{Churchland2012,CunninghamYu2014,Gallego2017}.

\section{A research conclusion requires four layers}
For each study, the final report should distinguish (1) the original English claim; (2) the formal mathematical target and its assumptions; (3) the numerical algorithm and validation outcomes; and (4) the physiological interpretation and its remaining alternatives. A clean implementation establishes only the third layer. A statistically significant result speaks to the second within its null model. A mechanistic conclusion requires evidence beyond either one in isolation.

% --- expanded edition (October 2026) ---
\section{Project VI: a point-process model of a place cell}
Simulate a neuron whose intensity depends on position $x(t)$ along a linear track through a Gaussian tuning curve, $\lambda(t)=\lambda_0+\lambda_1\exp(-(x(t)-c)^2/(2w^2))$, together with a refractory history term. Fit (i) an occupancy-normalized rate map, (ii) a Poisson GLM with a spline basis in position, and (iii) the same GLM with a history filter. Compare estimates of $c$ and $w$, compute KS plots by time rescaling, and quantify how many laps are needed for a given precision. Then make the true field drift slowly across laps and identify which fits report this drift and which average it away. The mathematical tools are Chapters~\ref{ch:points} and~\ref{ch:glm}; the scientific lesson is that ``the place field'' is an estimand defined by a model.

\section{Project VII: source mixing and laminar profiles}
Place three dipolar sources at different depths of a homogeneous medium and compute the potential at a 32-contact linear probe with $50\,\mu$m spacing using the formulas of Chapter~\ref{ch:biophysics}. Give each source an independent time course. Compute the second-difference CSD, an inverse CSD with finite-radius discs, and a PCA of the laminar profiles. Vary the lateral extent of the sources and add independent noise at each contact. Determine when each method recovers the number, depth, and time courses of the sources, and when it creates spurious sinks at probe edges. The point of the project is to know the forward model before trusting any inverse.

\section{Project VIII: population information and correlations}
Simulate $N$ neurons with tuning curves to orientation and noise covariance built from three components: private Poisson-like variability, a shared multiplicative gain, and a small differential-correlation term along $\mu'(s)$. Estimate linear Fisher information with bias correction as a function of $N$ and of the number of trials, and compare with the value obtained after shuffling trials within stimulus. Identify which covariance component limits information and whether a realistic number of trials could detect the limit. Compare the results with the formulas of Chapter~\ref{ch:population}.

\section{A checklist for each project}
\begin{enumerate}[leftmargin=2em]
\item State the English claim and the formal estimand, with units.
\item List every assumption of the generative model and of the estimator.
\item Write down the null hypothesis, what it preserves, and what it randomizes.
\item Simulate data where the answer is known, including at least one case designed to fool the method.
\item Estimate bias, variance, coverage, and false-positive rate across many simulated datasets.
\item Only then apply the method to recordings, reporting the same quantities that were validated.
\item Separate in the conclusion what was established mathematically, numerically, statistically, and physiologically.
\end{enumerate}
Each item corresponds to a part of this book: language and logic for (1), models and probability for (2) and (3), numerical analysis for (4) and (5), and the scientific chapters for (6) and (7).
\begin{exercises}
\item Draft a preregistered null hypothesis for one project, including the sampling unit, statistic, and exclusion rule.
\item Create a matrix matching five methods to their scientific estimands, key assumptions, and primary failure modes.
\item Design one negative control and one positive control for your chosen method.
\item For Project VI, derive the expected spike count in a lap from the intensity and the animal's trajectory.
\item For Project VII, compute the potential at a contact $100\,\mu$m above a dipole of moment $Id=10^{-12}$\,A\,m aligned with the probe.
\item For Project VIII, state one result that would falsify the hypothesis that differential correlations limit information in your simulation.
\end{exercises}

\appendix
\chapter{Notation, assumptions, and units}
\begin{center}\begin{tabular}{p{0.29\textwidth}p{0.60\textwidth}}\toprule
Symbol & Meaning\\\midrule
$\in,\subseteq$ & Membership in a set; inclusion of one set in another\\
$\forall,\exists,\neg,\Rightarrow$ & For all; there exists; not; logical implication\\
$\mathbb N_0,\mathbb Z,\mathbb Q,\mathbb R,\mathbb C$ & Natural including zero, integer, rational, real, complex numbers\\
$\|x\|_2,\langle x,y\rangle$ & Euclidean norm and inner product\\
$E[X],\Var(X),\Cov(X,Y)$ & Expectation, variance, covariance\\
$\hat\theta,\theta$ & Estimator/estimate and model parameter\\
$f_s,\Delta t,N$ & Sampling frequency (Hz), interval (s), sample count\\
$x[n],x(t)$ & Discrete sampled signal and continuous-time signal\\
$X[k],S_{xx}(f)$ & DFT coefficient; power spectral density\\
$z(t),a(t),\phi(t)$ & Analytic signal, amplitude envelope, phase\\
$N(t),\lambda(t\mid\mathcal H_t)$ & Point-process count and conditional intensity\\
$\mathcal C_{xy}(f)$ & Magnitude-squared coherence\\
$X\in\mathbb R^{T\times N}$ & Population matrix: times by neurons\\
$\mathcal X\in\mathbb R^{N\times T\times R}$ & Tensor: neurons by time by trials\\
\bottomrule\end{tabular}\end{center}

\section*{Explicit conventions}
PSD values have signal-units-squared per Hz under a density convention; firing rate uses spikes/s; phase is measured in radians unless labeled otherwise. Distinguish the sample count $N$ from the point-process count $N(t)$ using context. The symbol $T$ may denote a recording duration in signal chapters but number of observations in matrix expressions when explicitly defined. We use $\log$ for natural logarithm. Complex inner products are linear in their first argument by the convention stated in Chapter~\ref{ch:functionspaces}. A reported result must disclose its observational unit and any filtering or binning convention.
\chapter{Provisional mathematical map to \texttt{arsenal/path}}\label{app:arsenal}
This is a \textbf{conceptual map, not a verified file inventory}. The GitHub repository and the actual \texttt{arsenal/path} tree were not available during compilation. Names below describe analysis families and do not assert that particular MATLAB files are implemented.
\begin{longtable}{p{0.25\textwidth}p{0.34\textwidth}p{0.31\textwidth}}\toprule
Method family & Mathematical prerequisites & Checks before interpreting results\\\midrule\endhead
Binning and firing rates & Natural numbers; functions; point processes & Window edges, exposure, refractory period\\
FFT and PSD & Complex numbers; sequences; convolution & Sample rate, leakage, normalization, 1/f background\\
Hilbert envelopes & Complex phase; Fourier; filtering & Narrowband validity, edges, amplitude near zero\\
Sliding moments & Probability; expectation; central moments & Window overlap, variance conventions, autocorrelation\\
Coherence and PLV & Cross-spectra; circular statistics & Reference, common input, sample size, null design\\
PAC & Analytic signals; nonlinear statistics & Harmonics, transients, nonstationarity, surrogates\\
Granger and VAR & Linear algebra; stochastic processes; likelihood & Lag order, stationarity, latent common input\\
Spike synchrony & Point processes; conditional probability & Coarse counts, jitter window, precise null\\
PCA and SVD & Linear maps; covariance; spectral theorem & Scaling, noise, held-out stability\\
State-space models & EDOs; conditional inference; linear algebra & Observability, latent-coordinate ambiguity\\
CSD and inverse models & Multivariable derivatives; forward physics & Geometry, conductivity, regularization\\
\bottomrule\end{longtable}
For a future verified edition, each row should link to the exact MATLAB function, an explicit function signature, the corresponding mathematical equation, a reference, a tested example, and at least one synthetic negative control. An automatically generated index from the real repository would be preferable to copying guessed filenames.
\chapter{Selected answers, proof sketches, and further problems}
\section{Selected answers}
\paragraph{Why empty collections matter (Chapter~\ref{ch:english}).} The statement $\forall x\in\varnothing\,P(x)$ is true because no element provides a counterexample. This is not evidence that any spike occurred: it is a logical consequence of quantifying over an empty domain.
\paragraph{Conditional order (Chapter~\ref{ch:logic}).} ``Every cell fires at some time'' has the form $\forall u\exists tF(u,t)$. ``There exists a time when every cell fires'' is $\exists t\forall uF(u,t)$. The second implies the first for a nonempty common time domain, but not conversely.
\paragraph{Count normalization (Chapter~\ref{ch:rationals}).} Twenty spikes in 5 s gives 4/s; thirty spikes in 15 s gives 2/s. Their unweighted mean is 3/s; the pooled rate is $50/20=2.5$/s. The different answers correspond to different weighting rules.
\paragraph{Projection (Chapter~\ref{ch:linear}).} $P^2=UU^\top UU^\top=U I U^\top=P$. Also $P^\top=P$. For any $v$ in the column span of $U$, $Pv=v$.
\paragraph{Gain model (Chapter~\ref{ch:population}).} Because $N\mid G$ is Poisson with mean $G\mu$, total variance yields $\Var N=E[G\mu]+\Var(G\mu)=\mu+\mu^2\Var G$ when $E G=1$.
\section{Capstone proof and analysis problems}
\begin{enumerate}
\item Start from the definition of a function and develop the meaning of a sampled time series, including its domain, range, units, and missing data.
\item Prove that the covariance matrix of a finite centered real data matrix is positive semidefinite. Identify when it fails to be positive definite.
\item From a linear state-space model, construct two different latent coordinate systems that generate identical observations.
\item Compare a phase-locking measure with an amplitude-weighted PAC measure on signals with very low amplitude. Characterize instability.
\item Write a full conditional resampling test for spike synchrony, explicitly stating what remains fixed under the null.
\item Derive the memory footprint and main arithmetic complexity of a covariance-based versus an SVD-based approach for $T\ll N$.
\item Design and report a blinded synthetic benchmark where an algorithm is required to distinguish an oscillation from a steep aperiodic spectrum.
\end{enumerate}
These problems are deliberately cumulative. They require assembling definitions from several parts rather than treating methods as isolated formulas.
\backmatter
\begin{thebibliography}{99}
\bibitem{Amarasingham2012} Amarasingham, Asohan and Harrison, Matthew T. and Hatsopoulos, Nicholas G. and Geman, Stuart. \emph{Conditional modeling and the jitter method of spike resampling}. Journal of Neurophysiology, 2012. \url{https://doi.org/10.1152/jn.00633.2011}
\bibitem{Aru2015} Aru, Juhan and Aru, Jaan and Priesemann, Viola and others. \emph{Untangling cross-frequency coupling in neuroscience}. Current Opinion in Neurobiology, 2015. \url{https://doi.org/10.1016/j.conb.2014.08.002}
\bibitem{Bastos2016} Bastos, Andre M. and Schoffelen, Jan-Mathijs. \emph{A Tutorial Review of Functional Connectivity Analysis Methods and Their Interpretational Pitfalls}. Frontiers in Systems Neuroscience, 2016. \url{https://doi.org/10.3389/fnsys.2015.00175}
\bibitem{Berens2009} Berens, Philipp. \emph{CircStat: A MATLAB Toolbox for Circular Statistics}. Journal of Statistical Software, 2009. \url{https://doi.org/10.18637/jss.v031.i10}
\bibitem{Bokil2010} Bokil, Hemant and Andrews, Peter and Kulkarni, Jayant E. and Mehta, Samar and Mitra, Partha. \emph{Chronux: A Platform for Analyzing Neural Signals}. Journal of Neuroscience Methods, 2010. \url{https://doi.org/10.1016/j.jneumeth.2010.06.020}
\bibitem{Buccino2020} Buccino, Alessio P. and others. \emph{SpikeInterface, a unified framework for spike sorting}. eLife, 2020. \url{https://doi.org/10.7554/eLife.61834}
\bibitem{Buccino2026} Buccino, Alessio Paolo and Sridhar, Arjun and Feng, David and Svoboda, Karel and Siegle, Joshua H.. \emph{Efficient and reproducible pipelines for spike sorting large-scale electrophysiology data}. eLife, 2026. \url{https://doi.org/10.7554/eLife.110170}
\bibitem{Buzsaki2012} Buzsaki, Gyorgy and Anastassiou, Costas A. and Koch, Christof. \emph{The origin of extracellular fields and currents---EEG, ECoG, LFP and spikes}. Nature Reviews Neuroscience, 2012. \url{https://doi.org/10.1038/nrn3241}
\bibitem{Cekic2018} Cekic, Sezen and Grandjean, Didier and Renaud, Olivier. \emph{Time, frequency, and time-varying Granger-causality measures in neuroscience}. Statistics in Medicine, 2018. \url{https://doi.org/10.1002/sim.7621}
\bibitem{Churchland2012} Churchland, Mark M. and others. \emph{Neural population dynamics during reaching}. Nature, 2012. \url{https://doi.org/10.1038/nature11129}
\bibitem{CohenKohn2011} Cohen, Marlene R. and Kohn, Adam. \emph{Measuring and interpreting neuronal correlations}. Nature Neuroscience, 2011. \url{https://doi.org/10.1038/nn.2842}
\bibitem{CohenTsuchiya2018} Cohen, Dror and Tsuchiya, Naotsugu. \emph{The Effect of Common Signals on Power, Coherence and Granger Causality}. Frontiers in Systems Neuroscience, 2018. \url{https://doi.org/10.3389/fnsys.2018.00030}
\bibitem{ColeVoytek2018} Cole, Scott and Voytek, Bradley. \emph{Cycle-by-cycle analysis of neural oscillations}. bioRxiv, 2018. \url{https://doi.org/10.1101/302000}
\bibitem{CheveigneNelken2019} de Cheveigne, Alain and Nelken, Israel. \emph{Filters: When, Why, and How (Not) to Use Them}. Neuron, 2019. \url{https://doi.org/10.1016/j.neuron.2019.02.039}
\bibitem{Donoghue2020} Donoghue, Thomas and others. \emph{Parameterizing neural power spectra into periodic and aperiodic components}. Nature Neuroscience, 2020. \url{https://doi.org/10.1038/s41593-020-00744-x}
\bibitem{Einevoll2013} Einevoll, Gaute T. and Kayser, Christoph and Logothetis, Nikos K. and Panzeri, Stefano. \emph{Modelling and analysis of local field potentials for studying the function of cortical circuits}. Nature Reviews Neuroscience, 2013. \url{https://doi.org/10.1038/nrn3599}
\bibitem{Gallego2017} Gallego, Juan A. and Perich, Matthew G. and Miller, Lee E. and Solla, Sara A.. \emph{Neural Manifolds for the Control of Movement}. Neuron, 2017. \url{https://doi.org/10.1016/j.neuron.2017.05.025}
\bibitem{Gao2017} Gao, Peiran and others. \emph{A theory of multineuronal dimensionality, dynamics and measurement}. bioRxiv, 2017. \url{https://doi.org/10.1101/214262}
\bibitem{Gerster2022} Gerster, Moritz and others. \emph{Separating Neural Oscillations from Aperiodic 1/f Activity: Challenges and Recommendations}. Neuroinformatics, 2022. \url{https://doi.org/10.1007/s12021-022-09581-8}
\bibitem{Goris2014} Goris, Robbe L. T. and Movshon, J. Anthony and Simoncelli, Eero P.. \emph{Partitioning neuronal variability}. Nature Neuroscience, 2014. \url{https://doi.org/10.1038/nn.3711}
\bibitem{CunninghamYu2014} Cunningham, John P. and Yu, Byron M.. \emph{Dimensionality reduction for large-scale neural recordings}. Nature Neuroscience, 2014. \url{https://doi.org/10.1038/nn.3776}
\bibitem{HalmosNaive} Halmos, Paul R.. \emph{Naive Set Theory}. Springer, 1974. \url{https://doi.org/10.1007/978-1-4757-1645-0}
\bibitem{MendelsonNumbers} Mendelson, Elliott. \emph{Number Systems and the Foundations of Analysis}. Dover Publications, 2008. 
\bibitem{Canolty2006} Canolty, R. T. and others. \emph{High Gamma Power Is Phase-Locked to Theta Oscillations in Human Neocortex}. Science, 2006. \url{https://doi.org/10.1126/science.1128115}
\end{thebibliography}
\end{document}